\documentclass[11pt,a4paper]{article}
\usepackage[T1]{fontenc}
\usepackage[utf8]{inputenc}
\usepackage{lmodern,amsmath,amssymb,amsthm,mathtools,microtype,mathrsfs}
\usepackage[margin=25mm,top=23mm,bottom=25mm]{geometry}
\usepackage{booktabs,array,longtable,enumitem,xcolor,fancyhdr}
\usepackage[colorlinks=true,linkcolor=blue!45!black,urlcolor=blue!45!black,citecolor=blue!45!black]{hyperref}
\hypersetup{pdftitle={Reading the Alpoge--Claude RH and Dirichlet-GRH proof: working version 0.54},pdfauthor={ChatGPT},pdfsubject={Partial human-readable reconstruction with explicit arithmetic inputs},pdfkeywords={Riemann hypothesis, Lean, theta lifts, prime sums, Arakelov slopes}}
\setlength{\parindent}{0pt}
\setlength{\parskip}{5pt plus 1pt}
\setlist{topsep=3pt,itemsep=3pt,leftmargin=*}
\setlength{\headheight}{14pt}
\pagestyle{fancy}\fancyhf{}
\fancyhead[L]{\small RH and Dirichlet GRH: a reconstruction}
\fancyhead[R]{\small Working version 0.54}
\fancyfoot[C]{\thepage}
\renewcommand{\headrulewidth}{0.3pt}
\numberwithin{equation}{section}
\newtheorem{theorem}{Theorem}[section]
\newtheorem{proposition}[theorem]{Proposition}
\newtheorem{lemma}[theorem]{Lemma}
\newtheorem{corollary}[theorem]{Corollary}
\theoremstyle{remark}\newtheorem{remark}[theorem]{Remark}
\newcommand{\Rea}{\operatorname{Re}}
\newcommand{\HS}{\mathrm{HS}}
\newcommand{\tr}{\operatorname{tr}}
\newcommand{\codim}{\operatorname{codim}}
\newcommand{\covol}{\operatorname{covol}}
\newcommand{\OK}{\mathcal O_K}
\newcommand{\src}[2]{\href{https://github.com/CocoSgt/riemann-hypothesis-lean/blob/7cbf7a0c32fca1badf05931879a37c3f8a2b03b4/src/#1}{\texttt{#2}}}
\newcommand{\status}[1]{\textbf{Status:} #1}
\begin{document}
\begin{center}
{\LARGE\bfseries RH and Dirichlet GRH: new human proofs}\\[7pt]
{\large Version 0.54 --- 23 November 2026}\\[5pt]
{\large Actual sections, connected cores, exact orbital measures and Wach inclusions}
\end{center}
\textbf{Scope.} This edition contains the current audit and all new chapters.
It is a partial human reconstruction of the accepted formal proof.
Earlier named claims are in the cumulative manuscript and preserved components.
Each result below states its original inputs and exact remaining producers.
\setcounter{tocdepth}{1}
\tableofcontents
\clearpage
\begingroup\emergencystretch=3em
\section{Current proof obligations and completed suppliers: version 0.54}
\label{sec:dependency-audit-v54}

This continues the human reconstruction with the formal proof accepted
as correct. The selected sufficient route has 121 nodes and 161 edges.
There are 24 open whole obligations: 17 class C and 7 class B.
Of these, 17 are shared by ordinary RH and unrestricted Dirichlet
GRH; 7 are additional Dirichlet-GRH inputs, and none are RH-only.
The complete version-0.49 \texttt{GLOBAL\_\allowbreak LOCAL\_\allowbreak GAIN} closure remains
in force.

The four new chapters contain 90 numbered claims and bring the
cumulative inventory to 3451. Their endpoints are original
mathematical subtheorems, with the supplied objects and retained
hypotheses stated explicitly. Actual sections now give the full
smooth-projective perfect-complex producer over arbitrary base rings.
The rank-two core is preconnected, and finite points descend through
the original coefficient extension. The archimedean orbital measure
has its exact gauge, compact-volume and Gram-determinant factors.
The actual Wach coefficient map is injective, reflects units, has
prime-saturated image and has the prescribed integral intersection.

These results remove specific premises from their next uses. The
independent polarized base change, all-cover cohomology vanishing,
geometric rank and square-law producers are still needed for the
global section laws. The arithmetic core does not by itself supply
flat local charts, component dimensions, seed incidence or
\texttt{Off\allowbreak Fin}. The archimedean scalar does not determine the
unipotent discriminant covolumes or arithmetic quotient volumes.
The Wach return still retains the actual prism and the remaining
cyclotomic-kernel-fixed Frobenius matrix descent, together with the
original local compatibility and crystalline-realization premises.
The separate faithful-invariant, 'etale-basis, Herr and quantitative
comparison work follows that return. The original outer-layer
constructor requires no additional unspecified transition system.

One choice of each common object is carried through its argument:
the actual section cube and its finite model, the marked arithmetic
frame, the prescribed infinity measures and bridge, and the actual
period maps. Simplifications are accepted only where their implication
to the full original consumer is proved. In particular the arithmetic
polynomial family needs only its generic point and two specified
endpoints inside the core; it does not require every specialization
to remain there.

A class-B claim is a written human deduction at its stated scope.
The classifications and counts are inventories of this selected route,
not a minimality assertion or a completion percentage. Earlier proof
records and historical text are retained unchanged. The following
register governs the current remaining scope.

\subsection{Complete current whole-obligation register}

\paragraph{1. GLOBAL\_\allowbreak DENOMINATOR — C — Actual LinkN local table and binding producer}

Closed(Des $\mathcal C$ 3) LinkN: the actual etaSum is bounded with exact dZ0 and same-host vvTerm, subject only to the literal ShaFinOdd and summability guards.

\emph{Shared objects and choice order.} one FamP; one class $\mathcal A$ and its target-compatible subclass; one coefficient field per family datum; one x and i at each application; original selected rows, words and levels; thresholds chosen before the common family; actual D.nd; actual etaJ and Aj; sha\_\allowbreak eq

\emph{Completed inputs (version 0.54).} WIS1–WIS26 (full four-clause actual inclusion and fixed/near-parameter completeness; exact two-supplier Wach return), WLD1-WLD26 (full generic Wach lattice constructors; exact conditional three-supplier all-layer return), SSC1-SSC17 (full original Selmer SignatureChoice and family setup retaining only coherent Wach comparison), DFD1-DFD15, DNE1-DNE14 (full original ConclF-to-LinkN consumer)

\emph{Remaining producer scope.} WIS supplies the full original four-clause ioE\_\allowbreak injective\_\allowbreak saturated\_\allowbreak inter: actual injectivity, unit reflection, prime saturation and exact integral intersection. A finite multiplication determinant proves the needed Jacobson property without assuming stronger unused parameter-adic completeness. It also supplies full original padic\_\allowbreak closed, fixed-subring prime-adic completeness, fixed lifts modulo p for arbitrary H, tilt parameter-separatedness/completeness, nearPi\_\allowbreak sep\_\allowbreak complete and fixSub\_\allowbreak nearPi\_\allowbreak complete. WLD26 now retains exactly two period-ring suppliers before full WachEx: the actual Wach prism and the remaining homogeneous cyclotomic-kernel-fixed Frobenius matrix descent. The latter still needs Frobenius-filtration, actual invariant approximation, near-closedness with const\_\allowbreak dvd and matrix lifting; completed WIS inputs must be reused. Original LocTied/IsStd and Bhatt–Scholze realization premises remain explicit. Full all-layer WachEx is not yet unconditional. FaithH, EtBasis, Herr class maps and quantitative comparison remain for the complete HW/Fon/Cmp triple. No extra outer-layer transition system is introduced. Then every original local row/table bound must hold on one common ConclF family, retaining ShaFinOdd, exceptional sets, vvTerm summability and uniform thresholds. SSC, DFD12 and DNE remain available; DNE fixes K=2000(|Ca|+1)\textasciicircum{}3 before the family. Whole GLOBAL\_\allowbreak DENOMINATOR remains open.

\emph{Source.} \href{https://github.com/CocoSgt/riemann-hypothesis-lean/blob/7cbf7a0c32fca1badf05931879a37c3f8a2b03b4/src/Formalism/SharpRestriction/GlobalTheoremInputs.lean#L45-L48}{LinkN}.

\emph{Source.} \href{https://github.com/CocoSgt/riemann-hypothesis-lean/blob/7cbf7a0c32fca1badf05931879a37c3f8a2b03b4/src/Formalism/SharpRestriction/GlobalTheoremInputs.lean#L40-L44}{LinkN\_\allowbreak at}.

\emph{Source.} \href{https://github.com/CocoSgt/riemann-hypothesis-lean/blob/7cbf7a0c32fca1badf05931879a37c3f8a2b03b4/src/Development/AutomorphicArithmetic/CharacterBoxes/DivinNOfHypV.lean#L117-L134}{divinN\_\allowbreak of\_\allowbreak hypV}.



\paragraph{2. GLOBAL\_\allowbreak RETAINED\_\allowbreak DEFECT — C — Original retained-row defect bound LinkIRR}

Rows.DelpNonneg and delp$\le$exp(-x\textasciicircum{}(ez/7)) uniformly for every late actual row datum. CURRENT VERSION0.47 QUALIFICATION: Finish the same-D/common-family adelic setting after ASG: normalized nonzero Bessel W and decay, measured-field block, original Siegel-Weil rescaling and measure/character ties, Bessel uniqueness/local ties, then the original AutTie/RowDictTie/LabTie. RowsFinDim is now proved universally by RFD. Supply the actual CountIdempotentsClass theorem or its three sufficient arithmetic cylinder counts, with one simultaneous family threshold; reuse the completed ARO structure. Completed new original producers: ARO1-ARO30, RFD1-RFD9, ASG1-ASG18. Any earlier narrative naming these completed producers as outstanding is historical; this qualification is the current scope authority. CURRENT VERSION0.48 QUALIFICATION: Finish the same-D/common-family adelic setting after ASG and AMD: normalized nonzero Bessel W and decay, original orbital/fiber and Siegel-Weil measure/character ties, Bessel uniqueness/local ties, and the full original transport theorem, then AutTie/RowDictTie/LabTie. The complete MExists measured-data block and unipotent covolume normalization are now supplied, and universal RowsFinDim remains available. Supply the actual CountIdempotentsClass theorem or its three sufficient arithmetic cylinder counts with one simultaneous family threshold. Completed current producers or conditional consumers: ARO1-ARO30, RFD1-RFD9, ASG1-ASG18, AMD1-AMD17. Earlier frontier descriptions are historical; this qualification is the current scope authority. Original whole-contract premises and quantifiers are retained. CURRENT VERSION0.49 QUALIFICATION: Prove the exact original BesselNVGap nonzero integral on its supplied discrete-series, compact-representation and character data. CoefWeightedLoc, WlocExists with its original nonvanishing antecedent, WFromLoc, the normalized global product and all-later-measured-block decay are complete. On the inherited Forms0 scope this leaves BesselNVGap for ConcreteSettingExists. Continue original orbital/fiber and Siegel-Weil measure/character ties, the required Bessel uniqueness/local ties and the full original transport theorem, then AutTie/RowDictTie/LabTie. Supply the actual CountIdempotentsClass theorem or its three sufficient arithmetic cylinder counts with one simultaneous family threshold. Completed current producers or conditional consumers: ARO1-ARO30, RFD1-RFD9, ASG1-ASG18, AMD1-AMD17, BWP1-BWP24 at their exact BesselNVGap/Forms0 scope. Earlier frontier descriptions are historical; this qualification is the current scope authority. Original whole-contract premises and quantifiers are retained.

\emph{Shared objects and choice order.} one FamP; one class $\mathcal A$ and its target-compatible subclass; one coefficient field per family datum; one x and i at each application; original selected rows, words and levels; thresholds chosen before the common family; AllRows A; LabelCountGoodInjective; AllJoinA; CountAllRows

\emph{Completed inputs (version 0.54).} SSC full Selmer-signature comparison, DFD common-family subresults, DNE full separate denominator consumer, WRS1-WRS10 full selected-class local gain, Inherited first DivDef3 and row-pair integrability

\emph{Remaining producer scope.} Complete the original F3 local coefficient bridges and F4 common-overfield readings, with valuation/norm/GoodCoef/iq and the entire coefficient field, and actual complex Hecke standard/no-Jordan plus arithmetic character data on the same form datum. Construct the actual base rows, host/frame lengths, nonzero chains and group ties, MultEq support, eigensystem/carrier/parts, all-embedding logLam, HD/d3b and NEb. Keep the stronger second DivRows3 scope distinct from the completed first DivDef3 and row-pair integrability. The full arithmetic entry and conditional common-frame construction do not by themselves install HOST\_\allowbreak NUMBER\_\allowbreak S, HOST\_\allowbreak FRAME\_\allowbreak LENGTH, TASK\_\allowbreak HOST\_\allowbreak VAN\_\allowbreak E\_\allowbreak INSTALL, TASK\_\allowbreak HOST\_\allowbreak XPP\_\allowbreak INSTALL or TASK\_\allowbreak COUNT\_\allowbreak V\_\allowbreak INSTALL. Continue the distinct retained-defect, period-height and normalized spectral producers, and the original compatible second-layer selection on one common arithmetic/geometric family. GLOBAL\_\allowbreak LOCAL\_\allowbreak GAIN remains fully closed at its original v0.49 scope. Complete local comparison theorems do not themselves construct the whole common family. Common design existence, NBUILT compatibility, original true-column records, shared fields/embeddings, Fits and eta agreement, simultaneous rows and eventual-family thresholds remain explicit. A stronger historical WPt/TSel construction may serve other producers but is not needed to reopen the proved selected-class local-gain estimate.

\emph{Source.} \href{https://github.com/CocoSgt/riemann-hypothesis-lean/blob/7cbf7a0c32fca1badf05931879a37c3f8a2b03b4/src/Formalism/SharpRestriction/GlobalTheoremInputs.lean#L55-L58}{LinkIRR}.

\emph{Source.} \href{https://github.com/CocoSgt/riemann-hypothesis-lean/blob/7cbf7a0c32fca1badf05931879a37c3f8a2b03b4/src/Development/AbelianVarietyTheory/Proofs/GlobalInterpolation/LinkIRROfJoinEx.lean#L17-L38}{linkIRR\_\allowbreak of\_\allowbreak join\_\allowbreak ex}.

\emph{Source.} \href{https://github.com/CocoSgt/riemann-hypothesis-lean/blob/7cbf7a0c32fca1badf05931879a37c3f8a2b03b4/src/Development/AbelianVarietyTheory/Proofs/GlobalInterpolation/LinkIRROfJoinEx.lean#L143-L166}{link\_\allowbreak irr}.



\paragraph{3. GLOBAL\_\allowbreak PERIOD\_\allowbreak HEIGHT — C — Full actual period and height comparison LinkDIVt}

Closed conditional TuSize -> DGamSize -> (LinkEta -> DIVinN -> LinkDIV), with every original LinkDIV\_\allowbreak at clause.

\emph{Shared objects and choice order.} one FamP; one class $\mathcal A$ and its target-compatible subclass; one coefficient field per family datum; one x and i at each application; original selected rows, words and levels; thresholds chosen before the common family; same RR target; same Bu/besS; same Nu and Bk

\emph{Source.} \href{https://github.com/CocoSgt/riemann-hypothesis-lean/blob/7cbf7a0c32fca1badf05931879a37c3f8a2b03b4/src/Formalism/SharpRestriction/GlobalTheoremInputs.lean#L77-L90}{LinkDIV\_\allowbreak at}.

\emph{Source.} \href{https://github.com/CocoSgt/riemann-hypothesis-lean/blob/7cbf7a0c32fca1badf05931879a37c3f8a2b03b4/src/Formalism/SharpRestriction/GlobalTheoremInputs.lean#L94-L96}{LinkDIVt}.

\emph{Source.} \href{https://github.com/CocoSgt/riemann-hypothesis-lean/blob/7cbf7a0c32fca1badf05931879a37c3f8a2b03b4/src/Development/AbelianVarietyTheory/Proofs/GlobalInterpolation/LinkBrest.lean#L7-L25}{GlobalTheory.Top.link\_\allowbreak brest}.

\emph{Source.} \href{https://github.com/CocoSgt/riemann-hypothesis-lean/blob/7cbf7a0c32fca1badf05931879a37c3f8a2b03b4/src/Development/AbelianVarietyTheory/Proofs/GlobalInterpolation/WireSel.lean#L11-L41}{GlobalTheory.Top.wire\_\allowbreak sel}.



\paragraph{4. GLOBAL\_\allowbreak SPECTRAL — C — Original normalized kernel/Gram estimates LinkS}

Closed TuSize -> LinkS with exact Gamma normalization, off-diagonal Hilbert-Schmidt bound, diagonal error and Dall/Dstar estimate, errors exp(-x\textasciicircum{}(epsilon/8)). CURRENT VERSION0.48 QUALIFICATION: Complete the original same-family spectral/Bessel and Siegel-Weil inputs with all measure, character, transport and common-object ties. The complete MExists measured block and prescribed-measure covolume normalization are supplied; these are not the full normalized spectral theorem. Completed current producers or conditional consumers: ASG1-ASG18, AMD1-AMD17, GCF1-GCF18 at their stated scopes. Earlier frontier descriptions are historical; this qualification is the current scope authority. Original whole-contract premises and quantifiers are retained. CURRENT VERSION0.49 QUALIFICATION: Complete the original same-family spectral/Bessel and Siegel-Weil inputs with every measure, character, transport and common-object tie. The normalized packet and decay are now constructed conditional on the exact BesselNVGap, and the analytic setting uses the original supplied Forms0 scope. These do not imply the full normalized spectral theorem. Completed current producers or conditional consumers: ASG1-ASG18, AMD1-AMD17, GCF1-GCF18 at their stated scopes, BWP1-BWP24 at their exact BesselNVGap/Forms0 scope, WRS1-WRS10 original local gain. Earlier frontier descriptions are historical; this qualification is the current scope authority. Original whole-contract premises and quantifiers are retained.

\emph{Shared objects and choice order.} one FamP; one class $\mathcal A$ and its target-compatible subclass; one coefficient field per family datum; one x and i at each application; original selected rows, words and levels; thresholds chosen before the common family; actual KD; same S.Cu spaces; GamOf; Dstar; Dall

\emph{Completed inputs (version 0.54).} AOS1–AOS18 (full exact archimedean orbital scalar, gauge, compact volume and all-section integration clause), GMR1-GMR19 (full global/local measure suppliers and RescaleStep3 on the actual setting), BQS1-BQS23 (full QuotP/QuotM, StepStatement and BesselNVGap; setting at Forms0), BST1-BST11, BSE1-BSE9, BNV1-BNV23, BSP1-BSP10, BWP1-BWP24 at original Forms0 scope, AMD1-AMD24

\emph{Remaining producer scope.} AOS supplies full original PlaceJac, universal GaugeDens, GaugeVolU, HinfCompact, PushInf and ConstTest with the literal gaugeU*volU(k)\textasciicircum{}d*Tw\textasciicircum{}(-kd) scalar. It also supplies ExInf, HtotInf, WeilInf, AinfValue/AinfClause, GinfTest and AprimeInf at every original witness and section, preserving the supplied bridge and all infinity measures. The prior GMR full RescaleStep3 and continuous Gram-fiber measures are reused; its one standard character remains fixed before all words, places and tuples. The immediate remaining RealBracket producers are the actual reference unipotent covolume/discriminant exponent, broad JacGauge/PlaceSplit and ProdExact. An independent next packet is full VolStabFromSWS and InVolumesAt from the exact original Siegel–Weil/compatibility premises, using one common admissible positive test and the same measures for unfolding and scalar cancellation. Arithmetic theta/quotient-volume equalities and finite-place factors are not supplied by the archimedean scalar. BQS already gives full BesselNVGap. Forms0, all-family Witt/genus/level/intertwiner Transport/LevelMatch, Bessel uniqueness, AutTie/RowDictTie/LabTie, simultaneous arithmetic cylinder counts and the normalized GLOBAL\_\allowbreak SPECTRAL estimate remain.

\emph{Source.} \href{https://github.com/CocoSgt/riemann-hypothesis-lean/blob/7cbf7a0c32fca1badf05931879a37c3f8a2b03b4/src/Formalism/SharpRestriction/GlobalTheoremInputs.lean#L145-L149}{LinkS\_\allowbreak at}.

\emph{Source.} \href{https://github.com/CocoSgt/riemann-hypothesis-lean/blob/7cbf7a0c32fca1badf05931879a37c3f8a2b03b4/src/Formalism/SharpRestriction/GlobalTheoremInputs.lean#L152-L157}{LinkS}.

\emph{Source.} \href{https://github.com/CocoSgt/riemann-hypothesis-lean/blob/7cbf7a0c32fca1badf05931879a37c3f8a2b03b4/src/Development/AbelianVarietyTheory/Proofs/GlobalInterpolation/LinkAlt.lean#L16-L50}{GlobalTheory.Top.link\_\allowbreak alt}.



\paragraph{5. AUT\_\allowbreak UNRAM\_\allowbreak COEF — C — Actual unramified-coefficient automorphy}

For each actual Dat with coefficient ramification one and fixed iota, choose one marked g1 supplying FinA, AutB and LemW together. CURRENT VERSION0.48 QUALIFICATION: Construct the actual AutomorphyUnramifiedCoefficients package: for every original Dat with coefficient ramification one and fixed iota, choose one marked g1 that supplies FinA, AutB and LemW together. The original InpZ record, full RigII, CNR3, finite PropGK consumer and now the entire OneChar implication are supplied. Continue RH Cases I/V and the independent additional Dirichlet-GRH CaseIII/promodularity, SeedIV, tensor-factor and rank-four-chain tasks. Completed current producers or conditional consumers: OCH1-OCH30 (the full consumer implication), KIZ original InpZ, KTR1-KTR13, RNC16, PGK1-PGK6. Earlier frontier descriptions are historical; this qualification is the current scope authority. Original whole-contract premises and quantifiers are retained. CURRENT VERSION0.49 QUALIFICATION: Construct one common GdEnd for each original datum and fixed iota: the same marked g1 must supply FinEnd, AutEnd3 and AutEnd2. UCE now proves all ordinary local/curve/coefficient consumers and the exact implication to FinA/AutB/LemW, including arbitrary entry maps and every original outer-corner extension. The producer still needs the actual common abelian-CM GdFrameEx and LinkedEndE with both rank-two and rank-three runs, full ModUAlt/PfromAut and the two actual seeds, followed by the original local/automorphic descent. At this pin the alternative upstream route gives GdEndAlt; its additional local/descent fields must be matched to the actual original points, so no name-only GdEndAlt-to-GdEnd identification is allowed. Keep the supplied InpZ, RigII, CNR3, finite PropGK and whole OneChar implication. Continue RH Cases I/V and the independent additional Dirichlet-GRH CaseIII/promodularity, SeedIV, tensor-factor and rank-four-chain tasks. Completed current producers or conditional consumers: UCE1-UCE19 (complete GdEnd consumer), OCH1-OCH30, KIZ original InpZ, KTR1-KTR13, RNC16, PGK1-PGK6. Earlier frontier descriptions are historical; this qualification is the current scope authority. Original whole-contract premises and quantifiers are retained.

\emph{Completed inputs (version 0.54).} RTC1–RTC25 (upstream universal coefficient rings, full preconnected\_\allowbreak rank\_\allowbreak two and FinPtDown), ACP1-ACP19 (full FinPtOf and abstract propagation principles; rank-specific suppliers retained), AFE1-AFE21 (full Gen entries, conditional frame joins, NRep descent and FinPt-to-FinEnd), LFC1-LFC23, CAF1-CAF11, MFC1-MFC14, UCE1-UCE19 (full original GdEnd consumer), OCH1-OCH30, KIZ original InpZ, KTR1-KTR13, RNC16, PGK1-PGK6

\emph{Remaining producer scope.} RTC supplies full original wpres/co\_\allowbreak B/gr\_\allowbreak A/cont\_\allowbreak u upstream for both ranks two and three, without FinA. Formal-unit twists separate integral weights over nilpotent test rings; finite monomial decompositions and fixed-weight filtration cofinality identify the actual entry and graded coefficient rings. It supplies the complete rank-two arbitrary-k-algebra triangular point correspondence, the original smooth/nonzero strict-cocycle producers, and full preconnected\_\allowbreak rank\_\allowbreak two. It also supplies full original FinPtDown at arbitrary rank and all allowed primes, and the exact ACP11-to-FinPtDown-to-AFE20 return. Remaining are actual loc\_\allowbreak cpl and LocFlat/LocRho charts, sharp component and dimension bounds, density/seed incidence, complete rank-two propagation, the rank-two-to-three corner and exceptional ideal, rank-three propagation and OffFin. Actual automorphic seeds, patching/modularity, classicality and the stronger Alt WitA/WitAQ outputs also remain. Keep rank-two and rank-three rings/ideals separate on the same AFE marked frame, with original P where used. ACP abstract propagation and FinPtOf are complete; completed generic gr\_\allowbreak A must replace the former opaque upstream input, while loc\_\allowbreak cpl and actual flatness are not inferred from it. The broad IsBar realization is proved directly; ACP’s diagonal IsLiftA shortcut is not silently enlarged. Host ell=3 and the original coefficient descent stay in scope.

\emph{Source.} \href{https://github.com/CocoSgt/riemann-hypothesis-lean/blob/7cbf7a0c32fca1badf05931879a37c3f8a2b03b4/src/Development/AutomorphicArithmetic/AutomorphicVarieties/AutomorphicCases/AutomorphyUnramifiedCoefficients.lean#L24-L27}{AutomorphyUnramifiedCoefficients}.

\emph{Source.} \href{https://github.com/CocoSgt/riemann-hypothesis-lean/blob/7cbf7a0c32fca1badf05931879a37c3f8a2b03b4/src/Development/AutomorphicArithmetic/AutomorphicVarieties/AutomorphicCases/AutomorphyUnramifiedCoefficients.lean#L12-L13}{FinA}.

\emph{Source.} \href{https://github.com/CocoSgt/riemann-hypothesis-lean/blob/7cbf7a0c32fca1badf05931879a37c3f8a2b03b4/src/Development/AutomorphicArithmetic/AutomorphicVarieties/AutomorphicCases/AutomorphyUnramifiedCoefficients.lean#L15-L17}{AutB}.

\emph{Source.} \href{https://github.com/CocoSgt/riemann-hypothesis-lean/blob/7cbf7a0c32fca1badf05931879a37c3f8a2b03b4/src/Development/AutomorphicArithmetic/AutomorphicVarieties/AutomorphicCases/AutomorphyUnramifiedCoefficients.lean#L19-L22}{LemW}.



\paragraph{6. TASK\_\allowbreak CASE3\_\allowbreak AUT\_\allowbreak JOIN — B — Complete the original Case-III automorphy/image join}

Install the actual Case-III ProModAlt, SeedIVAlt and shared PMExAlt with the inherited CohomologyHyp/Ferrero-Washington/HFit/SeedII facts into the original whole Case-III automorphy/image conclusion.

\emph{Source.} \href{https://github.com/CocoSgt/riemann-hypothesis-lean/blob/7cbf7a0c32fca1badf05931879a37c3f8a2b03b4/src/Development/AutomorphicArithmetic/AutomorphicVarieties/ModularityArgument/ThresholdArguments/DenseThresholdCases/PromodularityCaseThree.lean#L43-L44}{AutomorphyRankFourDense}.

\emph{Source.} \href{https://github.com/CocoSgt/riemann-hypothesis-lean/blob/7cbf7a0c32fca1badf05931879a37c3f8a2b03b4/src/Development/AutomorphicArithmetic/AutomorphicVarieties/ModularityArgument/ThresholdArguments/DenseThresholdCases/CuFinal.lean#L44-L44}{main4\_\allowbreak alt}.



\paragraph{7. TASK\_\allowbreak CASE3\_\allowbreak PROMOD\_\allowbreak JOIN — C — Reconstruct the actual Case-III chain-to-promodularity proof}

Use the original tensor-factor tower, ChainRankFour, ChainDesc4, seeded rank-four chain reduction, actual PMExAlt and SeedIVAlt to construct ProModAlt on every original Ctx3 star point at the same fixed iota, preserving all finite-coefficient, inertia and support maps.

\emph{Source.} \href{https://github.com/CocoSgt/riemann-hypothesis-lean/blob/7cbf7a0c32fca1badf05931879a37c3f8a2b03b4/src/Development/AutomorphicArithmetic/AutomorphicVarieties/ModularityArgument/ThresholdArguments/DenseThresholdCases/PromodularityCaseThreeOfLemmas.lean#L131-L153}{UpperTheory.V.Cu.promodularity\_\allowbreak case\_\allowbreak three\_\allowbreak of\_\allowbreak chain\_\allowbreak reduction}.

\emph{Source.} \href{https://github.com/CocoSgt/riemann-hypothesis-lean/blob/7cbf7a0c32fca1badf05931879a37c3f8a2b03b4/src/Development/AutomorphicArithmetic/AutomorphicVarieties/ModularityArgument/ThresholdArguments/DenseThresholdCases/PromodularityCaseThree.lean#L40-L42}{PromodularityCaseThreeDense}.



\paragraph{8. SEED\_\allowbreak IV\_\allowbreak ALT — C — Actual Case-III seed with prechosen residual datum}

Base Ctx3/Four tower, iota and Amb -> w,s before every later Dies tower -> compatible finite-coefficient seed with exact residual support/root entries.

\emph{Source.} \href{https://github.com/CocoSgt/riemann-hypothesis-lean/blob/7cbf7a0c32fca1badf05931879a37c3f8a2b03b4/src/Development/AutomorphicArithmetic/AutomorphicVarieties/ModularityArgument/ThresholdArguments/DenseThresholdCases/PromodularityCaseThree.lean#L19-L25}{SeedIVAlt}.



\paragraph{9. TENSOR\_\allowbreak FACTOR\_\allowbreak TOWER — B — Actual tensor-factor tower and characters}

Original finite-field, tame, factor and character constraints; source-level construction is sizeable and has not been reconstructed here.

\emph{Source.} \href{https://github.com/CocoSgt/riemann-hypothesis-lean/blob/7cbf7a0c32fca1badf05931879a37c3f8a2b03b4/src/Development/AutomorphicArithmetic/AutomorphicVarieties/ModularityArgument/ThresholdArguments/DenseThresholdCases/PromodularityCaseThreeOfLemmas.lean#L131-L153}{UpperTheory.V.Cu.promodularity\_\allowbreak case\_\allowbreak three\_\allowbreak of\_\allowbreak chain\_\allowbreak reduction}.



\paragraph{10. CHAIN\_\allowbreak RANK\_\allowbreak FOUR — C — Actual SomeIn-to-AllIn propagation}

Same source Space and PM ideal, residual point Good4, original bound 30(C+1)\textasciicircum{}2 <= d; geometry must preserve all identities.

\emph{Source.} \href{https://github.com/CocoSgt/riemann-hypothesis-lean/blob/7cbf7a0c32fca1badf05931879a37c3f8a2b03b4/src/Development/AutomorphicArithmetic/AutomorphicVarieties/ModularityArgument/ThresholdArguments/DenseThresholdCases/PromodularityCaseThree.lean#L26-L32}{ChainRankFour}.



\paragraph{11. CHAIN\_\allowbreak DESC4 — B — Original residual chain descent}

Finite residual descent retaining source Chain4 compatibility.

\emph{Source.} \href{https://github.com/CocoSgt/riemann-hypothesis-lean/blob/7cbf7a0c32fca1badf05931879a37c3f8a2b03b4/src/Development/AutomorphicArithmetic/AutomorphicVarieties/ModularityArgument/ThresholdArguments/DenseThresholdCases/PromodularityCaseThreeOfLemmas.lean#L131-L153}{UpperTheory.V.Cu.promodularity\_\allowbreak case\_\allowbreak three\_\allowbreak of\_\allowbreak chain\_\allowbreak reduction}.



\paragraph{12. RANK4\_\allowbreak CHAIN\_\allowbreak REDUCTION — C — Seeded reduction to the original residual chain}

Select T1 before every later Deep tower; retain one compatible seed and Chain4 for every original residual lift over Q.

\emph{Source.} \href{https://github.com/CocoSgt/riemann-hypothesis-lean/blob/7cbf7a0c32fca1badf05931879a37c3f8a2b03b4/src/Development/AutomorphicArithmetic/AutomorphicVarieties/ModularityArgument/ThresholdArguments/DenseThresholdCases/PromodularityCaseThree.lean#L33-L39}{RankFourChainReductionDense}.

\emph{Source.} \href{https://github.com/CocoSgt/riemann-hypothesis-lean/blob/7cbf7a0c32fca1badf05931879a37c3f8a2b03b4/src/Development/AutomorphicArithmetic/AutomorphicVarieties/ModularityArgument/ThresholdArguments/DenseThresholdCases/CuPChainIV.lean#L16-L82}{UpperTheory.V.Cu.chainIV\_\allowbreak alt}.



\paragraph{13. TASK\_\allowbreak COUNT\_\allowbreak V\_\allowbreak INSTALL — B — Complete the original CountV finite assembly}

From exactly the listed source host/cohomology inputs, complete countV\_\allowbreak of\_\allowbreak finiteness\_\allowbreak duality\_\allowbreak euler\_\allowbreak host, including original host/frame formation, normalized length, the CaseV guard and R1<=2\textasciicircum{}c2.

\emph{Completed inputs (version 0.54).} RTC full rank-two preconnected core and coefficient descent; actual propagation/counting inputs retained, ACP full finite-point and abstract propagation implications; actual rank-specific and counting inputs retained, AFE full Host entry/auxiliary frame/local witnesses including ell=3, LFC1-LFC23, CAF9-CAF11, MFC1-MFC14, Inherited CountV consumers

\emph{Remaining producer scope.} RTC supplies full original wpres/co\_\allowbreak B/gr\_\allowbreak A/cont\_\allowbreak u upstream for both ranks two and three, without FinA. Formal-unit twists separate integral weights over nilpotent test rings; finite monomial decompositions and fixed-weight filtration cofinality identify the actual entry and graded coefficient rings. It supplies the complete rank-two arbitrary-k-algebra triangular point correspondence, the original smooth/nonzero strict-cocycle producers, and full preconnected\_\allowbreak rank\_\allowbreak two. It also supplies full original FinPtDown at arbitrary rank and all allowed primes, and the exact ACP11-to-FinPtDown-to-AFE20 return. Remaining are actual loc\_\allowbreak cpl and LocFlat/LocRho charts, sharp component and dimension bounds, density/seed incidence, complete rank-two propagation, the rank-two-to-three corner and exceptional ideal, rank-three propagation and OffFin. Actual automorphic seeds, patching/modularity, classicality and the stronger Alt WitA/WitAQ outputs also remain. Keep rank-two and rank-three rings/ideals separate on the same AFE marked frame, with original P where used. ACP abstract propagation and FinPtOf are complete; completed generic gr\_\allowbreak A must replace the former opaque upstream input, while loc\_\allowbreak cpl and actual flatness are not inferred from it. The broad IsBar realization is proved directly; ACP’s diagonal IsLiftA shortcut is not silently enlarged. Host ell=3 and the original coefficient descent stay in scope. Complete the original F3 local coefficient bridges and F4 common-overfield readings, with valuation/norm/GoodCoef/iq and the entire coefficient field, and actual complex Hecke standard/no-Jordan plus arithmetic character data on the same form datum. Construct the actual base rows, host/frame lengths, nonzero chains and group ties, MultEq support, eigensystem/carrier/parts, all-embedding logLam, HD/d3b and NEb. Keep the stronger second DivRows3 scope distinct from the completed first DivDef3 and row-pair integrability. The full arithmetic entry and conditional common-frame construction do not by themselves install HOST\_\allowbreak NUMBER\_\allowbreak S, HOST\_\allowbreak FRAME\_\allowbreak LENGTH, TASK\_\allowbreak HOST\_\allowbreak VAN\_\allowbreak E\_\allowbreak INSTALL, TASK\_\allowbreak HOST\_\allowbreak XPP\_\allowbreak INSTALL or TASK\_\allowbreak COUNT\_\allowbreak V\_\allowbreak INSTALL.

\emph{Source.} \href{https://github.com/CocoSgt/riemann-hypothesis-lean/blob/7cbf7a0c32fca1badf05931879a37c3f8a2b03b4/src/Development/AutomorphicArithmetic/AutomorphicVarieties/HeightAndVanishing/Counts.lean#L69-L69}{CountV}.

\emph{Source.} \href{https://github.com/CocoSgt/riemann-hypothesis-lean/blob/7cbf7a0c32fca1badf05931879a37c3f8a2b03b4/src/Development/AutomorphicArithmetic/AutomorphicVarieties/ModularityArgument/ThresholdArguments/DenseThresholdCases/CountVOf.lean#L33-L60}{UpperTheory.V.UnitaryThreeInertia.countV\_\allowbreak of\_\allowbreak finiteness\_\allowbreak duality\_\allowbreak euler\_\allowbreak host}.



\paragraph{14. THM\_\allowbreak S6\_\allowbreak S — C — Original unitary-three ordinary support and automorphy}

Original host/frame and fixed iota -> one marked g1 -> every universal R,u -> weight-zero finiteness and every Type-I BPrime\_\allowbreak s point.

\emph{Source.} \href{https://github.com/CocoSgt/riemann-hypothesis-lean/blob/7cbf7a0c32fca1badf05931879a37c3f8a2b03b4/src/Development/AutomorphicArithmetic/AutomorphicVarieties/ModularityArgument/ThresholdArguments/DenseThresholdCases/ThmS6CoreS.lean#L17-L26}{ThmS6\_\allowbreak s}.

\emph{Source.} \href{https://github.com/CocoSgt/riemann-hypothesis-lean/blob/7cbf7a0c32fca1badf05931879a37c3f8a2b03b4/src/Development/AutomorphicArithmetic/AutomorphicVarieties/ModularityArgument/ThresholdArguments/DenseThresholdCases/ThmS6CoreS.lean#L12-L16}{BPrime\_\allowbreak s}.



\paragraph{15. HOST\_\allowbreak FRAME\_\allowbreak LENGTH — B — Original Case-V host/frame and normalized-length installation}

Actual HostData, frame, local-condition equality, abstract count and normalized original Selmer length must all refer to the same objects.

\emph{Completed inputs (version 0.54).} AFE/LFC full entry and conditional common frame, CAF and MFC, DNE weight/host-degree dictionary at supplied tables

\emph{Remaining producer scope.} Complete the original F3 local coefficient bridges and F4 common-overfield readings, with valuation/norm/GoodCoef/iq and the entire coefficient field, and actual complex Hecke standard/no-Jordan plus arithmetic character data on the same form datum. Construct the actual base rows, host/frame lengths, nonzero chains and group ties, MultEq support, eigensystem/carrier/parts, all-embedding logLam, HD/d3b and NEb. Keep the stronger second DivRows3 scope distinct from the completed first DivDef3 and row-pair integrability. The full arithmetic entry and conditional common-frame construction do not by themselves install HOST\_\allowbreak NUMBER\_\allowbreak S, HOST\_\allowbreak FRAME\_\allowbreak LENGTH, TASK\_\allowbreak HOST\_\allowbreak VAN\_\allowbreak E\_\allowbreak INSTALL, TASK\_\allowbreak HOST\_\allowbreak XPP\_\allowbreak INSTALL or TASK\_\allowbreak COUNT\_\allowbreak V\_\allowbreak INSTALL.

\emph{Source.} \href{https://github.com/CocoSgt/riemann-hypothesis-lean/blob/7cbf7a0c32fca1badf05931879a37c3f8a2b03b4/src/Development/AutomorphicArithmetic/AutomorphicVarieties/ModularityArgument/ThresholdArguments/DenseThresholdCases/CountVOf.lean#L33-L60}{UpperTheory.V.UnitaryThreeInertia.countV\_\allowbreak of\_\allowbreak finiteness\_\allowbreak duality\_\allowbreak euler\_\allowbreak host}.



\paragraph{16. HOST\_\allowbreak NUMBER\_\allowbreak S — C — Uniform original-host rank-three catalogue}

One exponent before all closures, coefficient rings, levels, characters, frames, embeddings and marked points; one separating trace function controls all later finite pairwise nonisomorphic collections.

\emph{Completed inputs (version 0.54).} AFE full Host entry at original scope, LFC full conditional common frame, CAF and MFC, Inherited host-count consumers

\emph{Remaining producer scope.} Complete the original F3 local coefficient bridges and F4 common-overfield readings, with valuation/norm/GoodCoef/iq and the entire coefficient field, and actual complex Hecke standard/no-Jordan plus arithmetic character data on the same form datum. Construct the actual base rows, host/frame lengths, nonzero chains and group ties, MultEq support, eigensystem/carrier/parts, all-embedding logLam, HD/d3b and NEb. Keep the stronger second DivRows3 scope distinct from the completed first DivDef3 and row-pair integrability. The full arithmetic entry and conditional common-frame construction do not by themselves install HOST\_\allowbreak NUMBER\_\allowbreak S, HOST\_\allowbreak FRAME\_\allowbreak LENGTH, TASK\_\allowbreak HOST\_\allowbreak VAN\_\allowbreak E\_\allowbreak INSTALL, TASK\_\allowbreak HOST\_\allowbreak XPP\_\allowbreak INSTALL or TASK\_\allowbreak COUNT\_\allowbreak V\_\allowbreak INSTALL.

\emph{Source.} \href{https://github.com/CocoSgt/riemann-hypothesis-lean/blob/7cbf7a0c32fca1badf05931879a37c3f8a2b03b4/src/Development/AutomorphicArithmetic/AutomorphicVarieties/ModularityArgument/ThresholdArguments/DenseThresholdCases/HostNumberS.lean#L18-L28}{HostNumber\_\allowbreak s}.



\paragraph{17. TASK\_\allowbreak HOST\_\allowbreak VAN\_\allowbreak E\_\allowbreak INSTALL — B — Transport K/duality to literal HostVanE}

Use the already supplied LemmaK/local-global duality and finiteness on the actual CountV host/frame and exact original coefficient modules.

\emph{Source.} \href{https://github.com/CocoSgt/riemann-hypothesis-lean/blob/7cbf7a0c32fca1badf05931879a37c3f8a2b03b4/src/Development/AutomorphicArithmetic/AutomorphicVarieties/UnitaryThreeDeformations/HostNumber.lean#L9-L12}{HostVanE}.

\emph{Source.} \href{https://github.com/CocoSgt/riemann-hypothesis-lean/blob/7cbf7a0c32fca1badf05931879a37c3f8a2b03b4/src/Development/AutomorphicArithmetic/AutomorphicVarieties/LevelCohomology/SoHostVanE.lean#L11-L18}{UpperTheory.V.X.So.hostVanE}.



\paragraph{18. TASK\_\allowbreak HOST\_\allowbreak XPP\_\allowbreak INSTALL — B — Transport CorX to literal HostXpp}

Preserve Psi\textasciicircum{}-2 eta, the prime-to-ell quotient, actual local conditions and the original CountV host while installing BR7.

\emph{Source.} \href{https://github.com/CocoSgt/riemann-hypothesis-lean/blob/7cbf7a0c32fca1badf05931879a37c3f8a2b03b4/src/Development/AutomorphicArithmetic/AutomorphicVarieties/UnitaryThreeDeformations/HostNumber.lean#L14-L17}{HostXpp}.

\emph{Source.} \href{https://github.com/CocoSgt/riemann-hypothesis-lean/blob/7cbf7a0c32fca1badf05931879a37c3f8a2b03b4/src/Development/AutomorphicArithmetic/AutomorphicVarieties/LevelCohomology/Xpp.lean#L11-L84}{UpperTheory.V.X.xpp}.



\paragraph{19. DESIGN\_\allowbreak CHOICES\_\allowbreak COMMON — C — Construct the original common design family and bundle}

For the one target HostExW, construct the original compatible DesignChoices on Cdes=WithTw(WithSm Cprime), where Cprime=WithChN(A intersect target class), A is the one Built construction class, and cT(eps,y)=cTmin(y). Retain every original field, tuple, word, metric and common bundle duty. DSZ10 and RMP11 give size/rank properties of class members, not member existence. The original RecExP/RecAtG and common object/row route remains coupled. Construct ONE original common design family R and its recursive/geometric half with original Inputs/Pins/Holds/ExAt/InGeo and required CoreP. PKT now reconstructs full exists\_\allowbreak pkT, rowPk\_\allowbreak nonempty and exists\_\allowbreak moved with Copy/Coh on those original given inputs and a supplied fitting later stage: one catalogue, both ChT, complete LfT, actual integral lattices and all metric/measure/period constants are retained; one numerical record precedes every eta and V/S/N/G. RDR supplies one Side0 before all readers, ReadPk/pins and full ReadL with exact o.eta=q.eta. The initial actual geometry, Fits/division/central conditions, primitive rational tests and eta agreement must still be produced. Inherited CSR supplies finite columns, one ReadC before later twists, period/evaluation and scalar/error transports, RowSideDetermined and same-R Bundle/BTw/BWit/BunU consumers. RPI supplies exact RowPairInt for actual pairs and finite translates with original mu\_\allowbreak b/fd0; first literal DivDef3 at the same updated R,F,Y,Z,C,sEE follows from its retained original readings, Pins/Holds, stored nonzero IsItW chain, five real-group predicates and group tie. Remaining common fields include host/frame, MultEq nonzero-part support, that original nonzero chain/group tie, eigensystem/carrier/parts and all-embedding logLam, HD/d3b, NEb, actual rational tests and base rows. Retain every simultaneous family hypothesis and eventual threshold. Stronger second DivRows3 remains separate; universal quotient-height CuspRapidDecay and actual common member existence are not inferred from the comparison constructors. Version 0.41 supplies the finite phase-correct rational theta route: exact coefficient extension, split representatives, square-law/PhiIs, Pins from CoversAt, class-sum/point-identity transfer to OneMapOverL, actual-good-projector RatG and both-level LiftTestB. Its original higher covers, all-tuple class-sum uniqueness including phase rescaling, one PointIdentCore before O, base ClassSumsSpan and geometric LatFull remain explicit producers. MLat is available from the original LatFull input. This sharpens the rational-test work without claiming actual common geometry, Fits or eta agreement. Version0.46 proves finite generation of the original whole integral cusp lattice by actual integral-frame evaluation. It also proves the exact identity-cover, analytic unfolding, original-field coefficient descent, global clearing-scalar and both-level lattice-transport implications. Full spanning retains the supplied original IsNab connection and the five explicit principal-chart, level-system, raw-Gaussian, short-rule/dictionary and ramified-jet producers. The higher tuple covers, actual common class sums, PointIdentCore and ClassSumsSpan remain separate. CURRENT VERSION0.47 QUALIFICATION: Construct one original common geometric family and the higher CoversAt/all-tuple class sums and uniqueness with original phase rescalings, one PointIdentCore with common lambda,z, and base ClassSumsSpan. CNC now supplies the full original IsNab; PCL supplies the conditional trace-dual principal-lattice and quantitative normalization. Actual uniformization/HasPrinc and the five principal-chart, finite-level-system, Gaussian/reproduction, short-rule/simultaneous-dictionary and ramified-jet producer groups remain. Keep all original both-level, coefficient-field, Fits and eta agreements. Completed new original producers: ICL1-ICL21, CNC1-CNC16, PCL1-PCL12. Any earlier narrative naming these completed producers as outstanding is historical; this qualification is the current scope authority. CURRENT VERSION0.48 QUALIFICATION: Construct the original finite holomorphic FamilyUnif/N1c, normalized theta families and local lifting, then the common principal charts and finite-level theta systems. PUN now supplies the exact fixed-space models, all-family center lattices, A=1 exponent matching and individual primitive points with the full point-equality criterion; CNC supplies the full canonical connection and PCL the quantitative trace-dual normalization. Complete higher CoversAt and all-tuple class sums/uniqueness with the original phase rescalings, one PointIdentCore with common lambda,z, and base ClassSumsSpan. The Gaussian/reproduction, short-rule/simultaneous-dictionary and ramified-jet producers, and all original both-level, coefficient-field, Fits and eta agreements remain explicit. Completed current producers or conditional consumers: ICL1-ICL21, CNC1-CNC16, PCL1-PCL12, PUN1-PUN24, GCF1-GCF18 at their stated scopes. Earlier frontier descriptions are historical; this qualification is the current scope authority. Original whole-contract premises and quantifiers are retained. CURRENT VERSION0.49 QUALIFICATION: Construct the adjoining simultaneous principal framed points and test-ring dictionary. Full finite holomorphic FamilyUnif/N1c, normalized theta families and local lifts are now supplied by HUF. The next bounded target is allTest\_\allowbreak of\_\allowbreak secBaseChange: represent FrPt by the Lie-eigenframe bundles and theta-section frame line, use the full n=1 SecBaseChangeSt on every nonflat coefficient base change, and prove finite-generated subalgebra descent with the exact TRing/FromTest transport. Then supply the full FramesOK/CoverTri/FibTri and IsUnf/hasPrinc package, including adapted mod-eight symplectic basis, normalized theta multiplier/weight and arithmetic invariance. Continue higher CoversAt, all-tuple class sums and uniqueness with original phase rescalings, one PointIdentCore with common lambda,z, base ClassSumsSpan, Gaussian/reproduction, short-rule dictionaries and ramified jets, retaining both-level, coefficient-field, Fits and eta agreements. Common-family nonemptiness remains distinct from compatibility of an already supplied member. Completed current producers or conditional consumers: ICL1-ICL21, CNC1-CNC16, PCL1-PCL12, PUN1-PUN24, HUF1-HUF21 (full original FamilyUnif/N1c), GCF1-GCF18 and WRS1-WRS10 at their stated scopes. Earlier frontier descriptions are historical; this qualification is the current scope authority. Original whole-contract premises and quantifiers are retained.

\emph{Completed inputs (version 0.54).} ASL1–ASL21 (actual section localization and full original arbitrary-base Perf producer), FCD1-FCD18 (full homogeneous-chart descent, graded presentation and Noetherian cube model), GSL1-GSL26 (full conditional section laws, universal basis and projective-model algebra), NFL1-NFL22, TDR1-TDR13, NFR1-NFR13, HUF1-HUF21, PUN1-PUN24, ICL1-ICL21, CNC1-CNC16, PCL1-PCL12, GCF/WRS at original supplied-member scope

\emph{Remaining producer scope.} ASL supplies the full original actual twisted-section power-killing and extension theorems, graded localization and dehomogenization, and all-vertex section-cube comparison. Affine line-bundle sections are finite projective over arbitrary rings and the actual chart algebras are smooth and flat. Joining these with FCD18 and GSL8–10 supplies the full original universal Perf producer, nonempty\_\allowbreak perf\_\allowbreak of\_\allowbreak smooth\_\allowbreak isProjBy, on the given scheme and line bundle over arbitrary R. One Noetherian model, one finite-projective complex and the actual section comparison precede every later coefficient algebra and pullback square. Do not reopen the algebraic or actual geometric perfect-complex model. The separate PolBaseChangeSt, Serre H1 with its all-cover conclusion, multiplication/invariant descent, geometric integrality, principal rank one, torsion section decomposition/growth and theorem-of-square generators remain. The next bounded geometry packet is the full original Serre H1 theorem: one finite homogeneous submodule for all degrees, a uniform large-degree Cech bound, and the reverse chosen-cover-to-all-cover argument. Supply the independent P,V,R,T producers before the relevant GSL26/NFL joins. Original Moduli and HUF/PUN foundations remain. Later finite-level systems, Gaussian/reproduction, simultaneous dictionaries, ramified jets, all-tuple class sums, phase-sensitive uniqueness, common PointIdentCore and base ClassSumsSpan remain, retaining both levels, whole fields, Fits and eta agreement.



\paragraph{20. LOCAL\_\allowbreak NBUILT\_\allowbreak COMMON — C — Actual local specialization on the common selected family}

Construct, with the required closed-class property on the same Des(C,3), one eventual family of actual NBuilt certificates for the original Fm.D(x,i), retaining the certificate before forgetting to TSel. HLP now reconstructs full HeckeLift0, all-index commutation, HeckeStd0 and HeckeIntTw from the original MtInput and complex Hecke facts, in the full ambient GenAt range k<=nV. Every maximal ideal is an actual supported residual character class; literal MtSub is the O-span of projections of the original MtLat generators, is finite free, and localizes onto the exact NCov.part. Literal Psi has its exact range and off-space projection error, both original span assertions hold, and ActEx/ActHK are constructed on the same lattice under original semisimplicity. Original MjLat relation quotients, universal coefficient-map relation fields, selected tame generators, coefficient embeddings and all identity transports give the complete nonempty\_\allowbreak spec constructor on its full printed input list. A single compositum combines the four original coefficient requirements with big fixed beforehand. Whole ndSpec\_\allowbreak exists still needs actual common geometric/choice MtInput, complex HeckeStdPlain and the no-Jordan producer, original F3 bridges with valuation/GoodCoef/iq, original F4 common-overfield reading of the entire coefficient field, and exact arithmetic character/reciprocity/admissibility suppliers. Original AutSpec, Galois specialization and the remaining NBuilt fields on that same installation also remain. Do not replace the actual projected lattice or maps by abstract objects with the same dimensions.

\emph{Completed inputs (version 0.54).} ASL full actual perfect-complex model; RTC common-frame arithmetic inputs; AOS exact unchanged infinity normalization, FCD full algebraic model, ACP abstract arithmetic suppliers, GMR same-setting measure normalization, GSL/NFL full conditional section laws and uniformization, AFE/LFC full entry and conditional common frame, TDR/NFR and CAF/MFC, WRS selected built-class extraction at membership

\emph{Remaining producer scope.} Continue the distinct retained-defect, period-height and normalized spectral producers, and the original compatible second-layer selection on one common arithmetic/geometric family. GLOBAL\_\allowbreak LOCAL\_\allowbreak GAIN remains fully closed at its original v0.49 scope. Complete local comparison theorems do not themselves construct the whole common family. Common design existence, NBUILT compatibility, original true-column records, shared fields/embeddings, Fits and eta agreement, simultaneous rows and eventual-family thresholds remain explicit. A stronger historical WPt/TSel construction may serve other producers but is not needed to reopen the proved selected-class local-gain estimate. ASL supplies the full original actual twisted-section power-killing and extension theorems, graded localization and dehomogenization, and all-vertex section-cube comparison. Affine line-bundle sections are finite projective over arbitrary rings and the actual chart algebras are smooth and flat. Joining these with FCD18 and GSL8–10 supplies the full original universal Perf producer, nonempty\_\allowbreak perf\_\allowbreak of\_\allowbreak smooth\_\allowbreak isProjBy, on the given scheme and line bundle over arbitrary R. One Noetherian model, one finite-projective complex and the actual section comparison precede every later coefficient algebra and pullback square. Do not reopen the algebraic or actual geometric perfect-complex model. The separate PolBaseChangeSt, Serre H1 with its all-cover conclusion, multiplication/invariant descent, geometric integrality, principal rank one, torsion section decomposition/growth and theorem-of-square generators remain. The next bounded geometry packet is the full original Serre H1 theorem: one finite homogeneous submodule for all degrees, a uniform large-degree Cech bound, and the reverse chosen-cover-to-all-cover argument. Supply the independent P,V,R,T producers before the relevant GSL26/NFL joins. Original Moduli and HUF/PUN foundations remain. Later finite-level systems, Gaussian/reproduction, simultaneous dictionaries, ramified jets, all-tuple class sums, phase-sensitive uniqueness, common PointIdentCore and base ClassSumsSpan remain, retaining both levels, whole fields, Fits and eta agreement.



\paragraph{21. SECOND\_\allowbreak LAYER\_\allowbreak COMPATIBLE\_\allowbreak SELECTION — C — Original second-layer selector and class duties}

For the same A, supply Duty on A and a Closed common original TSel family whose selected records are the original true-column rec7 records (or have the entire dependent tuple transported), with their real witness and original coefficient-field root data. Existence/compatibility remains arithmetic; merely abstractly isomorphic fields or independently selected columns do not suffice. CSR21–22 now construct the same-R Bundle/BTw/BWit/BunU outputs from the exact selected realization/row-side packet and NEb, with the original EvP/ParEv quantifier order. Actual common selection, the five stated realization packets and NEb remain upstream; none is replaced by an independently chosen compatible record. Version 0.38 closes the selected analytic extension and reconstructs the first updated DivDef3 consumer; it leaves the original common selection, geometric readings, rational tests, base rows, host/parts/NEb and exact retained DI premises upstream, with unchanged EvP/ParEv thresholds. Version 0.39 supplies the original given-data reader constructors and exact RowPairInt/first DivDef3 on original premises. It leaves one actual common geometric/comparison/test member, eta agreement, compatible selection, rational tests, base rows, host/parts/NEb and original EvP/ParEv thresholds upstream. Version 0.40 also supplies original comparison/RowPk/fitting-stage Moved algorithms and exact common data transport. The actual initial geometric member, fitting tests, eta agreement and compatible second-layer selector with original eventual thresholds remain upstream. Version 0.41 supplies the finite phase-correct rational theta route: exact coefficient extension, split representatives, square-law/PhiIs, Pins from CoversAt, class-sum/point-identity transfer to OneMapOverL, actual-good-projector RatG and both-level LiftTestB. Its original higher covers, all-tuple class-sum uniqueness including phase rescaling, one PointIdentCore before O, base ClassSumsSpan and geometric LatFull remain explicit producers. MLat is available from the original LatFull input. This sharpens the rational-test work without claiming actual common geometry, Fits or eta agreement. CURRENT VERSION0.49 QUALIFICATION: Continue the distinct global denominator, period/height, normalized spectral and second-layer compatible-selection producers, together with existence of the original common arithmetic/geometric family. GLOBAL\_\allowbreak LOCAL\_\allowbreak GAIN and its actual GroupCompatibilityFamily supplier are now complete on the unchanged selected built class; no additional WPt/TSel realization is needed for that estimate. Do not reopen local first-growth, mean, auxiliary admission, selected D3 tops/Adm2F, the single GIs2 stage or the exact comparison/kernel transport. A full RecT or TSel may remain useful for other tasks, but is not claimed by this local-gain closure. Completed current producers or conditional consumers: GCF1-GCF18 at their stated scopes, WRS1-WRS10 original selected-class compatibility/local gain. Earlier frontier descriptions are historical; this qualification is the current scope authority. Original whole-contract premises and quantifiers are retained.

\emph{Completed inputs (version 0.54).} WIS actual inclusion and invariant completeness; AOS exact scalar; ASL/RTC full named nested producers, WLD full generic lattice constructors, GMR actual-setting normalization, FCD/ACP exact nested suppliers, WRS1-WRS10 original selected-class compatibility/local gain, DNE full denominator consumer, DFD/SSC canonical-data and signature comparisons, GSL/NFL and AFE/LFC conditional constructors, Inherited exact same-record realization/transport consumers

\emph{Remaining producer scope.} Continue the distinct retained-defect, period-height and normalized spectral producers, and the original compatible second-layer selection on one common arithmetic/geometric family. GLOBAL\_\allowbreak LOCAL\_\allowbreak GAIN remains fully closed at its original v0.49 scope. Complete local comparison theorems do not themselves construct the whole common family. Common design existence, NBUILT compatibility, original true-column records, shared fields/embeddings, Fits and eta agreement, simultaneous rows and eventual-family thresholds remain explicit. A stronger historical WPt/TSel construction may serve other producers but is not needed to reopen the proved selected-class local-gain estimate.



\paragraph{22. BASE\_\allowbreak THETA\_\allowbreak READINGS — C — Actual base reading on the common bundle}

For every family in the original class, eventually on the SAME selected Bundle and for each i<m, supply the actual R,F,Read,ReadL, weak models/Weil data and phase-one pins. Identify their LiftTie.theta, mtRead, mjRead and ThDef3 basis with that Bundle.P.A3, P.Zs, P.S0 and its stored maps. Exact field, embeddings, metrics, measures, row family and weak choices are retained. This remains one part of the coupled BTH12 producer. PKT now supplies the complete original comparison packet, RowPk and fitting-stage Moved with Copy/Coh on original given geometry, Inputs/Pins/Holds/ExAt/InGeo and Fits. Its ChT/LfT/Reps, actual lattice carriers, original Haar and norm constants, period maps, representatives and all-place RrT are constructed together; its one numerical record precedes every eta/V/S/N/G. RDR gives the original level-zero side, moved-Read and ReadPk/pin constructors on those objects, with one Side0 before every reader and all twelve Read ties. Full ReadL retains the exact o.eta(j,v)=q.eta(j,v) guard for every 1<=j<=m and place v; arbitrary Over does not imply it. Inherited CSR fixes one ReadC before all TwArgs, transports both metrics and all thirteen reading maps, and gives actual period/evaluation with kappa=(c*nP)\textasciicircum{}(-1/2). RPI supplies exact row-pair integrability, so first updated DivDef3 follows from the retained original Pins/Holds, nonzero chain, five real-group predicates and group tie. Initial theta extension remains complete. Remaining work is actual initial common geometry and its original predicates, Fits/primitive rational tests, eta agreement and outer family/eventual/same-Bundle duties, reFace compatibility, MultEq support, original chain/group data, actual base rows and any stronger DivRows3 theta square. No new compatible reading, family threshold or Haar normalization is substituted. Version 0.41 supplies the finite phase-correct rational theta route: exact coefficient extension, split representatives, square-law/PhiIs, Pins from CoversAt, class-sum/point-identity transfer to OneMapOverL, actual-good-projector RatG and both-level LiftTestB. Its original higher covers, all-tuple class-sum uniqueness including phase rescaling, one PointIdentCore before O, base ClassSumsSpan and geometric LatFull remain explicit producers. MLat is available from the original LatFull input. This sharpens the rational-test work without claiming actual common geometry, Fits or eta agreement.

\emph{Completed inputs (version 0.54).} ASL1–ASL21 (actual section localization and full original arbitrary-base Perf producer), FCD1-FCD18 (full homogeneous-chart descent, graded presentation and Noetherian cube model), GSL1-GSL26, NFL1-NFL22, TDR1-TDR13, NFR1-NFR13, HUF1-HUF21, PUN1-PUN24, CNC1-CNC16, Inherited actual-reading consumers at original inputs

\emph{Remaining producer scope.} ASL supplies the full original actual twisted-section power-killing and extension theorems, graded localization and dehomogenization, and all-vertex section-cube comparison. Affine line-bundle sections are finite projective over arbitrary rings and the actual chart algebras are smooth and flat. Joining these with FCD18 and GSL8–10 supplies the full original universal Perf producer, nonempty\_\allowbreak perf\_\allowbreak of\_\allowbreak smooth\_\allowbreak isProjBy, on the given scheme and line bundle over arbitrary R. One Noetherian model, one finite-projective complex and the actual section comparison precede every later coefficient algebra and pullback square. Do not reopen the algebraic or actual geometric perfect-complex model. The separate PolBaseChangeSt, Serre H1 with its all-cover conclusion, multiplication/invariant descent, geometric integrality, principal rank one, torsion section decomposition/growth and theorem-of-square generators remain. The next bounded geometry packet is the full original Serre H1 theorem: one finite homogeneous submodule for all degrees, a uniform large-degree Cech bound, and the reverse chosen-cover-to-all-cover argument. Supply the independent P,V,R,T producers before the relevant GSL26/NFL joins. Original Moduli and HUF/PUN foundations remain. Later finite-level systems, Gaussian/reproduction, simultaneous dictionaries, ramified jets, all-tuple class sums, phase-sensitive uniqueness, common PointIdentCore and base ClassSumsSpan remain, retaining both levels, whole fields, Fits and eta agreement. Complete the original F3 local coefficient bridges and F4 common-overfield readings, with valuation/norm/GoodCoef/iq and the entire coefficient field, and actual complex Hecke standard/no-Jordan plus arithmetic character data on the same form datum. Construct the actual base rows, host/frame lengths, nonzero chains and group ties, MultEq support, eigensystem/carrier/parts, all-embedding logLam, HD/d3b and NEb. Keep the stronger second DivRows3 scope distinct from the completed first DivDef3 and row-pair integrability. The full arithmetic entry and conditional common-frame construction do not by themselves install HOST\_\allowbreak NUMBER\_\allowbreak S, HOST\_\allowbreak FRAME\_\allowbreak LENGTH, TASK\_\allowbreak HOST\_\allowbreak VAN\_\allowbreak E\_\allowbreak INSTALL, TASK\_\allowbreak HOST\_\allowbreak XPP\_\allowbreak INSTALL or TASK\_\allowbreak COUNT\_\allowbreak V\_\allowbreak INSTALL.



\paragraph{23. BASE\_\allowbreak THETA\_\allowbreak RATIONAL\_\allowbreak TESTS — C — Simultaneous original base rational tests}

On the SAME original Bundle/reading/theta map, for every i<m supply one finite separating complex-linear test family whose values lie in the fixed embedded coefficient field on theta of EVERY original rational source vector and on EVERY stored target-coordinate basis vector. BTH8/BTH11 use this same LiftTestB for rational adjoint image and matrix coordinates. RTG now reconstructs prescribed-field coefficient extension, original split class representatives/RepsOK, the self-dual split-lattice square law and finite residue PhiIs; every Pins field follows from the actual CoversAt. With the base coefficient record literally fixed, exact class-sum existence/uniqueness including phase-rescaled O, PointIdentCore with common lambda,z before every O, and base ClassSumsSpan imply all-tuple finite lift span and OneMapOverL. The phase is exactly cj(twAt embLL A3.A.v0)(R.phs i). RatG uses the actual good projector eOp(y), and full LiftTestB follows at BOTH original levels. These finite implication proofs are complete. Remaining original producers are the prescribed higher covers, those actual class sums/uniqueness, PointIdentCore, base ClassSumsSpan and the base/common geometric datum; exact MLat is supplied only from original LatFull, whose geometric production remains separate. Finite dimension or an unrelated rational theta model cannot replace these same-object identities. Original outer eventual-family/Bundle quantifiers remain. Version0.46 proves finite generation of the original whole integral cusp lattice by actual integral-frame evaluation. It also proves the exact identity-cover, analytic unfolding, original-field coefficient descent, global clearing-scalar and both-level lattice-transport implications. Full spanning retains the supplied original IsNab connection and the five explicit principal-chart, level-system, raw-Gaussian, short-rule/dictionary and ramified-jet producers. The higher tuple covers, actual common class sums, PointIdentCore and ClassSumsSpan remain separate. CURRENT VERSION0.47 QUALIFICATION: Construct one original common geometric family and the higher CoversAt/all-tuple class sums and uniqueness with original phase rescalings, one PointIdentCore with common lambda,z, and base ClassSumsSpan. CNC now supplies the full original IsNab; PCL supplies the conditional trace-dual principal-lattice and quantitative normalization. Actual uniformization/HasPrinc and the five principal-chart, finite-level-system, Gaussian/reproduction, short-rule/simultaneous-dictionary and ramified-jet producer groups remain. Keep all original both-level, coefficient-field, Fits and eta agreements. Completed new original producers: ICL1-ICL21, CNC1-CNC16, PCL1-PCL12. Any earlier narrative naming these completed producers as outstanding is historical; this qualification is the current scope authority. CURRENT VERSION0.48 QUALIFICATION: Construct the original finite holomorphic FamilyUnif/N1c, normalized theta families and local lifting, then the common principal charts and finite-level theta systems. PUN now supplies the exact fixed-space models, all-family center lattices, A=1 exponent matching and individual primitive points with the full point-equality criterion; CNC supplies the full canonical connection and PCL the quantitative trace-dual normalization. Complete higher CoversAt and all-tuple class sums/uniqueness with the original phase rescalings, one PointIdentCore with common lambda,z, and base ClassSumsSpan. The Gaussian/reproduction, short-rule/simultaneous-dictionary and ramified-jet producers, and all original both-level, coefficient-field, Fits and eta agreements remain explicit. Completed current producers or conditional consumers: ICL1-ICL21, CNC1-CNC16, PCL1-PCL12, PUN1-PUN24. Earlier frontier descriptions are historical; this qualification is the current scope authority. Original whole-contract premises and quantifiers are retained. CURRENT VERSION0.49 QUALIFICATION: Construct the adjoining simultaneous principal framed points and test-ring dictionary. Full finite holomorphic FamilyUnif/N1c, normalized theta families and local lifts are now supplied by HUF. The next bounded target is allTest\_\allowbreak of\_\allowbreak secBaseChange: represent FrPt by the Lie-eigenframe bundles and theta-section frame line, use the full n=1 SecBaseChangeSt on every nonflat coefficient base change, and prove finite-generated subalgebra descent with the exact TRing/FromTest transport. Then supply the full FramesOK/CoverTri/FibTri and IsUnf/hasPrinc package, including adapted mod-eight symplectic basis, normalized theta multiplier/weight and arithmetic invariance. Continue higher CoversAt, all-tuple class sums and uniqueness with original phase rescalings, one PointIdentCore with common lambda,z, base ClassSumsSpan, Gaussian/reproduction, short-rule dictionaries and ramified jets, retaining both-level, coefficient-field, Fits and eta agreements. Completed current producers or conditional consumers: ICL1-ICL21, CNC1-CNC16, PCL1-PCL12, PUN1-PUN24, HUF1-HUF21 (full original FamilyUnif/N1c and theta family). Earlier frontier descriptions are historical; this qualification is the current scope authority. Original whole-contract premises and quantifiers are retained.

\emph{Completed inputs (version 0.54).} ASL1–ASL21 (actual section localization and full original arbitrary-base Perf producer), FCD1-FCD18 (full homogeneous-chart descent, graded presentation and Noetherian cube model), GSL1-GSL26, NFL1-NFL22, TDR1-TDR13, NFR1-NFR13, HUF1-HUF21, PUN1-PUN24, ICL1-ICL21, CNC1-CNC16, PCL1-PCL12, Inherited both-level RatG/LiftTestB consumers

\emph{Remaining producer scope.} ASL supplies the full original actual twisted-section power-killing and extension theorems, graded localization and dehomogenization, and all-vertex section-cube comparison. Affine line-bundle sections are finite projective over arbitrary rings and the actual chart algebras are smooth and flat. Joining these with FCD18 and GSL8–10 supplies the full original universal Perf producer, nonempty\_\allowbreak perf\_\allowbreak of\_\allowbreak smooth\_\allowbreak isProjBy, on the given scheme and line bundle over arbitrary R. One Noetherian model, one finite-projective complex and the actual section comparison precede every later coefficient algebra and pullback square. Do not reopen the algebraic or actual geometric perfect-complex model. The separate PolBaseChangeSt, Serre H1 with its all-cover conclusion, multiplication/invariant descent, geometric integrality, principal rank one, torsion section decomposition/growth and theorem-of-square generators remain. The next bounded geometry packet is the full original Serre H1 theorem: one finite homogeneous submodule for all degrees, a uniform large-degree Cech bound, and the reverse chosen-cover-to-all-cover argument. Supply the independent P,V,R,T producers before the relevant GSL26/NFL joins. Original Moduli and HUF/PUN foundations remain. Later finite-level systems, Gaussian/reproduction, simultaneous dictionaries, ramified jets, all-tuple class sums, phase-sensitive uniqueness, common PointIdentCore and base ClassSumsSpan remain, retaining both levels, whole fields, Fits and eta agreement.



\paragraph{24. SLR\_\allowbreak BASE\_\allowbreak ROWS\_\allowbreak PRODUCER — C — Actual base row identities on that same bundle}

For the SAME selected Bundle, common original fields and readings used by the other three BTH12 inputs, supply actual RowsHold0\_\allowbreak 3 uniformly eventually for every family. Its original primitive row realization, spectrum/type and basis identifications must be retained. BTH supplies a theta consumer only; it does not construct the base row identity. No separate witness may be chosen merely to satisfy this row requirement.



\clearpage
\section{Actual sections, homogeneous localization and the smooth-projective complex}
\label{sec:actual-section-localization-v054}

\status{This chapter constructs the localization and extension maps for
the actual line-bundle sections on the original scheme.  It proves the
full original twisted-section power-killing and extension theorems,
the dehomogenization and section-cube comparisons, and finite
presentation and projectivity of the actual affine section modules.
Together with FCD18 and GSL8--GSL10, these give the original
\texttt{nonempty\_\allowbreak{}perf\_\allowbreak{}of\_\allowbreak{}smooth\_\allowbreak{}is\allowbreak{}Proj\allowbreak{}By} for every commutative
base ring.  The other polarized, vanishing, rank, multiplication and
square-law suppliers are not conclusions of this chapter.}

The scheme and its bundles remain fixed throughout each construction.
In particular, the section module on an affine chart is not replaced
by an unrelated projective module.  The key geometric argument has two
steps: first clear denominators on finitely many affine charts; then
multiply by a further common power to make those extensions agree on
the overlaps.  Agreement on the nonzero locus alone would not suffice
for gluing over a ring with zero divisors.

\paragraph{Notation and the original coordinate convention.}
For an open $V\subseteq Y$ write $\mathcal O(V)=\Gamma(V,\mathcal O_Y)$.
All restrictions in a displayed identity are understood to the indicated
common open.  A line bundle $M$ is described by opens $U_p$ and units
$a_{pq}$, with coordinate convention $z_p=a_{pq}z_q$.
Write $V_l$ and $b_{lk}$ for the corresponding data of a second
line bundle $H$.  For every $t\in\mathbb Z$, put
\[
 M(t)=M\otimes H^{\otimes t},\qquad
 z_{pl}=a_{pq}b_{lk}^{t}z_{qk}.                 \tag{ASL.1}
\]
Negative tensor powers mean powers of the dual bundle.  In the source
this is literally \texttt{tw M H t}: the transition of
$\mathrm{LB.div}(M,\mathrm{pw}(H,-t))$ is the unit in (ASL.1).
We write $\operatorname{Sec}(M(t),V)$ for its actual compatible
coordinate families, the source's \texttt{Sec\allowbreak{}On}.  Its $R$-module
structure for $f:Y\to\operatorname{Spec}R$ is restriction through
the actual scalar map
\[
 c_f(V):R\longrightarrow\mathcal O(V).
                                                        \tag{ASL.2}
\]
Thus \texttt{Sec\allowbreak{}On\allowbreak{}R} changes the specified scalar action, not the sections.
Multiplication, restriction, pullback and gluing below are operations on
these families.  Ordinary facts about the structure sheaf of an affine
scheme, finite subcovers, and module localization will be used with
their concrete content spelled out where it matters.  No Noetherian,
reducedness or flat-base-change hypothesis on the original $R$ is imposed.

\subsection{The local calculations and the two common exponents}

% BEGIN ASL1
\paragraph{ASL1: fractions on an affine open and power-killing on a compact open.}
\label{claim:ASL1}
Let $W$ be affine, $u\in\mathcal O(W)$, and let $D_W(u)$ be its
principal open.  Every $z\in\mathcal O(D_W(u))$ has the form
$v/u^n$ with $v\in\mathcal O(W)$ and $n\ge0$.  If $v$ restricts
to zero on $D_W(u)$, some power $u^nv$ is zero on $W$.
More generally the latter conclusion holds when $W$ is any
quasi-compact open of any scheme.

\begin{proof}
For $W=\operatorname{Spec}B$, the definition of the structure sheaf
on the basis of principal opens gives
$\mathcal O(D_W(u))=B[u^{-1}]$, with restriction $v\mapsto v/1$.
Every localization element is a fraction with a power of $u$ as
denominator.  The equality $v/1=0$ means exactly that
$u^nv=0$ for some $n$.  Neither assertion cancels $u$.

For a general quasi-compact $W$, choose finitely many affine opens
$W_1,\ldots,W_r$ covering it.  Restriction of $D_W(u)$ to $W_j$
is $D_{W_j}(u|_{W_j})$.  The affine statement therefore supplies
an exponent $n_j$ killing $v|_{W_j}$.  Choose $N\ge n_j$ for all
$j$.  Then $u^Nv$ is zero on each $W_j$, and hence zero by the
separatedness part of the sheaf axiom.  For the empty open the
function ring is zero and exponent zero works.
\end{proof}
% END ASL1

% BEGIN ASL2
\paragraph{ASL2: nonzero loci, units, and actual twisted coordinates.}
\label{claim:ASL2}
Let $\sigma\in\operatorname{Sec}(H^{\otimes d},Y)$, where
$d\in\mathbb Z$, and write $u_l$ for its coordinates.  Then
$u_l=b_{lk}^{d}u_k$ on overlaps.  Its nonzero locus $D(\sigma)$
satisfies, for every open $W\subseteq V_l$,
\[
 D(\sigma)\cap W=D_W(u_l|_W).                       \tag{ASL.3}
\]
Consequently $u_l$ is a unit on every open contained in both
$V_l$ and $D(\sigma)$.  Multiplication by $\sigma^n$ sends a
section of $M(t)$ to a section of $M(t+nd)$, with coordinates
$u_l^nz_{pl}$.  On $D(\sigma)$ this operation is invertible.
These maps commute with every restriction.

\begin{proof}
Multiplication by the transition unit $b_{lk}^d$ does not change
whether the germ belongs to the maximal ideal of the local ring.
Thus the coordinatewise principal opens agree on overlaps and glue
to exactly (ASL.3).  A regular function whose germ is a unit at every
point has a regular inverse: its local inverses agree wherever both
are defined and glue by the structure sheaf.  This proves the unit
assertion without an integrality assumption.

Using (ASL.1), on a common open one has
\[
 u_l^nz_{pl}
 =a_{pq}b_{lk}^{t+nd}u_k^nz_{qk}.                 \tag{ASL.4}
\]
This is the precise transition law for the new twist.  On
$D(\sigma)$ the inverse uses $u_l^{-n}$ and satisfies the inverse
transition identity.  Restricting functions preserves products,
units and their integer powers, so the identities are compatible
with arbitrary restrictions.  The same reasoning shows that
sections in the refined charts of $M(0)$ are naturally the original
sections of $M$: their coordinates are the restrictions of $z_p$,
and the redundant $H$-index glues away.  This identification is also
restriction-natural.
\end{proof}
% END ASL2

% BEGIN ASL3
\paragraph{ASL3: finite simultaneous trivializing charts and the shift identity.}
\label{claim:ASL3}
If $Y$ is quasi-compact and $M,H$ are line bundles, there is a
finite affine cover $W_\alpha$ with
$W_\alpha\subseteq U_{p_\alpha}\cap V_{l_\alpha}$.
If $Y$ is quasi-separated, all $W_\alpha\cap W_\beta$ are
quasi-compact.  On an overlap write $a=a_{p_\alpha p_\beta}$,
$b=b_{l_\alpha l_\beta}$, and $u_\alpha=b^du_\beta$.
For all $t,d\in\mathbb Z$ and $n\ge0$,
\[
 (ab^{t+nd})u_\beta^n=u_\alpha^n(ab^t).          \tag{ASL.5}
\]

\begin{proof}
At every point choose one chart for each bundle.  Affine opens
form a basis of the scheme topology, so an affine neighborhood
lies in their intersection.  Quasi-compactness extracts a finite
subcover.  Quasi-separatedness says that the intersection of
two quasi-compact opens is quasi-compact; apply it to the affine
members.  Finally $u_\alpha^n=b^{nd}u_\beta^n$, and all factors
commute.  The laws of integer powers of a unit give (ASL.5), also
when $d$ or $t$ is negative.  This identity is what changes the
transition functions when the exponent used to clear denominators
is increased.
\end{proof}
% END ASL3

% BEGIN ASL4
\paragraph{ASL4: a local power kills a section in every original chart.}
\label{claim:ASL4}
Let $a\in\operatorname{Sec}(M(t),Y)$ restrict to zero on
$D(\sigma)$, with $\sigma\in\operatorname{Sec}(H^d,Y)$ and
$t,d\in\mathbb Z$.  For every $x\in Y$ there exist an open
neighborhood $W$ and $n\ge0$ such that, for \emph{every} original
pair of charts $(p,l)$ and every
$O\subseteq W\cap U_p\cap V_l$,
\[
                    (u_l|_O)^n a_{pl}|_O=0.          \tag{ASL.6}
\]
This is the full neighborhood statement
\texttt{exists\_\allowbreak{}nhds\_\allowbreak{}pow\_\allowbreak{}mul\_\allowbreak{}tw\allowbreak{}Coord\_\allowbreak{}eq\_\allowbreak{}zero}.

\begin{proof}
Choose an affine neighborhood $W\subseteq U_q\cap V_k$ of $x$.
The function $a_{qk}|_W$ is zero on
$W\cap D(\sigma)=D_W(u_k|_W)$, by ASL2.  ASL1 gives
$(u_k|_W)^n a_{qk}|_W=0$ for some $n$.
Restrict to the open $O$ in the statement.  The transition law gives there
\[
 (u_l|_O)^n a_{pl}|_O
   =a_{pq}|_O\, b_{lk}|_O^{\,t+nd}
                           (u_k|_O)^n a_{qk}|_O=0.
\]
The exponent is independent of $(p,l)$ and $O$ because it was
chosen on the one neighborhood $W$.  This establishes the full
order of quantifiers, rather than a separate exponent for every
subchart.
\end{proof}
% END ASL4

% BEGIN ASL5
\paragraph{ASL5: one global power kills the actual twisted section.}
\label{claim:ASL5}
If $Y$ is quasi-compact, the hypotheses of ASL4 imply that some
$k\ge0$ satisfies
\[
 \forall p,l\qquad u_l^k a_{pl}=0
                 \quad\hbox{on }U_p\cap V_l.         \tag{ASL.7}
\]
Equivalently the actual section $\sigma^k a$ of $M(t+kd)$ is
zero.  No quasi-separatedness is needed.  This proves the full
original \texttt{exists\_\allowbreak{}pow\_\allowbreak{}mul\_\allowbreak{}tw\allowbreak{}Coord\_\allowbreak{}eq\_\allowbreak{}zero}.

\begin{proof}
Apply ASL4 at every point and choose a finite subcover
$W_1,\ldots,W_r$ of the resulting neighborhoods.  Let $k$ be
the maximum of their exponents.  On
$W_j\cap U_p\cap V_l$, multiply (ASL.6) by the remaining
power of $u_l$.  The result is zero for all $j$.
These opens cover $U_p\cap V_l$, so the sheaf axiom gives
(ASL.7) on the full original overlap.  The same single $k$
works for every $(p,l)$.  If the scheme is empty, take $k=0$.
\end{proof}
% END ASL5

% BEGIN ASL6
\paragraph{ASL6: local extensions with one denominator exponent.}
\label{claim:ASL6}
Let $Y$ be quasi-compact, put $V=D(\sigma)$, and let
$y\in\operatorname{Sec}(M(t),V)$.  On a finite simultaneous
affine cover $W_\alpha$ as in ASL3, there are functions
$z_\alpha\in\mathcal O(W_\alpha)$ and one $N\ge0$ such that
\[
 z_\alpha|_{V\cap W_\alpha}
      =u_\alpha^N y_\alpha.                         \tag{ASL.8}
\]
On $W_\alpha\cap W_\beta$, set
\[
 \tau_{\alpha\beta}=a_{p_\alpha p_\beta}
                  b_{l_\alpha l_\beta}^{\,t+Nd},
 \qquad
 \delta_{\alpha\beta}=z_\alpha-\tau_{\alpha\beta}z_\beta.
                                                        \tag{ASL.9}
\]
Then $\delta_{\alpha\beta}$ vanishes on
$V\cap W_\alpha\cap W_\beta$.

\begin{proof}
On $W_\alpha$, the open $V\cap W_\alpha$ is principal with
function $u_\alpha$.  ASL1 writes the regular function
$y_\alpha$ there as $v_\alpha/u_\alpha^{n_\alpha}$.
Take $N\ge n_\alpha$ for every member of the finite cover and
put $z_\alpha=u_\alpha^{N-n_\alpha}v_\alpha$.
This proves (ASL.8), with its equality on every smaller open
obtained by restriction.

On the part of an overlap contained in $V$, the transition law
for $y$ is $y_\alpha=a_{p_\alpha p_\beta}
b_{l_\alpha l_\beta}^{t}y_\beta$.  Equation (ASL.5) with
$n=N$ therefore gives
$u_\alpha^Ny_\alpha=\tau_{\alpha\beta}
u_\beta^Ny_\beta$.  Substitute (ASL.8) to obtain the required
vanishing.  This argument does not assert that the discrepancy
already vanishes on the whole overlap.
\end{proof}
% END ASL6

% BEGIN ASL7
\paragraph{ASL7: extension after one power, including all transition units.}
\label{claim:ASL7}
Let $Y$ be quasi-compact and quasi-separated, $H,M$ line bundles,
$d,t\in\mathbb Z$, and $\sigma\in\operatorname{Sec}(H^d,Y)$.
For every $y\in\operatorname{Sec}(M(t),D(\sigma))$ there are
$k\ge0$ and an actual global section
$a\in\operatorname{Sec}(M(t+kd),Y)$ such that
\[
                  a|_{D(\sigma)}=\sigma^k y.          \tag{ASL.10}
\]
In each original chart this is exactly
$a_{pl}|_{D(\sigma)\cap U_p\cap V_l}=u_l^ky_{pl}$.
This proves the full original \texttt{exists\_\allowbreak{}tw\_\allowbreak{}extend\_\allowbreak{}pow\_\allowbreak{}mul}.

\begin{proof}
Use ASL6.  Each $W_\alpha\cap W_\beta$ is quasi-compact by
ASL3.  On it, (ASL.3) identifies its intersection with
$D(\sigma)$ as the principal open of $u_\alpha$.
ASL1 applied to the discrepancy gives an exponent
$k_{\alpha\beta}$ with
$u_\alpha^{k_{\alpha\beta}}\delta_{\alpha\beta}=0$.
There are only finitely many ordered pairs, so choose one
$K\ge k_{\alpha\beta}$ for all of them.

Put $w_\alpha=u_\alpha^Kz_\alpha$.  The transition unit for
the final twist is
\[
 \tau'_{\alpha\beta}
   =a_{p_\alpha p_\beta}
            b_{l_\alpha l_\beta}^{\,t+(N+K)d}.
\]
Using (ASL.5), now with the starting twist $t+Nd$, gives on
the entire overlap
\[
 \begin{split}
 w_\alpha-\tau'_{\alpha\beta}w_\beta
  &=u_\alpha^Kz_\alpha
             -u_\alpha^K\tau_{\alpha\beta}z_\beta\\
  &=u_\alpha^K\delta_{\alpha\beta}=0.
 \end{split}                                            \tag{ASL.11}
\]
Thus the local coordinates genuinely glue to a section $a$ of
$M(t+(N+K)d)$ over $Y$.  To describe the gluing without suppressing
the original charts, on $W_\alpha\cap U_p\cap V_l$ define its
$(p,l)$-coordinate to be the transition unit from
$(p_\alpha,l_\alpha)$ to $(p,l)$ times $w_\alpha$.
Equation (ASL.11) and the cocycle identities make these functions
agree for different $\alpha$.  They therefore glue on every
$U_p\cap V_l$, and their resulting coordinates obey (ASL.1).

On $V\cap W_\alpha$, (ASL.8) says
$w_\alpha=u_\alpha^{N+K}y_\alpha$.
Transporting this identity to any $(p,l)$ uses (ASL.5) once
more and gives exactly the coordinate equation in the statement.
The finite cover then proves it on all of
$V\cap U_p\cap V_l$.  Take $k=N+K$.
All exponents can depend on the given section; no uniformity in
$y$ or in the twist has been asserted.  The argument works for
negative $t,d$, nilpotent coordinate functions and the empty scheme.
\end{proof}
% END ASL7

\subsection{The graded module and its actual chart comparisons}

Fix a commutative ring $R$, a morphism
$f:Y\to\operatorname{Spec}R$, and sections
$s_0,\ldots,s_m\in\operatorname{Sec}(H,Y)$.
Set $D_i=D(s_i)$ and, for a finite index set $I$,
\[
 D_I=\bigcap_{i\in I}D_i,\quad D_\varnothing=Y,
 \qquad X_I=\prod_{i\in I}X_i\in R[X_0,\ldots,X_m].
                                                        \tag{ASL.12}
\]
The actual graded module on an open $V$ is
\[
 G(V)=\bigoplus_{t\in\mathbb Z}\operatorname{Sec}(M(t),V),
 \qquad G=G(Y),\qquad S=R[X_0,\ldots,X_m].          \tag{ASL.13}
\]
Its elements have finite support in $t$.  Denote the inclusion of
the $t$th summand by $[z]_t$.  The source's supplied witness
\texttt{hx} only certifies commutation of the multiplication maps
defining this very module; it changes none of the coordinate formulas.

% BEGIN ASL8
\paragraph{ASL8: the product section and the exact polynomial action.}
\label{claim:ASL8}
The operators $X_i[z]_t=[s_i z]_{t+1}$ commute and define an
$S$-module structure on $G(V)$.  All restriction maps are graded
$S$-linear maps.  For every $I$ there is an actual section
\[
 \sigma_I=\prod_{i\in I}s_i\in\operatorname{Sec}(H^{|I|},Y),
 \quad (\sigma_I)_l=\prod_{i\in I}(s_i)_l,
 \quad D(\sigma_I)=D_I.                              \tag{ASL.14}
\]
For every $k\ge0$, multiplication by $X_I^k$ on a homogeneous
summand has exactly the coordinate action of $\sigma_I^k$.
More generally, if $F\in S$ is homogeneous of degree $d$, then
$F[z]_t=[z']_{t+d}$, where
\[
 z'_{pl}=F\bigl(c_f(U_p\cap V_l),
                      (s_0)_l,\ldots,(s_m)_l\bigr)z_{pl}.
                                                        \tag{ASL.15}
\]
All functions in (ASL.15) are restricted to $V\cap U_p\cap V_l$.

\begin{proof}
The coordinate for $s_i(s_jz)$ is
$(s_i)_l(s_j)_lz_{pl}$, which is symmetric in $i,j$.
Coordinatewise equality proves commutation, and the universal property
of the polynomial ring therefore gives the stated action.
Multiplication by a section commutes with restriction by ASL2,
so the direct-sum restriction is graded and $S$-linear.

On overlaps, each factor in $(\sigma_I)_l$ changes by the same
$H$-transition unit.  Their product changes by its $|I|$th power,
which is the transition for $H^{|I|}$.  In a local ring a finite
product is a unit if and only if each factor is a unit.  ASL2
therefore proves the equality of nonzero loci in (ASL.14).
The empty product is the constant section $1$ of $H^0$, with
nonzero locus $Y$.

For a monomial, repeated multiplication gives (ASL.15) directly;
its degree is the total number of factors.  Constants act through
(ASL.2).  Summing the monomial formulas proves the assertion for
every homogeneous polynomial, including zero.  This also proves
the source's \texttt{exists\_\allowbreak{}smul\_\allowbreak{}of\_\allowbreak{}tw\allowbreak{}Coord} assertion with
the prescribed target degree, and specializes to the stated
$X_I^k$ formula.
\end{proof}
% END ASL8

% BEGIN ASL9
\paragraph{ASL9: the original very-ample charts give exactly the needed topology.}
\label{claim:ASL9}
Suppose \texttt{Very\allowbreak{}Ample\allowbreak{}By}$(f,H,(s_i))$ holds.  Thus the
$D_i$ cover $Y$, each $D_i$ is affine, and on $D_i$ there are
the given functions $r_{ij}$ with $s_j=r_{ij}s_i$ and with
$R[X_0,\ldots,X_m]\to\mathcal O(D_i)$, $X_j\mapsto r_{ij}$,
surjective.  Then $Y$ is quasi-compact and quasi-separated,
and every nonempty $D_I$ is affine.

\begin{proof}
The finite cover by affine opens proves quasi-compactness.
On $D_i$, the section $s_i$ is invertible in the line-bundle
sense of ASL2.  Hence
\[
 D_I=D_{D_i}\left(\prod_{j\in I\setminus\{i\}}r_{ij}\right)
                       \qquad(i\in I).                \tag{ASL.16}
\]
This proves affineness of every nonempty intersection, with its
actual inclusion in $D_i$.

For clarity the quasi-separatedness assertion follows already
from this affine cover and its affine pairwise intersections.
Given two quasi-compact opens, refine each by finitely many
principal affine opens lying in the $D_i$.  The intersection of
one member from $D_i$ with one from $D_j$ is a principal open
of the affine $D_i\cap D_j$, obtained by restricting their two
defining functions and inverting their product.  It is therefore
quasi-compact.  The intersection of the original two opens is a
finite union of such intersections.  This is precisely the needed
quasi-separatedness condition.  No separation over $R$, or
Noetherian property of $R$, was inserted.
\end{proof}
% END ASL9

% BEGIN ASL10
\paragraph{ASL10: full twisted-section localization at every finite chart set.}
\label{claim:ASL10}
Under ASL9, for every $I$, every $t\in\mathbb Z$, and every
$a\in\operatorname{Sec}(M(t),Y)$,
\[
 a|_{D_I}=0\quad\Longrightarrow\quad
       \exists k\ge0:\ X_I^k[a]_t=0.                 \tag{ASL.17}
\]
For every $y\in\operatorname{Sec}(M(t),D_I)$ there are
$k\ge0$ and $a\in\operatorname{Sec}(M(t+k|I|),Y)$ such that
\[
             [a|_{D_I}]_{t+k|I|}=X_I^k[y]_t.         \tag{ASL.18}
\]
These prove the whole original
\texttt{twisted\_\allowbreak{}sections\_\allowbreak{}localise}, with the source's given
commuting-operator witnesses and every finite $I$.

\begin{proof}
Apply ASL5 to the actual product section $\sigma_I$ in ASL8,
using quasi-compactness from ASL9.  The resulting coordinate
identity is exactly (ASL.17) by (ASL.15) and injectivity of each
summand into the direct sum.

Apply ASL7 to the same section, using both topological properties
from ASL9.  It gives $a|_{D_I}=\sigma_I^ky$ in precisely the
twist $t+k|I|$.  ASL8 translates this to (ASL.18).  In
particular the original existential target degree $t'$ may be
taken to be this explicit integer.  If $I=\varnothing$, take
$k=0$ and use the identity restriction $Y\to Y$ in both statements.
Thus neither the empty vertex nor negative twists are omitted.
\end{proof}
% END ASL10

% BEGIN ASL11
\paragraph{ASL11: a canonical isomorphism from each homogeneous localization.}
\label{claim:ASL11}
For $I$ finite and $t\in\mathbb Z$, let $(G[X_I^{-1}])_t$
be the degree-$t$ homogeneous localization.  An element can be
represented by
\[
       a/X_I^k,\qquad a\in G_{t+k|I|}.
\]
There is an actual $R$-linear isomorphism
\[
 \gamma_{I,t}:(G[X_I^{-1}])_t
                 \xrightarrow{\ \sim\ }
                       \operatorname{Sec}(M(t),D_I),
 \qquad \gamma_{I,t}(a/X_I^k)=\sigma_I^{-k}(a|_{D_I}).
                                                        \tag{ASL.19}
\]
It is characterized by
$\sigma_I^k\gamma_{I,t}(a/X_I^k)=a|_{D_I}$.

\begin{proof}
Multiplication by every $s_i$ with $i\in I$ is invertible on
the graded section module $G(D_I)$, by ASL2.  Their product
and all its powers are therefore invertible.  The graded
restriction $G\to G(D_I)$ consequently extends uniquely to
localization; its degree-$t$ formula is (ASL.19).

More explicitly, equality of fractions $a/X_I^k=a'/X_I^{k'}$
means that a further power of $X_I$ kills
$X_I^{k'}a-X_I^ka'$.  Restricting to $D_I$, all these powers
are invertible operations on sections, so division gives the same
value in (ASL.19).  This also checks the formula directly in the
presence of zero divisors.  Addition uses a common denominator;
scalar multiplication uses (ASL.2).

If the value of $a/X_I^k$ is zero, then $a|_{D_I}=0$.
ASL10 kills the homogeneous numerator by a power of $X_I$,
so its fraction is zero.  Conversely ASL10 applied to a given
$y$ supplies $a$ and $k$ with $a|_{D_I}=\sigma_I^ky$;
the fraction $a/X_I^k$ maps to $y$.  Thus the displayed map
is injective and surjective.  For $I=\varnothing$ this is the
identity on the actual $t$th global-section summand.  These
are exactly the source's \texttt{hom\allowbreak{}Loc} modules and their
fraction identities; no abstract identification of equal ranks is used.
\end{proof}
% END ASL11

% BEGIN ASL12
\paragraph{ASL12: the full section-cube isomorphism, with its restriction maps.}
\label{claim:ASL12}
For every integer $t$, the maps in (ASL.19) form an isomorphism
of the actual cubes
\[
 \gamma_t:\operatorname{grCube}(G,t)
              \xrightarrow{\ \sim\ }
                  \operatorname{secCube}(f,M(t),(D_i)).
                                                        \tag{ASL.20}
\]
Every vertex, including $\varnothing$, is bijective.  This is
the full original
\texttt{exists\_\allowbreak{}gr\allowbreak{}Cube\_\allowbreak{}twisted\allowbreak{}Sections\_\allowbreak{}hom\_\allowbreak{}sec\allowbreak{}Cube\_\allowbreak{}bijective}.
Its coordinate formula commutes with pullback along every morphism
of the original schemes and bundles, and hence along every
specified base change.

\begin{proof}
Let $I\subseteq J$ and put $c=X_{J\setminus I}$.
The localization restriction sends
$a/X_I^k$ to $c^ka/X_J^k$, since $X_J=cX_I$.
Using the actual product sections, its image under $\gamma_{J,t}$
is
\[
 \sigma_J^{-k}\sigma_{J\setminus I}^{k}a|_{D_J}
       =\bigl(\sigma_I^{-k}a|_{D_I}\bigr)|_{D_J}.
                                                        \tag{ASL.21}
\]
All cancelled factors are units on $D_J$, so this identity
does not require them to be non-zero-divisors on $Y$.
It proves commutation with each original restriction map.
Identity and composition follow either from these formulas or
from the corresponding restriction identities for sections.
The vertexwise bijections of ASL11 therefore give the cube
isomorphism, with the precise maps required by the source.

Pullback of a section restricts its coordinate functions by the
specified scheme morphism.  It preserves products, transition
units and their inverses, and sends $D(s_i)$ to the nonzero
locus of the pulled section.  Applying pullback to (ASL.19)
therefore gives the same formula on the pulled data.  Equivalently,
the natural scalar map from the old graded section module to
the new one makes the localized comparison square commute.
This assertion does \emph{not} say that global sections commute
with an arbitrary nonflat base change.  The affine, nonempty
vertices do commute with it, as proved in ASL17 below; the
empty vertex requires no such assertion for the unaugmented
\v Cech complex.
\end{proof}
% END ASL12

% BEGIN ASL13
\paragraph{ASL13: division by the original section and semilinear dehomogenization.}
\label{claim:ASL13}
Let $V\subseteq D(s_i)$ and suppose functions
$r_j\in\mathcal O(V)$ satisfy the literal coordinate identities
$s_j=r_js_i$.  Then $r_i=1$.  Division by $s_i^t$ gives an
injective $\mathcal O(V)$-linear map
\[
 \operatorname{div}_{i,t}:\operatorname{Sec}(M(t),V)
                         \longrightarrow\operatorname{Sec}(M(0),V),
 \qquad z_{pl}\longmapsto(s_i)_l^{-t}z_{pl}.
                                                        \tag{ASL.22}
\]
In fact it is an isomorphism, with inverse multiplication by
$s_i^t$ on these unit coordinates.  The finite-support sum of
these maps defines a semilinear map
\[
 \varepsilon_V:G(V)\longrightarrow\operatorname{Sec}(M(0),V),
 \qquad \varepsilon_V(Fz)=\pi(F)\varepsilon_V(z),
 \quad \pi=\operatorname{eval}(c_f(V),(r_j)).        \tag{ASL.23}
\]
The degree-zero map is the identity.

\begin{proof}
On a common chart $(s_i)_l=b_{lk}(s_i)_k$, and
$z_{pl}=a_{pq}b_{lk}^{t}z_{qk}$.  Thus
\[
 (s_i)_l^{-t}z_{pl}
       =a_{pq}(s_i)_k^{-t}z_{qk},
\]
the transition law for $M(0)$.  All inverse integer powers
are defined because $V\subseteq D(s_i)$; they commute with
restriction.  The inverse formula, injectivity and the
degree-zero identity follow on each chart and hence globally.
The identity $s_i=r_is_i$ implies $r_i=1$ by cancellation
of the unit coordinate, followed by sheaf equality.

For each homogeneous $z\in G(V)_t$,
\[
 \operatorname{div}_{i,t+1}(s_jz)
               =r_j\operatorname{div}_{i,t}(z).
\]
This proves (ASL.23) for $F=X_j$.  It holds for constants by
the specified scalar map (ASL.2), then for monomials by iteration,
and for polynomials by addition.  Every element of $G(V)$ is a
finite sum of homogeneous elements, so this proves the full
semilinearity statement, not only its homogeneous restriction.
\end{proof}
% END ASL13

% BEGIN ASL14
\paragraph{ASL14: full dehomogenization on the actual very-ample chart.}
\label{claim:ASL14}
Let $B_i=\mathcal O(D_i)$ and
$P_i=\operatorname{Sec}(M(0),D_i)$.  For \emph{every} family
$(r_{ij})_j$ in $B_i$ satisfying $s_j=r_{ij}s_i$ in each original
$H$-chart over $D_i$, the actual map
\[
 \varepsilon_i:G\longrightarrow P_i,
 \qquad \varepsilon_i\left(\sum_t[a_t]_t\right)
                     =\sum_t s_i^{-t}a_t|_{D_i}
                                                        \tag{ASL.24}
\]
is semilinear for the original surjection
$\pi_i:S\to B_i$, $X_j\mapsto r_{ij}$, and has all the
source outputs:
\[
 \begin{gathered}
 r_{ii}=1,\qquad \varepsilon_i\text{ is surjective},\\
 a\in G_t,\ \varepsilon_i(a)=0
       \ \Longrightarrow\ \exists k\ge0:\ X_i^ka=0,\\
 \varepsilon_i([a]_0)=a|_{D_i}.
 \end{gathered}                                         \tag{ASL.25}
\]
Thus \texttt{exists\_\allowbreak{}dehom\_\allowbreak{}chart} holds at its full original
scope.  Moreover the actual induced quotient is
$G/(X_i-1)G\simeq P_i$.

\begin{proof}
The coordinate of $s_i$ is a unit on every $D_i\cap V_l$ by
ASL2.  Hence the ratio identities force $r_{ij}$ there to be
the coordinate of $s_j$ divided by that of $s_i$.  Every such
family therefore agrees on these opens with the family supplied
by \texttt{Very\allowbreak{}Ample\allowbreak{}By}.  These opens cover $D_i$,
so sheaf equality proves agreement globally.  In particular the
evaluation $\pi_i$ is the same surjection for every allowed family;
no choice of that family is an additional hypothesis.

Compose graded restriction $G\to G(D_i)$ with ASL13.
This gives (ASL.24), its semilinearity, $r_{ii}=1$, and the
last identity in (ASL.25).
For $y\in P_i$, apply ASL10 at degree zero and $I=\{i\}$.
It gives $a\in G_k$ with $a|_{D_i}=s_i^ky$; hence
$\varepsilon_i(a)=y$.  This proves surjectivity.
For homogeneous $a\in G_t$ in the kernel, injectivity of
division in (ASL.22) gives $a|_{D_i}=0$.  The first part of
ASL10 supplies the stated power of $X_i$.

Finally these are exactly the hypotheses of the already proved
FCD13 quotient argument: $X_i$ acts as $1$ in $P_i$, the map
is surjective, and its homogeneous kernel is killed by elementwise
powers of $X_i$.  Consequently its full kernel is $(X_i-1)G$.
One can see the last step directly by representing any finite sum
modulo $(X_i-1)G$ in one sufficiently high homogeneous degree;
if its image is zero, a power of $X_i$ kills that representative,
and $X_i$ is $1$ in the quotient.  Thus the quotient isomorphism
uses the same semilinear map, rather than a newly chosen module
isomorphism.
\end{proof}
% END ASL14

\subsection{The actual affine section modules}

% BEGIN ASL15
\paragraph{ASL15: ordinary line-bundle sections localize on every qcqs open.}
\label{claim:ASL15}
Let $M$ be a line bundle on any scheme $Y$, let $V$ be a
quasi-compact quasi-separated open, let $f\in\mathcal O(V)$,
and put $W=D_V(f)$.  Then
\[
 \begin{aligned}
 x|_W=0&\ \Longrightarrow\ \exists n\ge0:\ f^nx=0
                    &&(x\in\operatorname{Sec}(M,V)),\\
 \forall y\in\operatorname{Sec}(M,W)&\quad
 \exists n\ge0,\ x\in\operatorname{Sec}(M,V):\quad
                    x|_W=(f|_W)^ny.
 \end{aligned}                                            \tag{ASL.26}
\]
In particular the actual restriction map induces an isomorphism
\[
 \operatorname{Sec}(M,V)[f^{-1}]
                      \simeq\operatorname{Sec}(M,W),
 \qquad x/f^n\longmapsto(f|_W)^{-n}x|_W.            \tag{ASL.27}
\]
This proves the whole original \texttt{sec\allowbreak{}On\_\allowbreak{}basic\allowbreak{}Open\_\allowbreak{}localization}.

\begin{proof}
Work on the open subscheme $V$.  Take $H$ to be its trivial
line bundle, take $d=t=0$, and regard $f$ as a section of
$H^0$.  The actual sections of $M\otimes H^0$ are identified
with those of $M$ by the restriction-natural coordinate map
of ASL2.  ASL5 gives the first assertion of (ASL.26); ASL7
gives the second, with no change of the bundle because $d=0$.
This use of the trivial bundle changes no function or restriction
on the original open $V$.

The displayed fraction formula is well-defined because $f|_W$
is a unit.  The first assertion of (ASL.26) says that a fraction
in its kernel has numerator killed by a further power of $f$,
and is therefore zero.  The second assertion produces a
preimage for every $y$.  Applying the first assertion to a
difference also proves the exact equality criterion for two
restricted sections.  These are precisely the three localization
conditions used by the source's affine chart map.
\end{proof}
% END ASL15

% BEGIN ASL16
\paragraph{ASL16: finite presentation and projectivity of the actual affine section module.}
\label{claim:ASL16}
For every affine open $V$ of any scheme $Y$ and every actual
line bundle $M$, the module
\[
                    P=\operatorname{Sec}(M,V)
              \quad\hbox{over}\quad B=\mathcal O(V)
\]
is finitely presented and projective.  This proves the full
original \texttt{sec\allowbreak{}On\_\allowbreak{}finite\allowbreak{}Presentation\_\allowbreak{}projective} over
arbitrary affine coordinate rings.

\begin{proof}
Identify $V$ with $\operatorname{Spec}B$ by its canonical affine
identification.  Principal opens form a basis.  Refining the
given bundle charts and using quasi-compactness gives a finite
cover $D(f_i)$ on which $M$ has an actual generator $e_i$.
The dual bundle $M^\vee$ has the inverse transition units;
let $e_i^\vee$ be the dual local generator.  Coordinatewise
evaluation gives $M^\vee\otimes M\to\mathcal O_V$ and identifies
the endomorphism given by $u\otimes v$ with multiplication by
the function $v(u)$.

Apply ASL15 to extend $e_i$ and $e_i^\vee$ after powers of $f_i$.
Thus there are global sections $u_i\in P$ and
$v_i\in\operatorname{Sec}(M^\vee,V)$ and exponents $n_i,m_i$
with
\[
 u_i|_{D(f_i)}=f_i^{n_i}e_i,\qquad
 v_i|_{D(f_i)}=f_i^{m_i}e_i^\vee.
\]
Set $h_i=v_i(u_i)\in B$.  On $D(f_i)$ it equals
$f_i^{n_i+m_i}$.  ASL1 therefore gives $l_i\ge0$ such that
\[
               f_i^{l_i}h_i=f_i^{l_i+n_i+m_i}.       \tag{ASL.28}
\]
Put $\widetilde u_i=f_i^{l_i}u_i$,
$q_i=l_i+n_i+m_i$, and define the actual $B$-linear functional
$\lambda_i:P\to B$ by $\lambda_i(z)=v_i(z)$.
On every bundle chart, evaluation in rank one and (ASL.28) give
\[
                 \widetilde u_i\lambda_i(z)=f_i^{q_i}z
                       \qquad(z\in P).               \tag{ASL.29}
\]
This equality holds globally.  The extra factor $f_i^{l_i}$
was included precisely because equality on $D(f_i)$ alone
need not imply equality on $V$.

The functions $f_i^{q_i}$ generate the unit ideal.  If all
$q_i>0$, a prime containing them would contain every $f_i$,
contrary to the covering property.  If any $q_i=0$, one
generator is already $1$.  Hence choose $b_i\in B$ with
$\sum_i b_if_i^{q_i}=1$.  Summing (ASL.29) yields the
explicit finite dual-basis identity
\[
                    z=\sum_i\widetilde u_i\,
                                      b_i\lambda_i(z).  \tag{ASL.30}
\]
Let $\rho:B^r\to P$ send its standard basis to the
$\widetilde u_i$, and let
$j:P\to B^r$ have coordinates $b_i\lambda_i$.
Equation (ASL.30) is $\rho j=1$.  Thus $P$ is a direct
summand of a finite free module, hence finite projective.
For finite presentation explicitly, $\ker\rho$ is the image
of the idempotent $1-j\rho$ on $B^r$ and is generated by
its finitely many columns.  This is a finite presentation
of $P$ over $B$.

When $V$ is empty, $B$ and $P$ are zero and the conclusion
holds directly.  Otherwise the argument uses neither
Noetherianity of $B$ nor a pre-existing equivalence between
line bundles and projective modules.  It constructs the
splitting for this very section module.
\end{proof}
% END ASL16

% BEGIN ASL17
\paragraph{ASL17: actual affine pullback and the nonempty section vertices.}
\label{claim:ASL17}
For every ring map $B\to B'$ and every line bundle $M$ on
$\operatorname{Spec}B$, the canonical section map is an isomorphism
\[
 B'\otimes_B\operatorname{Sec}(M,\operatorname{Spec}B)
       \xrightarrow{\ \sim\ }
       \operatorname{Sec}(M',\operatorname{Spec}B'),
 \qquad b'\otimes z\longmapsto b'\,\pi^*z.          \tag{ASL.31}
\]
It commutes with the actual restriction maps between affine
opens.  In the very-ample situation, this gives, for every
$R\to R'$ and every nonempty $I$,
\[
 R'\otimes_R\operatorname{Sec}(M(t),D_I)
       \simeq\operatorname{Sec}(M'(t),\pi^{-1}D_I).
                                                        \tag{ASL.32}
\]
These comparisons use the supplied pullback square, not merely
one preferred Cartesian model.

\begin{proof}
Use the finite principal trivializing cover $D(f_i)$ in ASL16.
ASL15 identifies the localization of the actual global-section
module with the free rank-one section module on $D(f_i)$.
The same result on the pulled bundle identifies the target of
(ASL.31), localized at the image of $f_i$, with its pulled
rank-one module.  On these localizations the map sends the
actual generator $e_i$ to its pullback.  It is therefore an
isomorphism.  The images of the $f_i$ still generate the unit
ideal, so their principal opens cover $\operatorname{Spec}B'$.

To justify this last local test algebraically, a module whose
localizations at finitely many unit-ideal generators are zero
is zero: each of its elements is killed by a power of every
generator, and these finitely many powers still generate the
unit ideal.  Apply this to the kernel and cokernel of (ASL.31),
using exactness of localization.  That proves the global
isomorphism without flatness of $B'$ over $B$.
The canonical formula on pure tensors shows directly that
it commutes with restriction and with all local transition units.

For (ASL.32), write $D_I=\operatorname{Spec}B$ using ASL9.
Its pullback is canonically
$\operatorname{Spec}(R'\otimes_RB)$; associativity identifies
the left side with
$(R'\otimes_RB)\otimes_B\operatorname{Sec}(M(t),D_I)$.
Apply (ASL.31).  Any specified pullback square has its unique
comparison isomorphism with this affine model, and the pure
tensor formula transports through that isomorphism.  Hence the
coordinate identity holds for the actual specified square.
In particular these are the section comparisons in GSL9.
No assertion of an isomorphism at the global empty vertex is
needed or made for an arbitrary coefficient change.
\end{proof}
% END ASL17

% BEGIN ASL18
\paragraph{ASL18: the exact scalar map on a smooth or flat affine chart.}
\label{claim:ASL18}
Let $R$ be any commutative ring, $f:Y\to\operatorname{Spec}R$,
and $V\subseteq Y$ affine.  Then
\[
 \begin{split}
 f\text{ smooth}&\ \Longrightarrow\ c_f(V):R\to\mathcal O(V)
                                    \text{ is smooth},\\
 f\text{ flat}&\ \Longrightarrow\ c_f(V):R\to\mathcal O(V)
                                    \text{ is flat}.
 \end{split}                                            \tag{ASL.33}
\]
This is the full original \texttt{cst\_\allowbreak{}smooth\_\allowbreak{}and\_\allowbreak{}flat}.

\begin{proof}
Restrict $f$ to the affine source $V$ and the affine target
$\operatorname{Spec}R$.  The affine-local definition of a
smooth scheme morphism says that the induced homomorphism
\[
 \Gamma(\operatorname{Spec}R,\mathcal O)
                              \longrightarrow\mathcal O(V)
\]
is smooth; the analogous affine-local definition of flatness
says that its target is a flat module over the source.
Equivalently one can restrict the local smooth presentations,
or the stalkwise flat module condition, to these affine charts.
These are the usual scheme-to-ring meanings of the two
properties, with no finite-type assumption added to flatness.

The scalar map actually used in the section module is the
composite of this homomorphism with the canonical isomorphism
$R\simeq\Gamma(\operatorname{Spec}R,\mathcal O)$.
Transporting along a ring isomorphism preserves smoothness
and flatness: the polynomial lifting and finite-presentation
conditions, respectively the exact tensor functor, are simply
identified under restriction of scalars.  This composite is
exactly (ASL.2), including the original target-global-sections
identification.  It proves (ASL.33) for that map.
\end{proof}
% END ASL18

\subsection{The exact join with the completed Noetherian model}

% BEGIN ASL19
\paragraph{ASL19: one Noetherian model for the actual section cube.}
\label{claim:ASL19}
Let $R$ be arbitrary, let $f:Y\to\operatorname{Spec}R$ be
smooth, and fix the original line bundles $H,M$ and the
presentation \texttt{Very\allowbreak{}Ample\allowbreak{}By}$(f,H,(s_i))$.
There exist a Noetherian subring $R_0\subseteq R$, a cube
$C_0$ over $R_0$, and an actual cube map
\[
 \Theta:R\otimes_{R_0}C_0
                       \longrightarrow\operatorname{secCube}(f,M,(D_i))
                                                        \tag{ASL.34}
\]
such that every nonempty vertex of $C_0$ is $R_0$-flat,
all cohomology modules of its unaugmented \v Cech complex
are finite over $R_0$, and $\Theta_I$ is bijective for every
nonempty $I$.  These objects are fixed before any later
coefficient ring is chosen.

\begin{proof}
Form the actual $G$ in (ASL.13).  For chart $i$, use
\[
 B_i=\mathcal O(D_i),\qquad
 P_i=\operatorname{Sec}(M(0),D_i),\qquad
 \pi_i(X_j)=r_{ij},\qquad \varepsilon_i\text{ of (ASL.24)}.
\]
The ring map $R\to B_i$ is the original scalar map and is
smooth by ASL18.  ASL16 proves that $P_i$ is finitely
presented and projective over this same $B_i$.
The surjectivity of $\pi_i$ is part of the given
\texttt{Very\allowbreak{}Ample\allowbreak{}By} record.  ASL14 gives
$\pi_i(X_i)=1$, surjectivity of $\varepsilon_i$, and the
required elementwise $X_i$-power annihilation on its
homogeneous kernel.  Its degree-zero restriction is the
actual restriction of sections.

Thus every hypothesis of the full original algebraic theorem
FCD18 has now been supplied for these specific objects.
Reuse that theorem to obtain one $R_0,C_0$ and
\[
          \theta:R\otimes_{R_0}C_0
                          \longrightarrow\operatorname{grCube}(G,0)
\]
with its three original outputs.  ASL12 gives the actual
cube isomorphism $\gamma_0$ to the sections of $M(0)$.
ASL2 identifies those sections with the sections of $M$,
compatibly with every restriction.  Define $\Theta$ to be
this specific composite of $\theta$, $\gamma_0$ and the
inverse of the $M\to M(0)$ identification.

Flatness and finite cohomology of $C_0$ are the outputs of
FCD18.  For nonempty $I$, every factor of $\Theta_I$ is
bijective.  Their compatibility with restrictions makes
(ASL.34) the actual cube map, rather than a list of unrelated
vertex isomorphisms.  No geometric section comparison or
affine projectivity premise remains unproved in this application
of FCD18.  In particular $R_0\to R$ need not be flat.
\end{proof}
% END ASL19

% BEGIN ASL20
\paragraph{ASL20: the full universally valid finite-projective section complex.}
\label{claim:ASL20}
Under the original hypotheses of ASL19, there are a
nonnegative cochain complex $K$ of $R$-modules and an actual
map
\[
                 \phi:K\longrightarrow\check C((D_i),M)
                                                        \tag{ASL.35}
\]
such that each $K^p$ is finite projective, $K^p=0$ for
$p\ge m+1$, and for \emph{every} ring map $R\to R'$,
$R'\otimes_R\phi$ is a quasi-isomorphism.  The target has
the original alternating restriction differential and is the
source's \texttt{cech\allowbreak{}LB}.  This proves all conclusions of
\texttt{exists\_\allowbreak{}finite\_\allowbreak{}projective\_\allowbreak{}complex\_\allowbreak{}cech\allowbreak{}LB}.

\begin{proof}
Apply the already proved GSL8 to the model and actual cube
map (ASL.34).  Its assumptions are exactly the three
outputs of ASL19.  The unaugmented \v Cech term in degree
$p$ uses the $(p+1)$-element subsets of
$\{0,\ldots,m\}$; hence it uses only nonempty vertices.
GSL8 gives (ASL.35), finite projectivity and the asserted
vanishing bound.  The morphism is the scalar extension of
the one Noetherian complex map, followed by the actual
section comparison already fixed in ASL19.

For any subsequent $R\to R'$, the associativity of scalar
extension identifies this with the same Noetherian model
tensored directly by $R'$ over $R_0$.  The universal
quasi-isomorphism assertion in GSL8 applies, without assuming
either ring map is flat.  ASL17 identifies the tensor-extended
section complex with the actual section complex on every
specified pulled-back scheme: finite direct sums and the
restriction differentials commute with the comparisons.
There is therefore no additional assertion about base change
of global sections hidden in (ASL.35).
\end{proof}
% END ASL20

% BEGIN ASL21
\paragraph{ASL21: the full original smooth-projective \texttt{Perf} producer.}
\label{claim:ASL21}
For every commutative ring $R$, every scheme $Y$ and map
$f:Y\to\operatorname{Spec}R$, if $f$ is smooth and satisfies
the original \texttt{Is\allowbreak{}Proj\allowbreak{}By}, then for every actual line
bundle $M$ the literal record $\mathrm{Perf}(f,M)$ exists.
This is the full original
\texttt{nonempty\_\allowbreak{}perf\_\allowbreak{}of\_\allowbreak{}smooth\_\allowbreak{}is\allowbreak{}Proj\allowbreak{}By}, with arbitrary
later coefficient maps and arbitrary specified pullback squares.

\begin{proof}
Unpack \texttt{Is\allowbreak{}Proj\allowbreak{}By} once to obtain $H,m,(s_i)$ and its
\texttt{Very\allowbreak{}Ample\allowbreak{}By} presentation.  ASL9 supplies the finite
cover with affine nonempty intersections.  ASL20 now supplies
the universally valid finite-projective model for its actual
$M$-section complex.  These are all the inputs to the already
proved GSL10, so its construction applies without a remaining
geometric producer.

To identify precisely the output being reused, GSL10 applies
the fixed finite-matrix construction to $K^0\to K^1\to K^2$.
It chooses finite sizes $a,b,c$ and matrices
\[
 D_0:R^a\to R^b,\qquad D_1:R^b\to R^c,
                         D_1D_0=0.                  \tag{ASL.36}
\]
It refines the cover to $W_{ip}=D_i\cap U_p$, where $U_p$
are the original $M$-charts.  Its fixed functions $u_{k,ip}$
are the coordinates on $W_{ip}$ of the $i$th components
of $\phi^0$ applied to the chosen degree-zero elements.
Thus $a,b,c,D_0,D_1,W,\mathrm{ch},u$ are chosen over the
original $R$ before any coefficient change.

Let $g:R\to R'$ be arbitrary, and let
$\pi:Y'\to Y$, $f':Y'\to\operatorname{Spec}R'$ be any
specified Cartesian square.  The GSL10 matrix comparison,
the universal quasi-isomorphism in ASL20, and the actual
section comparison of ASL17 give a bijection
\[
 E_g:\ker(D_0\otimes_RR')
                       \xrightarrow{\ \sim\ }\operatorname{Sec}(M',Y').
                                                        \tag{ASL.37}
\]
In the original pulled coordinates it satisfies
\[
 E_g(v)|_{\pi^{-1}W_{ip}}
       =\sum_{k=1}^{a}c_{f'}(\pi^{-1}W_{ip})(v_k)
                                          \,\pi^*u_{k,ip}.
                                                        \tag{ASL.38}
\]
Indeed each map used to construct (ASL.37) acts on the
original coordinates and the original restriction maps.
ASL12 and ASL17 verify their naturality, including transport
through the unique isomorphism from the supplied pullback
square to the standard affine tensor-product model.
Consequently (ASL.38) is the printed \texttt{Perf.\allowbreak sec} field,
not just an abstract isomorphism of modules.

If the pulled bundle satisfies the original all-cover
$\mathrm{FirstCohomologyVanishes}(M')$ predicate, GSL9
makes its section \v Cech complex exact in degree one.
ASL20 and the GSL10 finite-matrix comparison then give
\[
 \ker(D_1\otimes_RR')=\operatorname{im}(D_0\otimes_RR').
                                                        \tag{ASL.39}
\]
This is exactly \texttt{Perf.\allowbreak exact}, with its own stated
vanishing antecedent.  No such antecedent was needed to
construct the record.  The refinements cover $Y$ and lie
in their indicated original charts, so the remaining
\texttt{le}, \texttt{cover} and composition fields also hold.
All literal fields of \texttt{Perf} have thus been supplied
in the original quantifier order.
\end{proof}
% END ASL21

\paragraph{Completed geometric supplier and retained boundaries.}
ASL21 removes the smooth-projective perfect-section-complex
producer from the remaining inputs of the GSL section-law route.
The Noetherian-model algebra is reused from FCD18, and its
comparison is now attached to the original sections and bundles.
Neither a field nor a reduced base was substituted for the
original arbitrary ring.  The exponent choices are the actual
finite-cover choices, and the complex and its coordinate functions
are fixed before all later base changes.

The independent GSL inputs concerning polarized base change,
Serre or geometric $H^1$ vanishing, the multiplication map and
its invariants, geometric integrality, principal rank one,
torsion-section decomposition and finite growth, and the square
generators remain.  The original \texttt{Moduli} record and
the retained HUF/PUN foundations are unchanged.  Thus this
chapter supplies a full original \texttt{Perf} producer;
it does not by itself supply global \texttt{Sec\allowbreak{}Base\allowbreak{}Change\allowbreak{}St},
\texttt{Sec\allowbreak{}Base\allowbreak{}Change\allowbreak{}Any\allowbreak{}St}, \texttt{Sec\allowbreak{}Twist\allowbreak{}St}, or a whole
selected RH/GRH obligation.

\subsection{Exact source correspondence and inherited proof inputs}

All primary source references are at commit
\texttt{7cbf7a0c32fca1badf05931879a37c3f8a2b03b4} of the
accepted repository.  In the following list $C$ abbreviates
\texttt{src/\allowbreak{}Development/\allowbreak{}Abelian\allowbreak{}Variety\allowbreak{}Theory/\allowbreak{}Proofs/\allowbreak{}}
\texttt{Coherent\allowbreak{}Cohomology\allowbreak{}Proofs/\allowbreak{}}.  Full files and exact
range hashes are supplied in the component manifest.

\begin{itemize}
\item ASL1--ASL7:
$C$\texttt{Exists\allowbreak{}Nhds\allowbreak{}Pow\allowbreak{}Mul\allowbreak{}Tw\allowbreak{}Coord\allowbreak{}Eq\allowbreak{}Zero.lean},
lines 13--62, and
$C$\texttt{Exists\allowbreak{}Tw\allowbreak{}Extend\allowbreak{}Pow\allowbreak{}Mul.lean}, lines 15--171.
The two denominator exponents, all original chart indices
and the integer twists match these complete source statements.
\item ASL8--ASL10:
$C$\texttt{Exists\allowbreak{}Smul\allowbreak{}Of\allowbreak{}Tw\allowbreak{}Coord.lean}, lines 13--49;
$C$\texttt{Mem\allowbreak{}Nonzero\allowbreak{}Locus\allowbreak{}Iff\allowbreak{}Mem\allowbreak{}Basic\allowbreak{}Open\allowbreak{}Chart.lean},
lines 13--59;
$C$\texttt{Very\allowbreak{}Ample\allowbreak{}Compact\allowbreak{}Quasi\allowbreak{}Separated\allowbreak{}Affine\allowbreak{}Charts.lean},
lines 17--50; and
$C$\texttt{Twisted\allowbreak{}Sections\allowbreak{}Localise\allowbreak{}To\allowbreak{}Charts.lean}, lines 24--50.
\item ASL11--ASL12:
$C$\texttt{Exists\allowbreak{}Hom\allowbreak{}Loc\allowbreak{}Bijective\allowbreak{}Gr\allowbreak{}Sum.lean}, lines 45--162,
and $C$\texttt{Exists\allowbreak{}Twisted\allowbreak{}Sections\allowbreak{}Cube\allowbreak{}Bijective.lean},
lines 21--63.  The proof gives the explicit fraction map
whose characterization is used by those source constructors.
\item ASL13--ASL14:
$C$\texttt{Exists\allowbreak{}Dehomogenisation\allowbreak{}Twisted\allowbreak{}Chart.lean},
lines 22--189.  The final quotient-kernel step reuses FCD13.
\item ASL15--ASL18:
$C$\texttt{Sec\allowbreak{}On\allowbreak{}Basic\allowbreak{}Open\allowbreak{}Localization.lean}, lines 14--21;
$C$\texttt{Sec\allowbreak{}On\allowbreak{}Finite\allowbreak{}Presentation\allowbreak{}Projective.lean},
lines 16--65; and $C$\texttt{Cst\allowbreak{}Smooth\allowbreak{}And\allowbreak{}Flat.lean},
lines 12--20.  ASL16 replaces the source's local
finite-presentation/flatness criterion by a proved finite
dual-basis identity on the same module.
The canonical affine pullback map in ASL17 is the map in
\texttt{Proofs/\allowbreak{}Abelian\allowbreak{}Scheme\allowbreak{}Law\allowbreak{}Proofs/\allowbreak{}}
\texttt{Is\allowbreak{}Localized\allowbreak{}Module\allowbreak{}Res\allowbreak{}Basic\allowbreak{}Open\allowbreak{}Of\allowbreak{}Qcqs.lean},
lines 92--115 and 187--261; its actual restriction and
\v Cech comparisons appear there in lines 294--375.
\item ASL19--ASL20:
$C$\texttt{Exists\allowbreak{}Finite\allowbreak{}Projective\allowbreak{}Cech\allowbreak{}Complex\allowbreak{}Line\allowbreak{}Bundle.lean},
lines 19--43.  These use the already reviewed FCD18 and GSL8;
the actual-section base-change input agrees with GSL9.
\item ASL21:
\texttt{Proofs/\allowbreak{}Abelian\allowbreak{}Scheme\allowbreak{}Law\allowbreak{}Proofs/\allowbreak{}}
\texttt{Nonempty\allowbreak{}Perfect\allowbreak{}Sections\allowbreak{}Complex\allowbreak{}Of\allowbreak{}Smooth\allowbreak{}Proj.lean},
lines 15--22, and the literal \texttt{Perf} contract in
\texttt{Definitions/\allowbreak{}Abelian\allowbreak{}Scheme\allowbreak{}Laws/\allowbreak{}}
\texttt{First\allowbreak{}Cohomology\allowbreak{}Vanishes\allowbreak{}Spec\allowbreak{}Of\allowbreak{}Is\allowbreak{}Local\allowbreak{}Ring.lean},
lines 47--81.  The finite-matrix and record construction is
the unchanged reviewed GSL10.
\end{itemize}


\section{The connected rank-two core and descent of finite points}
\label{sec:rank-two-connectedness}

\paragraph{The arithmetic purpose.}
The rank-two propagation argument needs the actual core of its universal
representation space to be preconnected.  This is a statement about
the space already attached to the fixed frame; a connected auxiliary
space would not suffice.  We construct its points from strict extension
cocycles and join them by polynomial lines.  In doing so we prove the
upstream weighted-coordinate and continuity results needed to recover
an arbitrary point of that core.  These results use the universal
deformation property, with no finiteness conclusion from a later
arithmetic endpoint.

The point construction differs from ACP's finite-point shortcut.
Here a residual representation may have a nonzero lower off-diagonal
entry, and its target may be an arbitrary, possibly infinite-dimensional
or nonreduced, algebra over the residue field.  We first realize a
finite family of its cocycles over a formal power-series ring with the
original residue field.  Only the polynomial image of the entry algebra
is then evaluated in the prescribed target.  No evaluation of an
arbitrary formal power series at a nonnilpotent element is used.

The last part proves the original coefficient-descent theorem
\(\mathrm{FinPtDown}\), at every original rank.  It is the transport
between the extended coefficient frame returned by the actual two-run
argument and the original frame required by AFE20.  It does not supply
the rank-three propagation or automorphic hypotheses of that argument.

\paragraph{The fixed data and the original rings.}
Let \(O\) be the original complete discrete valuation ring, let
\(k=O/\mathfrak m_O\) be its finite residue field of characteristic
\(\ell\), and fix an original frame \(F\) with
\(F.\mathrm{Ok}(C_0)\).  Put \(G_0=\ker\delta\), where
\(\delta=F.\mathrm{del}\), write \(c=F.c\), and put
\[
 g^*=cg^{-1}c^{-1},\qquad \Psi=F.\Psi,
 \qquad w(2)=(1,-1),\quad w(3)=(1,0,-1).
                                                        \tag{RTC.1}
\]
The original hypotheses give \(c^2=1\), \(\delta(c)=-1\),
\(\Psi(c)=1\), odd \(\ell\), and continuous local maps
\(\operatorname{loc}_{0,v}:G_v\to G_0\).  In particular
\(\Psi(g^*)=\Psi(g)^{-1}\).  Both \(G_0\) and the local groups
are compact: \(G_0\) is a closed open subgroup of the original
profinite group.  The kernel of \(\Psi\) is open.

For \(n=2\) or \(3\), a supplied
\(S:\mathrm{Space}(O,F,n)\) consists of a complete Noetherian
local \(O\)-algebra \(R\), with its specified residue field \(k\),
and a universal normalized representation \((r,\psi)\).  Its
matrix and character entries are adically continuous, its residual
representation is literally diagonal with characters
\(\Psi^{w(n)_i}\), and it satisfies
\[
 r(g^*)=J_nr(g)^tJ_n,\qquad
 r(F.g_0)-r((F.g_0)^*)\text{ is diagonal},
                                                        \tag{RTC.2}
\]
and, for every local tuple \((s_0,\ldots,s_{n-1})\),
\[
 \prod_{i=0}^{n-1}
 \bigl(r(\operatorname{loc}_{0,v}s_i)-\psi_{v,i}(s_i)I_n\bigr)=0.
                                                        \tag{RTC.3}
\]
The order in this product is part of the condition.  Universality
holds for every complete Noetherian local \(O\)-algebra with the
same residue field and its actual \(O\)-structure.  We call such
an algebra a coefficient algebra; this is exactly \(\mathrm{IsCO}\).

The weight of \(r(g)_{ij}\) is \(w(n)_i-w(n)_j\).
With closure taken in the \(\mathfrak m_R\)-adic topology, define
\[
 \begin{aligned}
 B_0&=O[\text{weight-zero entry monomials},\ \psi_{v,i}(s)],
       &B&=\overline{B_0}\subset R,\\
 A&=B[r(g)_{ij}:g,i,j]\subset R,
       &\mathfrak m_B&=B\cap\mathfrak m_R,
       &x_0&=\mathfrak m_B A.
 \end{aligned}                                             \tag{RTC.4}
\]
Write \(u\) for the representation over \(A\).  The module
\(\operatorname{gr}_j A\) is the \(B\)-span of entry monomials
of total weight \(j\).  The ideals \(\mathfrak t_+\) and
\(\mathfrak t_-\) are generated by the entries below and above
the diagonal, respectively.  The actual rank-two core is
\[
 \mathcal C= S.\mathrm{BB}
 =\{\mathfrak p\in\operatorname{Spec}A:
       x_0\subset\mathfrak p,\quad
       \mathfrak t_+\not\subset\mathfrak p\}.
                                                        \tag{RTC.5}
\]
The notation \(R_j\) below means a homogeneous submodule of \(R\),
not a new deformation ring.  In general \(R\) itself is not the
direct sum of these submodules.

\subsection{Upstream weighted coordinates and continuity}

% BEGIN RTC1
\paragraph{RTC1. Locality and uniqueness for coefficient-algebra maps.}
\label{claim:RTC1}
Every \(O\)-algebra map between coefficient algebras preserves the
specified residue map to \(k\), is local, and is adically continuous.
Two maps out of \(R\) that agree on all entries of \(r\) and all
values of \(\psi\) are equal.

\emph{Proof.}
Let \(f:R\to C\) be such a map.  If \(x\in\mathfrak m_R\),
choose \(a\in O\) with \(f(x)-a\in\mathfrak m_C\), using the
specified surjectivity of the residue map from \(O\).  If \(a\)
were a unit of \(O\), then \(x-a\) would be a unit of \(R\),
whereas \(f(x-a)\in\mathfrak m_C\), a contradiction.  Thus
\(a\in\mathfrak m_O\), so \(f(x)\in\mathfrak m_C\).
For arbitrary \(y\), choose \(a\in O\) with
\(y-a\in\mathfrak m_R\).  The preceding observation shows that
both residues of \(y\) and \(f(y)\) are the residue of \(a\).
It also gives
\(f(\mathfrak m_R^N)\subset\mathfrak m_C^N\) for every \(N\),
which proves continuity.  Units map to units, so the map is local.

Mapping the universal normalized representation through any such
\(f\) preserves all identities, adic continuity and the exact
diagonal reduction.  If \(f\) and \(f'\) have the stated equalities,
they give the same normalized lift over \(C\), including the local
characters.  The uniqueness clause in the supplied universal property
therefore gives \(f=f'\).  No assertion about an arithmetic
specialization is needed. \hfill\(\square\)
% END RTC1

% BEGIN RTC2
\paragraph{RTC2. Finitely many actual entries give formal coordinates.}
\label{claim:RTC2}
There is a finite list \(z_1,\ldots,z_s\in\mathfrak m_R\), each
either an off-diagonal entry of \(r\), a diagonal entry minus a
chosen lift in \(O\) of its residue, or a character value minus
such a lift, for which evaluation is a surjection
\[
 \pi:P=O[[X_1,\ldots,X_s]]\longrightarrow R,
       \qquad X_i\longmapsto z_i.
                                                        \tag{RTC.6}
\]
Give each off-diagonal coordinate its entry weight and each other
coordinate weight zero.

\emph{Proof.}
Let \(e_a\) denote the entire family of centered entries just
described.  All lie in \(\mathfrak m_R\).  Consider the finite
dimensional \(k\)-space
\[
 V=\mathfrak m_R/(\mathfrak m_R^2+\mathfrak m_OR).
\]
The classes of the \(e_a\) span \(V\).  Indeed, otherwise there
is a nonzero \(k\)-linear functional \(\lambda:V\to k\)
vanishing on those classes.  It defines an \(O\)-algebra map to
\(k[\epsilon]/(\epsilon^2)\): for \(r-a\in\mathfrak m_R\),
\(a\in O\), put
\[
 f_\lambda(r)=\bar a+\epsilon\lambda([r-a]).
                                                        \tag{RTC.7}
\]
Different choices of \(a\) differ by an element of
\(\mathfrak m_O\), so the formula is independent of that choice.
Addition follows by choosing the sum of two lifts.  For multiplication,
write \(r=a+x\), \(r'=a'+x'\); the term \(xx'\) vanishes in
\(V\), and multiplication by \(a\) acts there by \(\bar a\).
This proves the ring and \(O\)-algebra identities in (RTC.7).

The maps \(f_\lambda\) and \(f_0\) agree on every centered entry,
hence on every matrix entry and character value.  Their target is a
coefficient algebra.  RTC1 makes them equal, which forces
\(\lambda=0\), a contradiction.  Since \(V\) is finite
dimensional, finitely many \(e_a\), now denoted \(z_i\), span it.
Nakayama's lemma applied to the finite module
\(\mathfrak m_R/(\mathfrak m_OR+(z_1,\ldots,z_s))\) gives
\[
 \mathfrak m_R=\mathfrak m_OR+(z_1,\ldots,z_s).
                                                        \tag{RTC.8}
\]

Evaluation of a power series at the \(z_i\) converges in the
complete ring \(R\), since its terms of variable degree at least
\(N\) lie in \(\mathfrak m_R^N\).  It respects products by
finite truncation and passage to the limit.  To prove surjectivity,
let \(\varpi\) generate \(\mathfrak m_O\).  Equation (RTC.8)
allows a successive approximation to any \(r\in R\): choose an
\(O\)-lift of its residue, then correct the error modulo
\(\mathfrak m_R^2\) by a homogeneous polynomial of degree one
in \(\varpi,z_1,\ldots,z_s\), and continue.  At step \(N\),
reduce the finitely many coefficients in an expression for the error
in \(\mathfrak m_R^N\) modulo \(\mathfrak m_R\), and lift them
to \(O\).  The corrections in \(P\) lie successively in
\((\varpi,X_1,\ldots,X_s)^N\).

This series of corrections converges in \(P\).  Explicitly, for
each fixed monomial its coefficients form a \(\varpi\)-adic
Cauchy sequence in \(O\); their limits define a formal series,
and any fixed finite total-order jet stabilizes to the required
precision.  Its image is \(r\), by separatedness of \(R\).
The same argument covers an empty list of coordinates. \hfill\(\square\)
% END RTC2

% BEGIN RTC3
\paragraph{RTC3. Universal diagonal twists make the kernel homogeneous.}
\label{claim:RTC3}
Write \(d_i\in\mathbb Z\) for the weights of the coordinates in
(RTC.6).  If \(P_j\) is the set of formal series supported on
monomials \(X^\alpha\) with \(\sum_i d_i\alpha_i=j\), then
\(\ker\pi\) is stable under every projection
\(f\mapsto f_j\) onto \(P_j\).

\emph{Proof.}
For a coefficient algebra \(C\), a map \(f:R\to C\), and a
unit \(t\in C^\times\), conjugate the mapped representation by
\[
 D_t=\operatorname{diag}(t^{w(n)_0},\ldots,t^{w(n)_{n-1}}).
\]
Its \((i,j)\)-entry becomes
\(t^{w(n)_i-w(n)_j} f(r(g)_{ij})\), while every local character
is unchanged.  Since \(w(n)_{n-1-i}=-w(n)_i\), one has
\(J_nD_t^{-1}=D_tJ_n\).  This verifies the polarization identity
for the conjugate.  The diagonal normalization is preserved, and
each factor in (RTC.3) is conjugated by the same \(D_t\).
Continuity and diagonal residual reduction follow from RTC1 and the
fact that only off-diagonal entries are rescaled.  Universality
therefore supplies a map \(\tau_{f,t}:R\to C\) with these
prescribed images.  In particular it sends each coordinate \(z_i\)
to \(t^{d_i}f(z_i)\).

Fix \(N\geq1\), put \(D=R/\mathfrak m_R^N\), and use the
coefficient algebra \(C=D[[Z]]\), the constant-coefficient map
\(f:R\to C\), and \(t=1+Z\).  Its residue field is still
\(k\).  If \(F\in\ker\pi\), applying \(\tau_{f,t}\) gives
\[
 0=\sum_{|\alpha|<N}
       \overline{a_\alpha z^\alpha}\,(1+Z)^{d\cdot\alpha},
       \qquad F=\sum_\alpha a_\alpha X^\alpha.
                                                        \tag{RTC.9}
\]
The sum is finite because the images of the \(z_i\) generate a
nilpotent ideal in \(D\).  Negative exponents are legitimate since
\(1+Z\) is a unit.  The substitution
\[
 D[Y,Y^{-1}]\longrightarrow D[[Z]],\qquad Y\longmapsto1+Z
                                                        \tag{RTC.10}
\]
is injective over any commutative ring \(D\).  Multiply a Laurent
polynomial by a power of \(Y\), apply the invertible polynomial
change of variable \(Y=1+Z\), and then use the coefficientwise
inclusion \(D[Z]\subset D[[Z]]\).  This proves injectivity even
when \(D\) has nilpotents or residue characteristic divides some
weight differences.

The coefficient of \(Y^j\) in (RTC.9) is
\(\pi(F_j)\) modulo \(\mathfrak m_R^N\), so it is zero.
This holds for every \(N\).  Separatedness of the coefficient
algebra \(R\) yields \(\pi(F_j)=0\), proving the claim.
\hfill\(\square\)
% END RTC3

% BEGIN RTC4
\paragraph{RTC4. Closed homogeneous parts and the actual entry weights.}
\label{claim:RTC4}
Put \(R_j=\pi(P_j)\).  Projection onto weight \(j\) induces a
continuous \(O\)-linear map \(p_j:R\to R_j\), preserving every
power of \(\mathfrak m_R\).  Each \(R_j\) is closed, finite
sums of distinct weights are direct, and
\(R_iR_j\subset R_{i+j}\).  Every universal matrix entry belongs
to its prescribed part, and every universal local character value
belongs to \(R_0\).

\emph{Proof.}
RTC3 makes \(p_j(\pi(F))=\pi(F_j)\) well-defined.  Let
\(\mathfrak n=(\varpi,X_1,\ldots,X_s)\subset P\).
Surjectivity and locality of \(\pi\) give
\(\pi(\mathfrak n^N)=\mathfrak m_R^N\).
Projection onto a set of monomials preserves \(\mathfrak n^N\):
the coefficient of \(X^\alpha\) in an element of this ideal is
divisible by
\(\varpi^{\max(0,N-|\alpha|)}\), and deleting other coefficients
preserves that condition.  For completeness, this coefficient
criterion follows by expanding products of \(N\) generators of
\(\mathfrak n\); conversely one partitions the indicated monomials
by a divisor of total order \(N\), giving a finite sum of such
products times formal series.  Hence \(p_j\) preserves
\(\mathfrak m_R^N\) and is continuous.

Since \(p_j^2=p_j\), its image is the kernel of the continuous
map \(1-p_j\), and is closed in the separated ring \(R\).
Applying \(p_j\) to a finite relation among different weights
proves directness.  The multiplication assertion follows from the
formal-series coefficient formula.  It includes \(1\in R_0\).

Let \(x\) be any entry of weight \(j\), or a local character
value with \(j=0\), and choose \(F\in P\) with \(\pi(F)=x\).
The maps \(\tau_{f,1+Z}\) in RTC3 send \(x\) to
\((1+Z)^j\bar x\).  Comparing coefficients in (RTC.10), now
with this expression on the right of (RTC.9), gives
\[
 \pi(F_j)=x\pmod{\mathfrak m_R^N},\qquad
 \pi(F_i)=0\pmod{\mathfrak m_R^N}\quad(i\ne j).
\]
Letting \(N\) vary proves \(x=\pi(F_j)\in R_j\).
Thus the grading refers to the actual universal entries, including
entries not chosen among the finitely many coordinates.
\hfill\(\square\)
% END RTC4

% BEGIN RTC5
\paragraph{RTC5. Finite monomial generators at every weight.}
\label{claim:RTC5}
For the finite integer weight list \(d_1,\ldots,d_s\), set
\[
 H_j=\{\alpha\in\mathbb N^s:d\cdot\alpha=j\}.
\]
The monoid \(H_0\) has a finite set \(H\) of nonzero generators.
For each \(j\), there is a finite set \(E_j\subset H_j\) such
that every \(\alpha\in H_j\) is \(\beta+\gamma\), with
\(\beta\in E_j\) and \(\gamma\in H_0\).  Consequently
\[
 P_j=\sum_{\beta\in E_j}P_0X^\beta.
                                                        \tag{RTC.11}
\]

\emph{Proof.}
We use the elementary finite-minimal-elements form of Dickson's
lemma.  Its sequential form says that every infinite sequence in
\(\mathbb N^s\) has two terms, in their original order, with
the first coordinatewise at most the second.  For \(s=1\), an
infinite sequence either has an infinitely repeated value or has an
increasing subsequence.  Induct on \(s\), first taking such a
nondecreasing subsequence in the first coordinate and then applying
the induction hypothesis to the remaining coordinates.  An infinite
set of pairwise incomparable minimal elements is therefore impossible.
Moreover, every member of a subset dominates a minimal member, by
descending the sum of its coordinates.

Take \(H\) to be the coordinatewise minimal elements of
\(H_0\setminus\{0\}\).  Given any nonzero \(\gamma\in H_0\),
choose \(h\in H\) with \(h\leq\gamma\).  Then
\(\gamma-h\in H_0\), with smaller total degree.  Repeating
proves that \(H\) generates \(H_0\).  Take \(E_j\) to be the
minimal elements of \(H_j\).  If \(\beta\leq\alpha\) are in
\(H_j\), their difference is in \(H_0\), proving the second
assertion.  Empty fibers cause no difficulty.

Finally, for each exponent \(\alpha\in H_j\), choose one such
\(\beta\).  Partition the monomials of a series in \(P_j\)
according to that choice, remove the factor \(X^\beta\), and
collect the remaining coefficients.  Each collected series has
weight zero.  There are only finitely many \(\beta\), and every
coefficient identity is finite, so this is an identity of formal
series.  It proves (RTC.11). \hfill\(\square\)
% END RTC5

% BEGIN RTC6
\paragraph{RTC6. The degree-zero ring is a complete coefficient algebra.}
\label{claim:RTC6}
The ring \(R_0\) is a complete Noetherian local \(O\)-algebra
with residue field \(k\), and its maximal ideal is
\(R_0\cap\mathfrak m_R\).  It is the image of the continuous
surjection
\[
 O[[U_h:h\in H]]\longrightarrow R_0,
       \qquad U_h\longmapsto\pi(X^h).
                                                        \tag{RTC.12}
\]

\emph{Proof.}
First map the source in (RTC.12) to \(P_0\) by
\(U_h\mapsto X^h\).  All \(h\) are nonzero, so the substitution
of formal series is defined.  It is surjective: for each
\(\alpha\in H_0\), choose a decomposition
\(\alpha=\sum_h n_h(\alpha)h\), and put the coefficient of
\(X^\alpha\) at the source monomial
\(\prod_hU_h^{n_h(\alpha)}\).  Different \(\alpha\) give
different chosen source monomials.  Since every \(h\) has positive
total degree, only finitely many source monomials contribute to a
fixed target coefficient.  Substitution recovers the prescribed
series.  Composing with \(\pi\) proves (RTC.12).

The source is a complete Noetherian local ring: Noetherianity is
the formal power-series theorem over the Noetherian DVR \(O\),
and completeness follows coefficientwise as in RTC2.  Its quotient
by the kernel in (RTC.12) is also complete and Noetherian.  Here
the quotient completeness is for its maximal-ideal topology: a
Cauchy sequence in the quotient can be lifted successively so that
the error at the \(N\)th stage lies in the \(N\)th power upstairs;
the lifted sequence converges.  Ideals are closed by Krull
intersection in the Noetherian local quotient, which also gives
separatedness.

All variables \(U_h\) map into \(\mathfrak m_R\), and the
residue of a constant is its original residue in \(k\).
Consequently the preimage of \(\mathfrak m_R\) is exactly the
maximal ideal of the source of (RTC.12): an element outside that
ideal is a unit, and one inside it has zero residue.  Thus the
quotient has residue field \(k\) and maximal ideal
\(R_0\cap\mathfrak m_R\).  This also proves locality of its
specified \(O\)-map and all the original \(\mathrm{IsCO}\)
clauses.  If \(H\) is empty, the source of (RTC.12) is simply
\(O\), with the same argument. \hfill\(\square\)
% END RTC6

% BEGIN RTC7
\paragraph{RTC7. The original weighted presentation and entry algebra.}
\label{claim:RTC7}
For either original rank \(n=2\) or \(n=3\), the presentation
(RTC.6) has every field of the original \(\mathrm{wpres}\), and
\[
 B=R_0,\qquad A=\bigoplus_{j\in\mathbb Z}R_j,
       \qquad \operatorname{gr}_jA=R_j\quad\text{inside }A.
                                                        \tag{RTC.13}
\]
In particular the original \(\mathrm{co\_\allowbreak B}\) and
\(\mathrm{gr\_\allowbreak A}\) conclusions hold: the grading is internal
and multiplicative, its degree-zero part is the unit \(B\)-module,
each graded piece is a finite \(B\)-module, \(A\) is a finite
type \(B\)-algebra and a Noetherian ring, and every entry has
its stated weight.

\emph{Proof.}
All generators of \(B_0\) have weight zero by RTC4, so the
closedness of \(R_0\) gives \(B\subset R_0\).  Conversely,
each selected weight-zero coordinate is in \(B_0\), while each
selected coordinate of nonzero weight is an actual off-diagonal
entry.  Therefore the image of every weight-zero coordinate monomial
is in \(B_0\): collect the weight-zero coordinates separately,
and recognize the remaining product as a weight-zero entry monomial.
Truncations of a series in \(P_0\) have images in \(B_0\) and
converge to its image.  Thus \(R_0\subset B\), and RTC6 gives
the complete coefficient-algebra structure on the original \(B\),
with exactly the original ideal \(\mathfrak m_B\).

By RTC4, the generators of the original \(A\) lie in the finite
sum algebra of the \(R_j\).  Conversely RTC5 expresses each
\(R_j\) as a finite \(B\)-span of selected coordinate monomials.
Those coordinates lie in \(A\), so all \(R_j\) lie in \(A\).
Directness and multiplicativity were proved in RTC4; this gives
the middle equality of (RTC.13).

An actual entry monomial of weight \(j\) lies in \(R_j\), so
\(\operatorname{gr}_jA\subset R_j\).  In the reverse inclusion,
factor each selected coordinate monomial of weight \(j\) into
its weight-zero coordinates and its actual off-diagonal entries.
The former contribute a coefficient in \(B\), and the latter
are an entry monomial of weight \(j\).  RTC5 then gives equality.
This proves all the identifications required by \(\mathrm{wpres}\),
not just an abstract grading with unidentified pieces.

The same finite-span argument proves finiteness of each graded
piece over \(B\).  In fact
\(A=B[z_i:d_i\ne0]\): every selected coordinate of weight zero
already belongs to \(B\), and the argument above generates every
homogeneous part from the selected coordinates.  Hilbert's basis
theorem and RTC6 give Noetherianity of \(A\).  All remaining
\(\mathrm{gr\_\allowbreak A}\) fields now follow from (RTC.13) and RTC4.
The proof has used the original universal property only.  No
\(\mathrm{FinA}\), \(\mathrm{FinPt}\), \(\mathrm{FinEnd}\),
modularity ideal or automorphic representation occurs as an input.
\hfill\(\square\)
% END RTC7

% BEGIN RTC8
\paragraph{RTC8. A cofinal filtration on every fixed weight.}
\label{claim:RTC8}
For every \(j\in\mathbb Z\) and \(N\geq0\), there is
\(M\geq0\) such that
\[
 R_j\cap\mathfrak m_R^M\ \subset\ \mathfrak m_B^N R_j.
                                                        \tag{RTC.14}
\]
More explicitly, let \(D\geq1\) bound the total degrees of
the zero-weight generators in \(H\), and let \(L\) bound the
degrees in \(E_j\).  If \(E_j\ne\varnothing\), one may use
\(M=DN+L\).

\emph{Proof.}
If \(E_j\) is empty, \(R_j=0\).  Otherwise take
\(x\in R_j\cap\mathfrak m_R^M\).  Lift it to an element of
\(\mathfrak n^M\subset P\), using the surjectivity in RTC2,
and project to weight \(j\).  RTC4 shows that this gives
\(F\in P_j\cap\mathfrak n^M\) with \(\pi(F)=x\).

For a nonzero coefficient \(a_\alpha\) of \(F\), set
\(a=\max(0,M-|\alpha|)\).  The coefficient criterion in
RTC4 gives \(a_\alpha=\varpi^a b_\alpha\).  Decompose
\[
 \alpha=\beta+h_1+\cdots+h_r,
       \qquad\beta\in E_j,\quad h_i\in H.
\]
Since \(|\beta|\leq L\) and \(|h_i|\leq D\),
\[
 M\leq a+|\alpha|\leq a+L+Dr\leq L+D(a+r).
\]
Thus \(a+r\geq N\).  From the factors consisting of \(a\)
copies of \(\varpi\) and \(X^{h_1},\ldots,X^{h_r}\),
choose \(N\).  Their product is one of finitely many monomials
of degree \(N\) in
\(\varpi,\{X^h:h\in H\}\), all of weight zero.  After its
removal the remaining term still has weight \(j\).

Partition the terms of \(F\) by this chosen product.  For each
fixed product the removed exponent is fixed; hence collecting the
remaining coefficients defines a formal series in \(P_j\).
There are only finitely many products, so this is an actual finite
expression for \(F\).  Applying \(\pi\), every removed factor
belongs to \(\mathfrak m_B\), by RTC6 and RTC7.  The expression
therefore puts \(x\) in \(\mathfrak m_B^N R_j\).
This argument does not appeal to closedness of that module or
assume the desired comparison of topologies. \hfill\(\square\)
% END RTC8

% BEGIN RTC9
\paragraph{RTC9. Full continuity in the original entry-algebra topology.}
\label{claim:RTC9}
For either original rank, \(u\) is adically continuous for the
ideal \(x_0=\mathfrak m_BA\).  The diagonal entries
\(a_{B,g,i}\) and character values \(\psi_{B,v,i}(s)\) are
adically continuous for \(\mathfrak m_B\).  Moreover, for every
\(j,N\) there is \(M\) such that
\[
 a\in\operatorname{gr}_jA,quad a\in\mathfrak m_R^M
       \quad\Longrightarrow\quad
 a\in\mathfrak m_B^N\operatorname{gr}_jA.
                                                        \tag{RTC.15}
\]
These are all the original \(\mathrm{cont\_\allowbreak u}\) conclusions.

\emph{Proof.}
The last assertion is RTC8 with the identifications in RTC7.
Fix an entry position \((i,j)\) and an exponent \(N\), and
apply RTC8 to its fixed weight.  Original continuity of the
universal representation in \(R\) says that, in a neighborhood
of any given group element \(g\), the difference
\(r(h)_{ij}-r(g)_{ij}\) lies in \(\mathfrak m_R^M\).
Both entries have the same weight, so (RTC.15) puts this difference
in \(\mathfrak m_B^N\operatorname{gr}_{w_i-w_j}A\), and hence
in \(x_0^N\).  This is precisely local constancy modulo
\(x_0^N\).

Diagonal entries and all character values have weight zero.
RTC8 applied to their differences puts those differences in
\(\mathfrak m_B^N R_0=\mathfrak m_B^N\).  Their original
continuity in \(R\) is part of the supplied lift.  This proves
all three continuity assertions, for every exponent including zero,
and supplies the full fixed-weight cofinality clause rather than
only its first residual quotient. \hfill\(\square\)
% END RTC9

% BEGIN RTC10
\paragraph{RTC10. Arbitrary fiber points are locally constant and determined by entries.}
\label{claim:RTC10}
Let \(Q\) be any commutative \(k\)-algebra, including a ring with
nilpotents or the zero ring.  The original condition
\(\mathrm{FPt}(\xi)\) on a ring map \(\xi:A\to Q\) is
\[
 x_0\subset\ker\xi,\qquad
 \xi(a)=\bar a\quad(a\in O),
                                                        \tag{RTC.16}
\]
where the right side uses the specified \(k\)-algebra structure.
Such a map gives a normalized representation whose diagonal entries
and local characters are exactly \(\Psi^{w(n)_i}\), and whose
matrix entries are locally constant.  If another fiber point agrees
with it on all matrix entries, the two ring maps are equal.

\emph{Proof.}
For \(b\in B\), choose \(a\in O\) with the same residue in
\(k\).  Then \(b-a\in\mathfrak m_B\), so (RTC.16) forces
\[
 \xi(b)=\operatorname{res}_R(b)\quad\text{in }Q.
                                                        \tag{RTC.17}
\]
This argument uses the actual residue of \(b\) in \(R\), and
works for every prescribed \(k\)-algebra structure on \(Q\).
The diagonal entries and character values belong to \(B\), with
their prescribed residues.  Equation (RTC.17) therefore gives
exactly the two clauses called \(\mathrm{IsBar}\); it does not
force the off-diagonal entries to vanish.

By RTC9, each universal matrix entry is locally constant modulo
\(x_0\).  The map \(\xi\) factors through that quotient,
so its image is locally constant in \(Q\), without imposing
any topology on \(Q\).  Finally \(A\) is generated by \(B\)
and the matrix entries.  All fiber points agree on \(B\) by
(RTC.17).  Agreement on the entries then gives agreement on the
entire generated algebra.  This proves the complete original
\(\mathrm{smoothCharacter\_\allowbreak extends\_\allowbreak to\_\allowbreak open\_\allowbreak subgroup}\)
statement at both allowed ranks. \hfill\(\square\)
% END RTC10

\subsection{Strict extension cocycles and their actual points}

\paragraph{The strict linear space.}
For a \(k\)-vector space \(M\), the original
\(F.\mathrm{strS}(M)\) consists of the functions
\(a:G_0\to M\) satisfying local constancy and
\[
 \begin{aligned}
 a(gh)&=\Psi(h)a(g)+\Psi(g)^{-1}a(h),\\
 a(g^*)&=a(g),\\
 a(\operatorname{loc}_{0,v}s)
   &=\bigl(\Psi(\operatorname{loc}_v s)
               -\Psi(\operatorname{loc}_v s)^{-1}\bigr)\gamma_v
       \quad\text{for some }\gamma_v\in M.
 \end{aligned}                                             \tag{RTC.18}
\]
All scalars here act through \(k\).  These conditions are linear,
including the existence of the local scalars \(\gamma_v\).
For each \(v\), the original \(f3\) field supplies \(t_v\)
with
\[
 d_v=\Psi(\operatorname{loc}_v t_v)
             -\Psi(\operatorname{loc}_v t_v)^{-1}\ne0.
                                                        \tag{RTC.19}
\]

For RTC11--RTC14, we work more generally with any field \(k\) of
characteristic \(\ell\) and any frame \(F\) satisfying
\(F.\mathrm{Ok}(C_0)\), without choosing a coefficient ring \(O\)
or a space \(S\).  The original \(\mathrm{Ok}\) premise itself
includes finiteness of \(k\).  The proofs below use the stated
frame properties and do not require a presentation of \(k\) as
the residue field of a particular ring.

% BEGIN RTC11
\paragraph{RTC11. A nonzero strict class has a normalized cocycle nonzero on \(G_0\).}
\label{claim:RTC11}
Let \(\chi:\Gamma\to k^\times\) be any character with
\(\chi(c)=-1\), nontrivial on every original local image.
If the original strict cone cohomology
\(H^1_C(\operatorname{loc},k(\chi),0)\) is nonzero, there
is a locally constant cocycle \(x:\Gamma\to k\) with
\[
 \begin{aligned}
 x(gh)&=x(g)+\chi(g)x(h),& x(c)&=0,\\
 x(\operatorname{loc}_v s)&=(\chi(\operatorname{loc}_v s)-1)Y_v
       &&\text{for some }Y_v\in k,
 \end{aligned}                                             \tag{RTC.20}
\]
and \(x(g)\ne0\) for at least one \(g\in G_0\).
No continuity of \(\chi\) is added to this assertion.

\emph{Proof.}
A nonzero class is represented by a pair \((z,(y_v)_v)\),
where \(z\) is a locally constant cocycle and
\(z(\operatorname{loc}_v s)=(\chi(\operatorname{loc}_v s)-1)y_v\).
The strict boundaries are exactly
\((g\mapsto(\chi(g)-1)X,(X)_v)\).  This is the original
cone definition, including the local primitives.

Since \(\ell\ne2\), put \(X=-z(c)/2\) and subtract this
boundary.  The new cocycle is
\[
 x(g)=z(g)-(\chi(g)-1)X
      =\frac{z(g)+z(gc)-z(c)}2.
                                                        \tag{RTC.21}
\]
The second formula follows from the cocycle identity.  It proves
local constancy directly, using continuity of right multiplication
by \(c\); thus no continuity assumption on \(\chi\) is hidden
in the normalization.  The cocycle identity gives \(z(1)=0\),
and \(c^2=1\) shows \(x(c)=0\).  The local primitives become
\(Y_v=y_v-X\), giving (RTC.20).

Suppose that \(x\) vanished on \(G_0\).  If
\(\delta(g)=-1\), then \(gc\in G_0\) and
\(g=(gc)c\).  The cocycle identity and \(x(c)=0\) imply
\(x(g)=0\) as well.  Hence \(x\) would vanish on all of
\(\Gamma\).  At each local place, choose \(s\) with
\(\chi(\operatorname{loc}_v s)\ne1\).  Equation (RTC.20)
then forces \(Y_v=0\).  The original pair would be the
boundary of \(X\), contradicting its nonzero class.  This proves
the assertion, with the full original scope of
\(\mathrm{smooth\_\allowbreak vector\_\allowbreak fixed\_\allowbreak by\_\allowbreak open\_\allowbreak subgroup}\).
\hfill\(\square\)
% END RTC11

% BEGIN RTC12
\paragraph{RTC12. The frame supplies one nonzero strict extension.}
\label{claim:RTC12}
There exists a nonzero \(b\in F.\mathrm{strS}(k)\).  It can
be chosen once, before any point of the core is considered.

\emph{Proof.}
Use
\(\chi=\Psi^{-2}\delta_k\), where \(\delta_k\) is the
original sign character with values in \(k^\times\).  Since
\(\Psi(c)=1\), \(\chi(c)=-1\).  On a local image contained
in \(G_0\), \(\chi=\Psi^{-2}\); (RTC.19) implies its
nontriviality there.  The original \(f4'\) field says precisely
that the strict cone cohomology for this \(\chi\) is nonzero.
Apply RTC11 and define
\[
 b(g)=x(g)\Psi(g),\qquad g\in G_0.
                                                        \tag{RTC.22}
\]
This is locally constant because \(x\) is locally constant and
\(\Psi\) has open kernel.  On \(G_0\), \(\chi(g)=\Psi(g)^{-2}\),
so multiplication of the cocycle identity by \(\Psi(gh)\)
gives the first equation of (RTC.18).

The inverse cocycle identity gives
\(x(g^{-1})=-\Psi(g)^2x(g)\).  Since \(x(c)=0\) and
\(\chi(c)=-1\), it follows that
\(x(g^*)=-x(g^{-1})=\Psi(g)^2x(g)\).  Together with
\(\Psi(g^*)=\Psi(g)^{-1}\), this proves \(b(g^*)=b(g)\).
Locally, (RTC.20) gives
\[
 b(\operatorname{loc}_{0,v}s)
  =(\Psi(\operatorname{loc}_v s)^{-1}
       -\Psi(\operatorname{loc}_v s))Y_v,
\]
which is the last condition of (RTC.18), with \(\gamma_v=-Y_v\).
Finally a nonzero value of \(x\) on \(G_0\) remains nonzero
after multiplication by \(\Psi(g)\).  This proves the full
original nonzero-\(\mathrm{strS}\) supplier. \hfill\(\square\)
% END RTC12

% BEGIN RTC13
\paragraph{RTC13. Finite coefficient expansion of a strict cocycle.}
\label{claim:RTC13}
For every \(k\)-vector space \(M\) and every
\(a\in F.\mathrm{strS}(M)\), there are finitely many
\(b_1,\ldots,b_t\in M\) and
\(a_1,\ldots,a_t\in F.\mathrm{strS}(k)\) such that
\[
 a(g)=\sum_{i=1}^t b_i a_i(g)\quad(g\in G_0).
                                                        \tag{RTC.23}
\]
The \(b_i\) may be taken to be a basis of the span of the values
of \(a\).  No finite-dimensional hypothesis on \(M\) is needed.

\emph{Proof.}
The locally constant function \(a\) has finite image: its open
fibers cover the compact group \(G_0\), so a finite subcover
suffices.  Let \(M_0\) be their finite-dimensional span and choose
a basis \(b_i\).  Write its coordinate functions as \(a_i(g)\).
They are locally constant, and the multiplication and star identities
in (RTC.18) descend to each coordinate by linear independence.

The local primitive in (RTC.18) also belongs to \(M_0\), because
(RTC.19) gives
\(\gamma_v=d_v^{-1}a(\operatorname{loc}_{0,v}t_v)\).
Applying the coordinate functionals therefore gives the required
local primitive for every \(a_i\).  This proves (RTC.23), including
the case \(a=0\), where the basis may be empty.  It also proves
that applying any \(k\)-linear map to a strict cocycle gives a
strict cocycle.  In particular, addition and multiplication by an
element of a \(k\)-algebra preserve this space. \hfill\(\square\)
% END RTC13

% BEGIN RTC14
\paragraph{RTC14. Strict cocycles are exactly the normalized triangular representations.}
\label{claim:RTC14}
For a commutative \(k\)-algebra \(Q\), every
\(a\in F.\mathrm{strS}(Q)\) defines the actual normalized
rank-two representation
\[
 r_a(g)=
 \begin{pmatrix}\Psi(g)&0\\a(g)&\Psi(g)^{-1}\end{pmatrix},
 \qquad
 \psi_{v,0}(s)=\Psi(\operatorname{loc}_v s),\quad
 \psi_{v,1}(s)=\Psi(\operatorname{loc}_v s)^{-1}.
                                                        \tag{RTC.24}
\]
It satisfies \(\mathrm{IsBar}\), has zero upper-entry ideal,
and has locally constant matrix entries.  Conversely, every normalized
rank-two representation with these three properties has its lower
entry in \(F.\mathrm{strS}(Q)\).

\emph{Proof.}
The multiplication identity for \(a\) is precisely the lower-left
entry of \(r_a(gh)=r_a(g)r_a(h)\).  At \(g=h=1\) it gives
\(a(1)=0\); hence the identity matrix is respected, and the
group law also gives the inverses.  The equality
\(a(g^*)=a(g)\) and the inverse relation for \(\Psi\) give
\(r_a(g^*)=J_2r_a(g)^tJ_2\).  They also give the prescribed
off-diagonal normalization at \(F.g_0\).

Write \(p_s=\Psi(\operatorname{loc}_v s)\),
\(q_s=p_s^{-1}\), and \(a_s=a(\operatorname{loc}_{0,v}s)\).
The only possibly nonzero entry of the ordered product
\((r_a(s)-p_sI)(r_a(t)-q_tI)\) is its lower-left entry,
\[
 a_s(p_t-q_t)-(p_s-q_s)a_t.
                                                        \tag{RTC.25}
\]
The local formula in (RTC.18) makes (RTC.25) zero.  This checks
the full ordered identity for every pair \((s,t)\), with the
original order.  The characters and diagonal entries in (RTC.24)
are exactly the prescribed ones; local constancy follows from the
open kernel of \(\Psi\), local continuity and the defining
local constancy of \(a\).

Conversely, \(\mathrm{IsBar}\) fixes the two diagonal entries
and both local characters, while the vanishing upper-entry ideal
sets the upper entry to zero.  Matrix multiplication and the
polarization identity give the first two identities of (RTC.18).
In (RTC.25), choose \(t=t_v\) from (RTC.19).  The scalar
\(d_v=p_{t_v}-q_{t_v}\) is a unit in every \(k\)-algebra.
Thus
\(a_s=(p_s-q_s)d_v^{-1}a_{t_v}\), proving the last identity.
This includes the zero target ring.  The converse is the full
original \(\mathrm{framePoint\_\allowbreak fiber\_\allowbreak nonempty}\) assertion.
\hfill\(\square\)
% END RTC14

From RTC15 through RTC20, we return to the supplied coefficient
ring \(O\), its specified residue field \(k\), and the original
rank-two space \(S\) of the chapter's fixed data.

% BEGIN RTC15
\paragraph{RTC15. Realization over every prescribed target algebra.}
\label{claim:RTC15}
For the supplied rank-two space \(S\), every commutative
\(k\)-algebra \(Q\) and every \(a\in F.\mathrm{strS}(Q)\)
admit a ring map \(\xi_a:A\to Q\) satisfying
\[
 \mathrm{FPt}(\xi_a),\qquad
 \mathfrak t_-\subset\ker\xi_a,\qquad
 \xi_a(u(g)_{10})=a(g)\quad(g\in G_0).
                                                        \tag{RTC.26}
\]
Its entire normalized representation is (RTC.24).  This is the
original \(\mathrm{pt\_\allowbreak ex}\) supplier; together with RTC14 it
also supplies the rank-two case of
\(\mathrm{smoothHom\_\allowbreak weight\_\allowbreak sum\_\allowbreak neg}\) at its original scope.

\emph{Proof.}
Use RTC13 to write \(a=\sum_i b_i a_i\), and form the
coefficient algebra
\[
 C=k[[T_1,\ldots,T_t]],\qquad
       a_C(g)=\sum_i T_i a_i(g),
                                                        \tag{RTC.27}
\]
with its \(O\)-structure through \(O\to k\).  The function
\(a_C\) is a strict cocycle.  The representation (RTC.24) over
\(C\) is an original continuous lift: its lower entries lie in
the maximal ideal, its upper entries are zero, and its diagonal
entries and characters are the specified residue characters.
Every entry is locally constant even before taking adic quotients,
since it is a finite combination of locally constant functions.
Universality supplies an \(O\)-map \(f:R\to C\) taking the
universal normalized representation to this one.

We claim that \(f(B)\subset k\).  A weight-zero entry monomial
containing an off-diagonal factor has image zero.  Indeed, all
nonzero off-diagonal entries in (RTC.24) have negative weight;
a zero-weight product containing such a factor must also contain
a positive-weight, hence zero, upper entry.  Products containing
only diagonal factors, and every character value, have images in
the constant field \(k\).  Thus \(f(B_0)\subset k\).
For \(b\in B\), approximate it modulo \(\mathfrak m_R^N\)
by elements of \(B_0\).  RTC1 maps these approximations modulo
\(\mathfrak m_C^N\).  For \(N\geq1\), their constant images
are necessarily the same constant, namely the residue of \(b\).
Separatedness of \(C\) proves
\(f(b)=\operatorname{res}_R(b)\in k\).

It follows that \(f(A)\) lies in the polynomial subring
\(k[T_1,\ldots,T_t]\subset C\): the other generators are the
matrix entries, which are constants or the linear expressions
in (RTC.27).  Because the inclusion of this polynomial ring in
the formal-series ring is injective, restriction of \(f\) gives
a well-defined ring map from \(A\) to the polynomial ring.
Now, and only now, evaluate \(T_i\) at the prescribed elements
\(b_i\in Q\).  The resulting map is \(\xi_a\).

It kills \(\mathfrak m_B\), has the prescribed \(O\)-map,
and takes every matrix entry and character value to (RTC.24).
This proves all clauses of (RTC.26), including every full local
character value and its inverse.  Nothing requires \(Q\) to be
local, Noetherian, finite over \(k\), reduced, or a field.
An empty coefficient list uses \(C=k\) and the same argument.
In particular the proof does not invoke the more general
\(\mathrm{point\_\allowbreak loc}\) construction and does not apply ACP's
diagonal-residual shortcut to a triangular residual point.
\hfill\(\square\)
% END RTC15

% BEGIN RTC16
\paragraph{RTC16. The strict point correspondence and its functoriality.}
\label{claim:RTC16}
For every commutative \(k\)-algebra \(Q\), the map
\[
 \{\xi:A\to Q:\mathrm{FPt}(\xi),\
                   \mathfrak t_-\subset\ker\xi\}
   \longrightarrow F.\mathrm{strS}(Q),\qquad
       \xi\longmapsto(g\mapsto\xi(u(g)_{10}))
                                                        \tag{RTC.28}
\]
is a bijection.  It commutes with every \(k\)-algebra map
\(Q\to Q'\).

\emph{Proof.}
RTC10 supplies \(\mathrm{IsBar}\) and local constancy for the
mapped normalized representation.  The upper-entry hypothesis
then permits the converse in RTC14, so (RTC.28) is well-defined.
Surjectivity is RTC15.  If two points have the same lower entry,
their upper entries are zero and their diagonal entries are the
same prescribed characters.  They therefore agree on every matrix
entry; RTC10 proves equality of the ring maps.  This is the
original \(\mathrm{pt\_\allowbreak uniq}\) argument with its entire
coefficient-ring scope, and the well-definedness just proved is
the original \(\mathrm{pt\_\allowbreak mem}\) argument.

Composition with a \(k\)-algebra map preserves (RTC.16), the
vanishing upper-entry ideal and the lower-entry values.  Applying
a \(k\)-linear map also preserves strict cocycles by RTC13.
These observations prove the asserted naturality, with no flatness
condition on the subsequent map. \hfill\(\square\)
% END RTC16

\subsection{The actual core and its polynomial lines}

% BEGIN RTC17
\paragraph{RTC17. Every core prime is a strict triangular point.}
\label{claim:RTC17}
For every \(\mathfrak p\in\mathcal C\), there is a field
\(K\) with its specified \(k\)-algebra structure and a fiber
point \(\xi:A\to K\) such that
\[
 \ker\xi=\mathfrak p,\qquad
 \mathfrak t_-\subset\mathfrak p,\qquad
 a(g)=\xi(u(g)_{10})\in F.\mathrm{strS}(K),\quad a\ne0.
                                                        \tag{RTC.29}
\]
The field \(K\) need not be finite or algebraic over \(k\).

\emph{Proof.}
Take \(K=\operatorname{Frac}(A/\mathfrak p)\) and let
\(\xi\) be the quotient map followed by the field inclusion.
Its kernel is \(\mathfrak p\).  The map from \(O\) kills
\(\mathfrak m_O\), because its image in \(B\) is contained
in \(\mathfrak m_B\), and \(x_0\subset\mathfrak p\).
It consequently induces a field map \(k\to K\).  With this
precise scalar structure, \(\xi\) satisfies (RTC.16).

Since \(\mathfrak t_+\not\subset\mathfrak p\), choose
\(g\) for which \(u(g)_{10}\notin\mathfrak p\).
For every \(h\), the product
\(u(g)_{10}u(h)_{01}\) is a weight-zero entry monomial.
It belongs to \(B\), and its residue in \(R\) is zero,
since the original residual representation is diagonal.
It therefore lies in \(\mathfrak m_B\subset\mathfrak p\).
Primality forces \(u(h)_{01}\in\mathfrak p\) for every
\(h\).  This is exactly \(\mathfrak t_-\subset\mathfrak p\).

RTC10 now supplies local constancy and the full residual character
identities, and RTC14 puts the lower entry in
\(F.\mathrm{strS}(K)\).  The chosen \(g\) makes it nonzero.
Thus the rank-two case of the original
\(\mathrm{smoothDual\_\allowbreak exact}\) supplier is proved directly;
no higher-rank triangular-type dichotomy has been assumed.
\hfill\(\square\)
% END RTC17

% BEGIN RTC18
\paragraph{RTC18. A polynomial strict family has a common generic prime.}
\label{claim:RTC18}
Let \(K\) be any field over \(k\), and let
\(a_T\in F.\mathrm{strS}(K[T])\).  There is a prime
\(\eta\in\operatorname{Spec}A\), over \(\mathfrak m_B\),
whose ideal is contained in the kernel of every strict triangular
fiber point \(\xi:A\to K\) with
\[
 \xi(u(g)_{10})=a_T(g)|_{T=t}\quad(g\in G_0)
                                                        \tag{RTC.30}
\]
for some \(t\in K\).

\emph{Proof.}
Apply RTC15 with the actual target algebra \(K[T]\), obtaining
\(\xi_T:A\to K[T]\).  Put \(\eta=\ker\xi_T\).
The target is an integral domain, so this is a prime ideal;
the fiber condition puts \(x_0\) in it.  Evaluation at
\(T=t\) is a \(k\)-algebra map.  By RTC16 its composition
with \(\xi_T\) is the unique strict triangular fiber point
with the values (RTC.30), and hence equals \(\xi\).
Its kernel contains \(\ker\xi_T\), proving the assertion.
This is the actual \(\mathrm{line}\) supplier, including its
universal quantification over the later evaluation point and map.
\hfill\(\square\)
% END RTC18

% BEGIN RTC19
\paragraph{RTC19. Full original preconnectedness of the rank-two core.}
\label{claim:RTC19}
For every original \(O,F,C_0\), every proof of
\(F.\mathrm{Ok}(C_0)\), and every supplied rank-two space
\(S\), its actual set \(S.\mathrm{BB}\) is preconnected.
In fact it is nonempty.  This is the full original theorem
\(\mathrm{preconnected\_\allowbreak rank\_\allowbreak two}\), with no
\(\mathrm{Big}\), modularity or finite-point hypothesis added.

\emph{Proof.}
Choose once the nonzero \(b\in F.\mathrm{strS}(k)\) of
RTC12 and realize it by \(\xi_b:A\to k\).  Its kernel
\(\mathfrak q_b\) is prime and contains \(x_0\).
Some lower-entry value \(b(g)\) is nonzero, so
\(\mathfrak t_+\not\subset\mathfrak q_b\).
Consequently \(\mathfrak q_b\in\mathcal C\).

Now take any \(\mathfrak p\in\mathcal C\), and use its
actual residue-field point \(\xi:A\to K\) from RTC17.
Let \(a(g)=\xi(u(g)_{10})\), and denote the image of \(b\)
in \(K\) by \(b_K\).  In the prescribed algebra \(K[T]\)
form
\[
 a_T(g)=(1-T)a(g)+Tb_K(g).
                                                        \tag{RTC.31}
\]
The scalar-extension, addition and scalar-multiplication assertions
of RTC13 show that this is a strict cocycle.  Let \(\eta\)
be the prime furnished by RTC18.  At \(T=0\), RTC16 identifies
the evaluated point with \(\xi\); hence
\(\eta\subset\mathfrak p\).  At \(T=1\), it identifies
the point with the composition \(A\xrightarrow{\xi_b}k\to K\).
Since the field map \(k\to K\) is injective, its kernel is
\(\mathfrak q_b\).  Thus
\[
 \eta\subset\mathfrak p,\qquad
 \eta\subset\mathfrak q_b,
 \qquad x_0\subset\eta.
                                                        \tag{RTC.32}
\]
The containment \(\eta\subset\mathfrak p\) and the fact
that \(\mathfrak p\) avoids \(\mathfrak t_+\) imply that
\(\eta\) also avoids \(\mathfrak t_+\).  Hence all three
primes in (RTC.32) lie in the original core.

For the Zariski topology,
\(\overline{\{\eta\}}=V(\eta)\).  The three-point subspace
\(\{\eta,\mathfrak p,\mathfrak q_b\}\) has \(\eta\) as
a dense point and is therefore irreducible, in particular
preconnected.  Indeed two nonempty disjoint open subsets covering
such a space cannot both meet its dense singleton.  As
\(\mathfrak p\) varies, these preconnected subsets cover
\(\mathcal C\) and all contain \(\mathfrak q_b\).
A separation of \(\mathcal C\) into two disjoint nonempty
relative open sets would put \(\mathfrak q_b\) in one of them;
preconnectedness of every three-point subset would put all its
points in that same set.  The other set would be empty.

This proves preconnectedness and the asserted nonemptiness.
The argument does not assert that every specialization of the
polynomial family (RTC.31) remains in the core; a specialization
can have zero lower entry.  Only its generic prime and the two
specified endpoints are used, and their membership was proved
above.  The anchor was chosen before \(\mathfrak p\), and
every construction took place on the same \(F,S,A,B\).
\hfill\(\square\)
% END RTC19

% BEGIN RTC20
\paragraph{RTC20. Installation in the actual rank-two propagation consumer.}
\label{claim:RTC20}
In the original \(\mathrm{ComponentsGoodRankTwo}\) construction,
the ambient Noetherianity and the full preconnected-core field of
the global propagation data are now supplied on the same actual
space.  Its open type-I locus and closed good locus have the
literal formulas
\[
 \mathrm{TI}
 =D(\mathfrak t_+)\cap D(\mathfrak t_-),\qquad
 \mathrm{Good}=V(MA)
                                                        \tag{RTC.33}
\]
for the original modularity ideal \(M\subset B\).
These facts can be passed to ACP's already proved abstract
propagation theorems without any change of rings or ideals.

\emph{Proof.}
The consumer takes \(X=\operatorname{Spec}A\), the above
\(\mathrm{TI}\) and \(\mathrm{Good}\), and
\(\mathrm{Core}=S.\mathrm{BB}\).  RTC7 proves that its actual
ring \(A\) is Noetherian, so \(X\) is Noetherian.  RTC19
is precisely the preconnectedness field used in the consumer.
In rank two, the extra rank-three exceptional-type clause in
\(\mathrm{TI}\) is vacuous, giving its first formula in
(RTC.33).  An ideal \(M\subset B\) lies in the contraction
of a prime of \(A\) if and only if its extended ideal lies
in that prime, giving the second formula.  These descriptions
prove openness and closedness respectively.

The separate local chart flatness and representation dictionary,
sharp component chains, nice-prime dimension bounds, closed-point
density and generic-component incidence, and seed-prime comparison
are still needed for the other original fields.  The current
installation supplies neither those fields nor \(\mathrm{AllGood}\).
It removes the complete connectedness supplier and the generic
weighted-algebra/continuity suppliers from that list.  In an
arithmetic application the original \(P\) used to construct the
frame remains the AFE/LFC premise; no new premise is substituted
for it here. \hfill\(\square\)
% END RTC20

\subsection{Returning finite points to the original coefficient frame}

% BEGIN RTC21
\paragraph{RTC21. Coefficient fields in the finite-residue test DVRs.}
\label{claim:RTC21}
Let \(A_*\) be a complete discrete valuation domain of positive
characteristic, with finite residue field \(k_*\) of size \(q\).
There is a unique section \(s:k_*\to A_*\) of the residue map
whose values satisfy \(s(a)^q=s(a)\).  Every given embedding
of a finite field into \(A_*\) agrees with this section on its
residual image.  If \(A_*\) is not a field and \(t\) is a
uniformizer, evaluation gives an isomorphism
\[
 k_*[[T]]\ \xrightarrow{\ \sim\ } A_*,
       \qquad T\longmapsto t,\quad a\longmapsto s(a).
                                                        \tag{RTC.34}
\]

\emph{Proof.}
Lift \(a\in k_*\) to \(a_0\in A_*\) and consider
\(a_m=a_0^{q^m}\).  Because \(a_0^q-a_0\in\mathfrak m_{A_*}\)
and the ring has the residue characteristic,
\[
 a_{m+1}-a_m=(a_0^q-a_0)^{q^m}
       \in\mathfrak m_{A_*}^{q^m}.
\]
The sequence converges by completeness; its limit lifts \(a\)
and is fixed by the \(q\)th-power map.  Two such lifts differ
by \(d\in\mathfrak m_{A_*}\) with \(d^q=d\).  Since
\(1-d^{q-1}\) is a unit, this forces \(d=0\).  Thus the
lift is independent of the first choice.  The \(q\)th-power
map is additive and multiplicative in this characteristic, so
uniqueness makes these lifts a field section.

If \(j:k\to A_*\) is an embedding of a finite field, its
residue map is a field embedding.  Thus \(k\) identifies with
a subfield of \(k_*\), and every \(j(a)\) satisfies
\(j(a)^q=j(a)\).  Uniqueness gives
\(j(a)=s(\operatorname{res}(j(a)))\).

In the nonfield case, extract residues successively.  Every
\(x\in A_*\) has a unique expansion
\[
 x=\sum_{m\geq0}s(a_m)t^m.
                                                        \tag{RTC.35}
\]
At each step subtract the lifted residue and divide the remainder
by \(t\); completeness proves existence.  A series with first
nonzero coefficient at index \(m\) is \(t^m\) times a unit,
so it cannot evaluate to zero.  This proves uniqueness and
injectivity as well as surjectivity of (RTC.34).  Addition and
multiplication follow by truncation, so it is the claimed ring
isomorphism.  If the valuation-ring convention permits the field
case, then \(A_*=k_*\) and the section is already an isomorphism;
that case will be handled separately below. \hfill\(\square\)
% END RTC21

% BEGIN RTC22
\paragraph{RTC22. A common test DVR for a prescribed finite coefficient extension.}
\label{claim:RTC22}
Let \(k,k_1\) be finite fields of the same characteristic,
let \(f:k\to k_1\) be a field map, and let \(j:k\to A_*\)
be the prescribed coefficient embedding in a complete finite-residue
test DVR.  There are a complete DVR \(B_*\) with finite residue
field, an injective local map \(\phi:A_*\to B_*\), and a
field map \(j_1:k_1\to B_*\) satisfying
\[
       \phi(j(a))=j_1(f(a))\qquad(a\in k).
                                                        \tag{RTC.36}
\]
This preserves the given maps; it does not replace them by
unrelated abstract embeddings.

\emph{Proof.}
The embedding \(j\) makes \(A_*\) have the same positive
characteristic as \(k\).  Let \(k_*\) be its residue field.
Inside an algebraic closure of \(k_*\), extend the prescribed
embedding of \(k\) to an embedding of \(k_1\), using its
finite algebraic extension and the factorization of its minimal
polynomial.  Let \(L\) be the field generated by \(k_*\) and
that image.  It is a finite extension of a finite field, hence
finite, and its two embeddings agree with the original maps on
\(k\).  This is the ordinary field-amalgamation step over the
specified common subfield.

If \(A_*\) is not a field, use the isomorphism (RTC.34) and
take \(B_*=L[[T]]\).  The inclusion of coefficients
\(k_*\subset L\) gives \(\phi\), with the same parameter
\(T\), and use the chosen \(k_1\)-embedding as constant
series for \(j_1\).  Equation (RTC.36) follows from the
compatibility assertion in RTC21.  Coefficientwise injectivity
proves injectivity of \(\phi\), and
\(\phi^{-1}((T))=(t)\) proves locality.  The power-series
ring \(B_*\) is a complete DVR with residue field \(L\).
In this case it is even finite free over \(A_*\), by a basis
of \(L\) over \(k_*\), but that stronger fact is not needed.

If the original test DVR is a field, take the same
\(B_*=L[[T]]\) and embed \(A_*=k_*\) as constants.
This map is injective and local, with zero preimage of \((T)\),
and (RTC.36) again holds.  This also shows why no extra finite
fraction-field-extension condition is being imposed on the
original test-ring quantifiers. \hfill\(\square\)
% END RTC22

% BEGIN RTC23
\paragraph{RTC23. Transport of the exact finite-point conditions.}
\label{claim:RTC23}
Let \(O=c.O\), \(O_1=c_1.O\), and
\(f:k=\operatorname{Res}(O)\to
 k_1=\operatorname{Res}(O_1)\) be the original coefficient
data of \(\mathrm{FinPtDown}\).  For every original
\(\mathrm{IsCD}(O,A_*)\) test with its specified scalar
tower \(O\to k\xrightarrow{j}A_*\), the ring \(B_*\)
in RTC22 has an \(\mathrm{IsCD}(O_1,B_*)\) structure, through
\(O_1\to k_1\xrightarrow{j_1}B_*\).  At every rank \(n\),
mapping a normalized \(\mathrm{IsLiftA}\) point through
\(\phi\) gives the required \(\mathrm{IsLiftA}\) point for
the frame \(\Phi.\mathrm{ext}(f)\), and preserves its original
local ordinary-character triviality condition.

\emph{Proof.}
The finite residue field of \(O_1\) is part of the original
coefficient datum.  Its map to \(B_*\) factors through that
field, so it sends precisely the nonunits of \(O_1\) to
nonunits.  Completeness, the domain and DVR properties, and the
finite residue field of \(B_*\) were proved in RTC22.  These
are exactly the \(\mathrm{IsCD}\) clauses.  The indicated
factorization is also the exact required scalar tower.

The frame extension changes only the residue character from
\(\Psi\) to \(f\Psi\); its groups, local maps, subgroups,
involution and distinguished elements are unchanged.  Therefore
the matrices and local characters mapped through \(\phi\)
still satisfy every normalization and ordered-product identity.
Since \(\phi\) is local,
\[
 \phi(\mathfrak m_{A_*}^N)\subset\mathfrak m_{B_*}^N
                                                        \tag{RTC.37}
\]
for every \(N\), which transports all entry and character
continuity.  Reduction commutes with \(\phi\).  For each
\(a\in k\), equation (RTC.36) identifies the reduction of
\(\phi(j(a))\) with that of \(j_1(f(a))\).
Apply this to \(\Psi(g)^{w(n)_i}\) and to each local value.
It gives the exact new diagonal and character residues, and
the zero off-diagonal residues remain zero.  This proves the
literal \(\mathrm{IsLiftA}\) condition, not only its
semisimplification.

Finally suppose that an original inertia element satisfies
\((f\Psi)(\operatorname{loc}_v s)=1\).  A field map is
injective, so \(\Psi(\operatorname{loc}_v s)=1\).
The original point condition then makes its ordinary character
value equal to one, and applying \(\phi\) preserves that
equality.  No smaller collection of local tests has been substituted.
\hfill\(\square\)
% END RTC23

% BEGIN RTC24
\paragraph{RTC24. Full original descent of finite points.}
\label{claim:RTC24}
For every prime \(\ell\), original coefficient data
\(c,c_1,f\), inclusion \(c.O\subset c_1.O\), and the
original compatibility of \(f\) with their residue maps,
every frame \(\Phi\) and every rank \(n\) satisfy
\[
 \mathrm{FinPt}\bigl(c_1.O,\Phi.\mathrm{ext}(f),n\bigr)
       \ \Longrightarrow\
 \mathrm{FinPt}(c.O,\Phi,n).
                                                        \tag{RTC.38}
\]
This is the full original \(\mathrm{FinPtDown}\) theorem.
There is no \(F.\mathrm{Ok}\) hypothesis, odd-prime
restriction, or restriction to ranks two and three in this claim.

\emph{Proof.}
Take any original test ring \(A_*\), scalar tower, normalized
point \(\rho\), and its local ordinary-character condition.
Use RTC22 and RTC23 to obtain the actual allowed \(B_*\)-point
for the extended frame.  Apply the assumed extended-frame
\(\mathrm{FinPt}\) to this point.  For every \(g\), it gives
\[
 \operatorname{charpol}(\phi(\rho(g)))
     =\prod_{i=0}^{n-1}
         \bigl(X-j_1(f(\Psi(g)^{w(n)_i}))\bigr).
                                                        \tag{RTC.39}
\]
Characteristic polynomials commute with every ring map, by the
determinant formula.  Equation (RTC.36) identifies the right side
of (RTC.39) with the image under \(\phi\) of
\(\prod_i(X-j(\Psi(g)^{w(n)_i}))\).  Since \(\phi\) is
injective, the induced map on polynomial rings is coefficientwise
injective.  Cancel it to obtain precisely the original finite-point
characteristic-polynomial identity over \(A_*\).

The test ring, scalar tower, point and \(g\) were arbitrary,
so this proves (RTC.38) at all its original quantifiers.  Empty
matrix size causes no exception: both determinant and empty product
are one.  The stated inclusion and residue compatibility of the
original coefficient data have been retained; the proof in fact
uses only their prescribed field map and its commutation in
(RTC.36), as does the actual transport call. \hfill\(\square\)
% END RTC24

% BEGIN RTC25
\paragraph{RTC25. The finite-end join keeps the original frame.}
\label{claim:RTC25}
In the original linked arithmetic construction, suppose that the
same extended coefficient frame has supplied a rank-three space
and its actual \(\mathrm{OffFin}\) conclusion.  Then ACP11,
RTC24 and AFE20 return the original-frame \(\mathrm{FinEnd}\)
on the same marking \(g_1\).  The finite extension determining
the subgroup in \(\mathrm{FinEnd}\) is fixed before every
later ordinary point.  The rank-two/rank-three rings and modularity
ideals are not identified in this join.

\emph{Proof.}
The actual two-run construction returns the finite-point conclusion
\[
 \mathrm{FinPt}(c_1.O,\Phi.\mathrm{ext}(f),3).
\]
At its supplied \(\mathrm{Space}\) and \(\mathrm{OffFin}\),
ACP11 supplies exactly this statement.  RTC24 transports it to
\(\mathrm{FinPt}(c.O,\Phi,3)\), which is the original
\(\mathrm{finPtDown}\) call in \(\mathrm{LinkedEndE}\).
AFE20 now applies on that original coefficient frame, using its
unchanged link and distinguished element.  Its proof chooses the
subgroup \(G_L\) before the later ordinary-point tests, and
RTC24 did not alter the groups or the marking.

The actual rank-two \(\mathrm{AllGood}\), rank-three
\(\mathrm{OffFin}\), spaces, modularity ideals, automorphic
seeds and arithmetic-point classicality are still original
suppliers.  RTC19 supplies the whole rank-two connectedness
theorem, not those other conclusions.  The original \(P\)
premise at the AFE frame entry and the Host case \(\ell=3\)
are unchanged.  ACP19's projection from the Alt endpoints also
remains a projection only: the terminal \(\mathrm{WitA}\)
and \(\mathrm{WitAQ}\) fields still need their actual producers.
Thus the coefficient-descent gap is closed while the larger
arithmetic endpoints and selected whole obligations remain open.
\hfill\(\square\)
% END RTC25

\paragraph{What the next arithmetic task can now use.}
The original generic weighted presentation, degree-zero coefficient
ring, graded entry-algebra finiteness and topology comparison are
available in both ranks, directly from the supplied universal space.
The full rank-two core is now preconnected, and the original
coefficient descent of finite points is available at arbitrary rank.
The next local-chart proof can use RTC7 and RTC9 upstream; it must
still establish the actual completion from the local universal
property and its flatness, rather than importing a later
\(\mathrm{FinA}\) conclusion.  The remaining sharp dimensions,
closed-point and seed incidence, rank-two-to-three corner, modularity
and automorphic constructions are separate mathematical suppliers on
the same original data.  No full rank-specific propagation,
\(\mathrm{GdEnd}\), \(\mathrm{HostEnd}\), Alt endpoint,
whole selected frontier obligation, or complete human RH/Dirichlet-GRH
proof is asserted by this chapter.


\section{The exact archimedean orbital scalar}
\label{sec:archimedean-orbital-scalar-v54}

This chapter computes an orbital measure for the actual setting.  It
does not choose another Haar measure at infinity.  The conclusion is
the original \texttt{Push\allowbreak Inf}, with the literal constant
\[
 K_X(u')=\operatorname{gaugeU}(X.B.I)\,
       \operatorname{volU}(k)^d\,T_w(u')^{-kd},
 \qquad
 \operatorname{volU}(k)=\prod_{j=1}^{k}\frac{(2\pi)^j}{(j-1)!}.
 \tag{AOS.1}
\]
Here $k=D.R.k$, $d=[F:\mathbb Q]$ is the number of real places of
$F=X.B.I.F$, and $T_w(u')$ is the source's positive rational integer
viewed in $\mathbb R$.  The proof explains separately the determinant,
the compact-unitary volume, and the character normalization.  Their
combination leaves no unspecified proportionality constant.

The compact-stabilizer argument also gives the original
\texttt{Ex\allowbreak Inf}, \texttt{Htot\allowbreak Inf}, \texttt{Weil\allowbreak Inf},
\texttt{Ainf\allowbreak Value}, and \texttt{Aprime\allowbreak Inf}.  These are archimedean
statements on a supplied setting.  Arithmetic quotient-volume
comparisons, the finite-place factors, and the ambient automorphic
family remain separate inputs to their later consumers.

\subsection{The supplied objects and the normalization to be computed}

Fix any original datum $D$, a Hida setting $X$, its actual adelic bridge
$h_b$, two sets $F_{dN},F_{dN'}$ in $X.B.I.G_\infty$, and words
$u,u'$.  The original hypothesis \texttt{Hyp\allowbreak X} says that $X$ has its
measure ties through this bridge, its level subgroup is normalized,
and the two specified sets are the required unipotent fundamental
domains.  In particular, it includes
\[
 \mu_{G,\infty}
   =(\operatorname{toInf}_{h_b})_*
           \left(\bigotimes_{v\mid\infty}\mu_{G,v}^{\rm old}\right),
 \qquad
 (\iota_H)_*\mu_{H,\infty}
    =(\operatorname{toInf}_{h_b})_*\mu_H(u').
 \tag{AOS.2}
\]
The first measures on the right are the original
\texttt{Formalism.\allowbreak REAL.\allowbreak mu\allowbreak G}; the second equation is the source's
ambient-measure tie for the actual stabilizer.  We retain the other
parts of \texttt{Hyp\allowbreak X}, even when an individual calculation needs only
these two equations.

Put $E=X.B.I.E$, $r=D.R.rS$, $a=2r+1$ and $n=X.B.I.n=k+a$.
The rational ordered tuple under consideration is
\[
 Y_0=(e_1,\ldots,e_r,z_{u'},f_1,\ldots,f_r),
 \qquad T_0=Y_0^*JY_0.
 \tag{AOS.3}
\]
The order in (AOS.3) is precisely \texttt{tup\allowbreak H(frame\allowbreak Of X u')};
we have merely written its vectors as columns.  At infinity let
$V_\infty=\operatorname{Mat}_{n\times a}(E\otimes_F F_\infty)$,
$\mathcal H_\infty=\operatorname{Herm}_a(E\otimes_F F_\infty)$,
$U_\infty$ be the independent-column locus, and
$\gamma(Y)=Y^*JY$.  The actual orbit map is
\[
 \pi_{Y_0}:G_\infty\longrightarrow V_\infty,
       \qquad g\longmapsto gY_0;
 \tag{AOS.4}
\]
it is \texttt{orb\allowbreak M X u'} in the source.

Let $\psi$ be any character satisfying the original
\texttt{Meas\allowbreak Compat\,$\psi$} and \texttt{Is\allowbreak Std\allowbreak Char}.  Its infinite
factor is
\[
       \psi_\infty(t)=\exp\left(2\pi i\sum_{v\mid\infty}t_v\right).
 \tag{AOS.5}
\]
For the setting obtained from GMR19 this is the one character already
chosen before all words, places and auxiliary tuples.  We do not
choose a character again in any calculation below.

The measures $\mu_\infty,\nu_\infty$ are \emph{arbitrary} witnesses of
the source's \texttt{Sd\allowbreak W} and \texttt{Sd\allowbreak Herm}.  Similarly
$\omega_T$ is any family satisfying \texttt{Is\allowbreak Fibre\allowbreak Family} for these
measures, the Gram map and $U_\infty$.  Thus it has the support,
disintegration and continuity properties stated in (GMR.3), on the
whole Hermitian parameter space.  GMR2--GMR5 prove existence and
uniqueness at this exact scope, including the finite product of real
places.  They will identify arbitrary witnesses with the explicit
densities calculated below.

Finally write $J_v=\operatorname{diag}(\epsilon_{v,1},\ldots,
\epsilon_{v,n})$, with $\epsilon_{v,i}\in\{1,-1\}$, for the original
real model.  Its Haar measure is $\mu_{G,v}^{\rm old}$, and its
source gauge is $g_v=\operatorname{Gauge.gaugeU}(\epsilon_v)$.
The original antecedent of \texttt{Push\allowbreak Inf} is
$\operatorname{GaugeDens}(\epsilon_v)$ for every $v$.
AOS5 also proves this universal gauge-density assertion directly;
in particular $g_v>0$.  The product in (AOS.1) is
$\operatorname{gaugeU}(X.B.I)=\prod_v g_v$.

\subsection{The two self-dual measures and the real Jacobian}

% BEGIN AOS1
\paragraph{AOS1. The exact infinite-place pairings and self-dual densities.}
\label{claim:AOS1}
Use the supplied embedding $\varphi_v:E\hookrightarrow\mathbb C$
and supplied matrix $P_v$ with
$P_v^*\varphi_v(J)P_v=J_v$.  The tuple coordinate is
$W_v=P_v^{-1}\theta_v(Y)$, column by column, where $\theta_v$ is
the actual completion of $\varphi_v$ to $E\otimes_F F_\infty$.
The Hermitian coordinate is $S_v=\theta_v(S)$, with no conjugation by
$P_v$.  These coordinates are real linear homeomorphisms onto
$\prod_v\operatorname{Mat}_{n\times a}(\mathbb C)$ and
$\prod_v\operatorname{Herm}_a(\mathbb C)$, respectively.  They carry
the source pairings to
\[
 \begin{split}
 B(W,Z)&=2\sum_v\operatorname{Re}\operatorname{tr}(W_v^*J_vZ_v),\\
 B_H(S,T)&=\sum_v\operatorname{Re}\operatorname{tr}(S_vT_v).
 \end{split}
 \tag{AOS.6}
\]
Consequently \emph{every} allowed self-dual pair becomes
\[
             d\mu_\infty=2^{nad}\prod_v dW_v,
       \qquad d\nu_\infty=\prod_v dT_v,
 \tag{AOS.7}
\]
where $dW_v$ is coordinate real-imaginary Lebesgue measure and $dT_v$
is Euclidean volume for the Frobenius inner product on the real
Hermitian subspace.

\begin{proof}
The CM algebra at a real place is $\mathbb C$, with involution complex
conjugation.  Its trace to $\mathbb R$ is $z+\overline z$.
One can verify this without a choice of complex coordinates: the two
extensions of the real embedding to $E$ are $\varphi_v$ and
$\overline{\varphi_v}$, and their sum is the trace.  Base change
preserves the trace.  The completed embedding is bijective because a
quadratic CM basis gives two real linearly independent complex
numbers.  Applying it at every real place gives the asserted finite
product homeomorphism.  Multiplication by the supplied invertible
$P_v^{-1}$ then gives the tuple coordinates.

The pairing \texttt{BR} is the sum over columns of the quadratic
algebra trace of $x^*Jy$.  Substituting $x=P_vw$, $y=P_vz$ gives the
first equation of (AOS.6).  The definition of \texttt{pair\allowbreak Herm} has a
factor $1/2$ in front of the quadratic trace.  Since
$\operatorname{tr}(ST)$ is real for Hermitian $S,T$, this gives the
second equation of (AOS.6), with no factor two remaining.

In an orthonormal real basis of tuple space, the first form has
diagonal coefficients $2\epsilon_{v,i}$, each repeated for the real
and imaginary coordinate of each column.  Its absolute determinant
is $2^{2nad}$.  GMR2's Fourier-inversion calculation therefore gives
self-dual density $2^{nad}$.  On the Hermitian space the second form
is precisely the Frobenius inner product, so its determinant in an
orthonormal basis is one.  In particular its self-dual measure is
Frobenius volume.  If one instead uses unscaled off-diagonal real and
imaginary entries as coordinates, this latter volume has density
$2^{a(a-1)/2}$; that factor is already contained in our use of an
orthonormal Hermitian basis and must not be inserted a second time.

Existence follows from Fourier inversion and finite products.  GMR2
gives uniqueness even for the source's conditional Fourier-inversion
predicate, so (AOS.7) holds for every original witness.  This argument
also shows why the determinant of $P_v$ does not create another
factor: the measures are transported together with their pairings.
\end{proof}
% END AOS1

% BEGIN AOS2
\paragraph{AOS2. Two determinant identities for complex matrices.}
\label{claim:AOS2}
For $B\in\operatorname{GL}_a(\mathbb C)$, right multiplication by $B$
on $\operatorname{Mat}_{b\times a}(\mathbb C)$ has absolute real
determinant $|\det B|^{2b}$.  The real linear map
\[
       H\longmapsto B^{-*}HB^{-1}
       \quad\hbox{on }\operatorname{Herm}_a(\mathbb C)
\]
has absolute determinant $|\det B|^{-2a}$, with Frobenius volume on
both sides.  The assertions include zero dimensions.

\begin{proof}
The real determinant of a complex linear map is its complex
determinant times its conjugate.  Apply this to each of the $b$ rows
to obtain the first formula.  For the second, use a singular-value
decomposition $B=U\operatorname{diag}(s_1,\ldots,s_a)V$ with $s_i>0$.
Unitary conjugations are Frobenius isometries, so they have absolute
real determinant one.  On the diagonal Hermitian coordinate indexed
by $i$, conjugation by $\operatorname{diag}(s_i^{-1})$ multiplies by
$s_i^{-2}$.  On each two-dimensional real off-diagonal coordinate
indexed by $i<j$, it multiplies both coordinates by $(s_is_j)^{-1}$.
Their product is
\[
       \prod_i s_i^{-2}\prod_{i<j}(s_is_j)^{-2}
          =\prod_i s_i^{-2a}=|\det B|^{-2a}.
 \tag{AOS.8}
\]
When $a=0$ or $b=0$, the relevant maps are maps of the zero vector
space and the empty determinant and product are one, as required.
\end{proof}
% END AOS2

% BEGIN AOS3
\paragraph{AOS3. The full original \texttt{Place\allowbreak Jac} determinant.}
\label{claim:AOS3}
Let $J=\operatorname{diag}(\epsilon_1,\ldots,\epsilon_n)$ with signs
$\epsilon_i\in\{1,-1\}$.  Let $S$ be any set of $k$ coordinate
indices, put $n=a+k$, and let $w\in\operatorname{Mat}_{n\times a}
(\mathbb C)$ vanish in the rows indexed by $S$.  Suppose
$T=w^*Jw$ is invertible.  Let $\mathfrak m$ consist of the matrices
$Z$ satisfying $Z^*J+JZ=0$ whose $S\times S$ block is zero, with its
Frobenius metric.  Then
\[
 \mathcal D_w:\mathfrak m\oplus\operatorname{Herm}_a(\mathbb C)
       \longrightarrow\operatorname{Mat}_{n\times a}(\mathbb C),
 \qquad
 \mathcal D_w(Z,H)=Zw+\tfrac12wT^{-1}H
 \tag{AOS.9}
\]
is a real linear isomorphism, and
\[
        |\det_{\mathbb R}\mathcal D_w|
                =2^{-na}|\det T|^k.
 \tag{AOS.10}
\]
Equivalently, its pushforward of product Euclidean volume is
$2^{na}|\det T|^{-k}$ times tuple volume.  This is the full original
\texttt{Place\allowbreak Jac}, for arbitrary finite orthonormal coordinate
spaces on its two factors.  No positivity assumption on the signs
indexed by $S$ is needed for this determinant statement.

\begin{proof}
Permute coordinates so that $S$ comes first.  These permutations are
Euclidean isometries.  Write
\[
       J=\begin{pmatrix}D_0&0\\0&D_1\end{pmatrix},
       \qquad w=\begin{pmatrix}0\\B\end{pmatrix}.
\]
Here $D_0,D_1$ are diagonal sign matrices and $B$ is invertible,
because $T=B^*D_1B$ is invertible.  The defining Lie algebra equations
give every $Z\in\mathfrak m$ uniquely as
\[
 Z=\begin{pmatrix}
       0&A\\-D_1A^*D_0&iD_1Q
    \end{pmatrix},\qquad Q^*=Q.
 \tag{AOS.11}
\]
Since the sign matrices are Euclidean unitary,
$\|Z\|_F^2=2\|A\|_F^2+\|Q\|_F^2$.
Thus $(R,Q)\mapsto Z$ with $A=R/\sqrt2$ is a real isometry
from coordinate $k\times a$ complex matrices and Frobenius Hermitian
matrices onto $\mathfrak m$.

In these coordinates (AOS.9) has upper block
$RB/\sqrt2$ and lower block
\[
       D_1\left(iQB+\tfrac12B^{-*}H\right)
       =D_1\left(iQ+\tfrac12B^{-*}HB^{-1}\right)B.
 \tag{AOS.12}
\]
The upper-block determinant is $2^{-ka}|\det B|^{2k}$ by AOS2.
For the lower block, the Hermitian congruence in $H$ contributes
$|\det B|^{-2a}$; multiplication by $1/2$ on its real $a^2$
coordinates contributes $2^{-a^2}$.  The map
$(Q,H')\mapsto iQ+H'$ is an isometry from two orthogonal Hermitian
spaces onto all complex matrices.  Finally right multiplication by
$B$ contributes $|\det B|^{2a}$, and left multiplication by $D_1$
contributes one.  The lower-block determinant is therefore
$2^{-a^2}$.  Multiplying the blocks gives
$2^{-a(k+a)}|\det B|^{2k}$, which is (AOS.10) because
$|\det T|=|\det B|^2$.

Each map used is invertible, which also proves that $\mathcal D_w$
is an isomorphism.  Any coordinates $L,H$ allowed in the original
\texttt{Place\allowbreak Jac} are onto Frobenius isometries by its exact norm and
surjectivity hypotheses.  Composing with them changes no absolute
determinant.  The usual linear change-of-variables formula takes the
reciprocal of (AOS.10), exactly as in the original measure identity.
The calculation uses empty products when $a=0$ or $k=0$, so both
boundary cases are included.
\end{proof}
% END AOS3

% BEGIN AOS4
\paragraph{AOS4. A Gram chart and the exact density of its fiber.}
\label{claim:AOS4}
Keep $J,w,T$ as in AOS3.  Near $T$ there is a real analytic matrix
function $b(Q)$ such that
\[
 b(T)=1,\qquad b(Q)^*Tb(Q)=Q,
       \qquad Db_T(H)=\tfrac12T^{-1}H.
 \tag{AOS.13}
\]
For sufficiently small neighborhoods of $0$ and $T$, the map
\[
       \Psi(Z,Q)=\exp(Z)wb(Q)
 \tag{AOS.14}
\]
is a diffeomorphism onto an open subset of the independent tuple
locus, and $\gamma(\Psi(Z,Q))=Q$.  If tuple and Hermitian measures
have the self-dual normalizations of AOS1, the fiber measure in the
chart $Z\mapsto\exp(Z)w$ has a continuous positive density $j_T(Z)$
and
\[
                     j_T(0)=|\det T|^k.
 \tag{AOS.15}
\]
This holds for every source fiber-family witness.

\begin{proof}
For $Q=T+H$ sufficiently close to $T$, the convergent matrix power
series for the square root defines
$b(Q)=(1+T^{-1}H)^{1/2}$, using the branch equal to one at $H=0$.
The matrix $A=T^{-1}H$ is $T$-self-adjoint: $A^*T=TA$.
Its real-coefficient power series therefore satisfies
$b(Q)^*T=Tb(Q)$.  Since $b(Q)^2=1+A$, (AOS.13) follows;
its derivative is the first-order term of the series.  This argument
does not assume that $T$ is positive definite.

The matrices $Z\in\mathfrak m$ are in the Lie algebra of $U(J)$,
so $\exp(Z)^*J\exp(Z)=J$.  Consequently (AOS.14) has Gram matrix
$Q$.  Its derivative at $(0,T)$ is exactly (AOS.9).  AOS3 and the
real inverse function theorem give the asserted local diffeomorphism.
Shrink to a product neighborhood on which its Jacobian does not
vanish and its tuples are independent.

Write $dZ$ for Euclidean volume on $\mathfrak m$ and $dQ$ for
Frobenius Hermitian volume.  The change-of-variables formula gives
\[
       d\mu(\Psi(Z,Q))
         =2^{na}|\det_{\mathbb R}D\Psi_{(Z,Q)}|\,dZ\,dQ.
 \tag{AOS.16}
\]
In these coordinates the Gram map is projection to $Q$.  GMR4's
fiber construction therefore has density
$2^{na}|\det D\Psi_{(Z,Q)}|\,dZ$ on each fiber in this chart.
GMR5 identifies any permitted fiber family with that construction
at \emph{every} $Q$, using continuity of compactly supported fiber
integrals; an almost-everywhere disintegration is not substituted
for that conclusion.  Evaluating (AOS.16) at $(0,T)$ and applying
(AOS.10) proves (AOS.15).  Positivity and continuity follow from
the nonzero continuous Jacobian.
\end{proof}
% END AOS4

\subsection{The prescribed Haar density and the compact volume}

% BEGIN AOS5
\paragraph{AOS5. The exact Haar density and the full original \texttt{Gauge\allowbreak Dens}.}
\label{claim:AOS5}
For any sign matrix $J$, let $d_G$ be the left invariant Riemannian
volume on $U(J)$ whose inner product at the identity is Frobenius
inner product on
$\mathfrak g_J=\{Z:Z^*J+JZ=0\}$.  If $\mu$ is any Haar measure
on this group, its density relative to $d_G$ is the positive scalar
\[
       c_\mu=\lim_{\rho\downarrow0}
          \frac{\mu(\exp B_{\mathfrak g_J}(0,\rho))}
               {\operatorname{vol}_{\mathfrak g_J}
                    (B_{\mathfrak g_J}(0,\rho))}.
 \tag{AOS.17}
\]
For the source measure $\mu_G^{\rm old}$ this is precisely its
$\operatorname{gaugeU}$, so $\mu_G^{\rm old}=g\,d_G$.
In particular the full original $\operatorname{GaugeDens}(\epsilon)$
holds for every dimension and every sign function $\epsilon$.
The comparison measure $d_G$ is used
to calculate the density; it does not replace $\mu_G^{\rm old}$.

\begin{proof}
The Lie algebra appearing here is the source's \texttt{g\allowbreak Lie}.
Indeed differentiating $\exp(tZ)^*J\exp(tZ)=J$ at zero gives
$Z^*J+JZ=0$.  Conversely this equation implies that the derivative
of $\exp(tZ)^*J\exp(tZ)$ is zero, so its value is $J$ for every
$t$.  Taking the real span in the source definition changes nothing.
The source realification map is a Frobenius isometry.

For completeness, smooth local group coordinates need no additional
structure theorem here.  The Cayley map
\[
       Z\longmapsto (1+Z/2)(1-Z/2)^{-1}
\]
is a chart from a neighborhood of zero in this Lie algebra to a
neighborhood of the identity in $U(J)$.  The equation
$Z^*J+JZ=0$ verifies unitarity by multiplication; its inverse is
$g\mapsto2(g-1)(g+1)^{-1}$ wherever the latter is defined.
Both maps are smooth, and the first derivative at zero is the
identity.  Translate this chart to obtain the group manifold.

The left-translated Frobenius inner product defines a smooth positive
volume density.  Its associated measure is locally finite, positive
on nonempty open sets, and left invariant.  It is therefore Haar.
Uniqueness of Haar measure gives $\mu=c_\mu d_G$ for a finite
$c_\mu>0$.  On a sufficiently small Lie algebra ball, exponential
is a diffeomorphism, and the pullback of $d_G$ is
$J_{\exp}(Z)\,dZ$, where $J_{\exp}$ is continuous and
$J_{\exp}(0)=1$ because $D\exp_0$ is the identity.  The ratio in
(AOS.17) is the average of $c_\mu J_{\exp}$ on this ball.  Uniform
continuity near zero makes that average tend to $c_\mu$.

The original \texttt{exp\allowbreak Set} is this exponential image, and its
\texttt{ratio\allowbreak At} uses exactly the Lie algebra volume just described.
Thus the limit exists and is $c_\mu>0$.  The source defines
$\operatorname{gaugeU}$ by \texttt{lim\allowbreak Under}, which equals any
existing limit.  Consequently $g=c_\mu$, and both the convergence
and strict positivity clauses of \texttt{Gauge\allowbreak Dens} follow.  This
proves the normalization from the literal gauge definition.  In
dimension zero the exponential ball and the group are singletons;
their Lie volume is one and the same assertion holds directly.
\end{proof}
% END AOS5

% BEGIN AOS6
\paragraph{AOS6. The compact unitary group has the required exact volume.}
\label{claim:AOS6}
For $U(j)$ with the Frobenius normalization of AOS5,
\[
                  d_{U(j)}(U(j))
                      =\prod_{i=1}^{j}\frac{(2\pi)^i}{(i-1)!}.
 \tag{AOS.18}
\]
In particular this proves the original \texttt{Gauge\allowbreak Vol\allowbreak U}: under its
stated \texttt{Gauge\allowbreak Dens} antecedent, the source measure
$\operatorname{muGauge}$ has total mass $\operatorname{volU}(j)$.

\begin{proof}
The group $U(j)$ is closed and bounded in its matrix space, hence
compact.  The map $q_j:U(j)\to S^{2j-1}$ taking the last column is
onto, and the stabilizer of the last standard basis vector is the
actual block $U(j-1)$.  Local orthonormal completion gives smooth
sections.  At the identity its vertical Lie algebra is the upper
$\mathfrak u(j-1)$ block.  Its Frobenius orthogonal complement has
matrices
\[
       \begin{pmatrix}0&z\\-z^*&it\end{pmatrix},
       \qquad \|Z\|_F^2=2\|z\|^2+t^2.
 \tag{AOS.19}
\]
The last-column derivative is $(z,it)$, whose round-sphere tangent
norm is $\|z\|^2+t^2$.  Its absolute determinant on orthonormal
horizontal coordinates is $2^{-(j-1)}$, since $z$ has real
dimension $2j-2$ and $t$ has dimension one.

Disintegrate the invariant group volume over the last-column map.
The fiber volume is the unchanged Frobenius volume on $U(j-1)$:
left multiplication is an isometry and the block inclusion preserves
Frobenius inner products.  The quotient density is invariant under
the transitive $U(j)$ action on the sphere.  Formula (AOS.19), or
ordinary change of variables in a local section times a fiber chart,
shows that this quotient volume is $2^{j-1}$ times round-sphere
area.  The equality of these densities at one point gives it
everywhere by invariance.  Thus
\[
 d_{U(j)}(U(j))
      =d_{U(j-1)}(U(j-1))\,2^{j-1}
                        \operatorname{area}(S^{2j-1}).
 \tag{AOS.20}
\]
The area is $2\pi^j/(j-1)!$: integrate $e^{-\|x\|^2}$ on
$\mathbb R^{2j}$, first as a product of $2j$ one-dimensional
Gaussian integrals, giving $\pi^j$, and then in polar coordinates.
The radial integral is
$\tfrac12\int_0^\infty e^{-s}s^{j-1}ds=(j-1)!/2$, by repeated
integration by parts.  Equation (AOS.20) therefore has successive
ratio $(2\pi)^j/(j-1)!$.

For $j=0$ the group and its Lie algebra have one point and the
normalized volume is one.  Induction proves (AOS.18), including
that boundary case.  Finally AOS5 gives
$\operatorname{muGauge}=g^{-1}\mu_G^{\rm old}=d_{U(j)}$,
so the source formula follows with its exact product convention.
\end{proof}
% END AOS6

% BEGIN AOS7
\paragraph{AOS7. The local density of the actual orbit pushforward.}
\label{claim:AOS7}
Suppose now that $J=\operatorname{diag}(I_k,D_1)$, and
$w=\binom{0}{B}$ with $B$ invertible.  Its pointwise stabilizer is
the actual block subgroup $H=U(k)$.  Let $\mu_G=g\,d_G$, with
$g>0$.  In the local orbit chart $Z\mapsto\exp(Z)w$ from AOS4,
$\pi_{w*}\mu_G$ has continuous positive density $q(Z)$, with
\[
                         q(0)=g\operatorname{volU}(k).
 \tag{AOS.21}
\]
The orbit map is proper, its pushforward is Radon, and every nonzero
invariant Radon measure on the same orbit is a positive multiple of
this pushforward.

\begin{proof}
An element fixing $w$ fixes the whole last $a$-dimensional coordinate
space, since $B$ is invertible.  Unitarity then makes it preserve
the orthogonal first $k$-space.  Hence it is exactly
$\operatorname{diag}(h,1)$, $h\in U(k)$; the converse is immediate.
The Frobenius splitting
$\mathfrak g_J=\mathfrak m\oplus\mathfrak u(k)$ is orthogonal.
Exponential and the inverse function theorem give a local section of
$G/H$ of the form $Z\mapsto\exp(Z)$.  AOS4 identifies this quotient
chart with the indicated orbit chart.

Pull back $\mu_G$ by $(Z,h)\mapsto\exp(Z)h$, using $dZ$ and
Frobenius Haar $d_H$ on the compact group $H$.  The resulting density
is a smooth positive function $D(Z,h)$.  At $(0,h)$, left translate
the derivative back by $h^{-1}$.  Its horizontal part is
$\operatorname{Ad}(h^{-1})Z$ and its vertical part is the usual
left-translated tangent vector in $\mathfrak h$.  Since $h$ is
Euclidean unitary and block diagonal, its adjoint action preserves
the Frobenius metric and the orthogonal splitting.  The derivative
therefore has absolute determinant one.  AOS5 implies
$D(0,h)=g$ for every $h\in H$.

This product map is a diffeomorphism onto the inverse image of a
sufficiently small quotient neighborhood: uniqueness of the local
section makes the $Z$ coordinate unique and then makes $h$ unique.
It follows that $q(Z)=\int_H D(Z,h)\,d_H(h)$.  Compactness gives
continuity and positivity, and AOS6 gives (AOS.21).

The same local-section description proves properness.  Cover a
compact subset of the orbit by finitely many smaller section
neighborhoods with compact closures inside section domains.  The
preimage of each corresponding compact piece lies in the product
of the compact section image with the compact $H$.  Their union is
compact.  Thus the orbit pushforward of the Radon Haar measure is
Radon.

For uniqueness, let $\rho$ be an invariant Radon measure on the
orbit and let $dh$ be probability Haar on $H$.  Integrating
$\int_H f(s(y)h)\,dh$ against $\rho(y)$ defines a left invariant
Radon measure on $G$, by the local-section argument of GMR11.
If $\rho\ne0$, invariance and transitivity make every nonempty open
set positive: otherwise its countably many translates would cover
the orbit with zero measure.  The lifted measure is therefore a
nonzero Haar measure, hence a finite positive multiple of $\mu_G$.
Pushing it back to the orbit recovers $\rho$, because $dh(H)=1$.
This proves the asserted uniqueness.  It also includes the zero
measure with scalar zero when that case is needed.
\end{proof}
% END AOS7

% BEGIN AOS8
\paragraph{AOS8. Exact one-place orbital measure, with no unknown scalar.}
\label{claim:AOS8}
Let $w\in\operatorname{Mat}_{n\times a}(\mathbb C)$ have invertible
Gram matrix $T=w^*Jw$, and suppose its orthogonal complement is
positive definite of dimension $k=n-a$.  Use the source character
$e^{2\pi it}$, its tuple and Hermitian self-dual measures, and any
permitted continuous fiber family $\omega$.  For the prescribed
Haar measure $\mu_G=g\,d_G$,
\[
            \pi_{w*}\mu_G
                 =g\operatorname{volU}(k)|\det T|^{-k}\,\omega_T.
 \tag{AOS.22}
\]

\begin{proof}
First suppose the orthogonal complement is a positive coordinate
$k$-space, so $w=\binom0B$.  Witt extension identifies its orbit
with the entire Gram fiber.  Since $T$ is invertible, that entire
fiber is inside the independent locus.  GMR4--GMR5 give an invariant
Radon fiber measure, positive on every nonempty relatively open set.
AOS7 therefore makes the orbit pushforward $c\omega_T$ for a finite
$c>0$.

Use the same local orbit coordinate $Z$ for both measures.  Its
fiber density is continuous with value $|\det T|^k$ by AOS4, and
its orbit-pushforward density is continuous with value
$g\operatorname{volU}(k)$ by AOS7.  Their proportionality forces
\[
       c=\frac{g\operatorname{volU}(k)}{|\det T|^k}.
\]
For example, integrate both densities over balls of radius $\rho$
and divide by their Euclidean volumes; as $\rho\downarrow0$ the
two averages tend to the stated positive values.  In a zero
dimensional chart the same calculation uses the singleton of
volume one.  This proves (AOS.22) without discarding a null set at
the base point.

For a general positive orthogonal complement, choose an orthonormal
basis of it.  Hermitian Gram--Schmidt on its nondegenerate
orthogonal complement extends this to an isometry from a coordinate
sign decomposition.  Equivalently, the elementary Witt extension
proved in GMR10 gives $c_0\in U(J)$ carrying a positive coordinate
$k$-space to this complement.  Thus $c_0^{-1}w$ has the preceding
block form, with the same Gram matrix $T$.

The unitary group is unimodular.  Here is the real version of
GMR14's trace-pairing argument.  The real Lie algebra
$\{Z:Z^*J+JZ=0\}$ is a real form of
$\operatorname{Mat}_n(\mathbb C)$: for the conjugate-linear
involution $\tau(Z)=-JZ^*J$, every complex matrix is the sum of
a fixed vector and $i$ times a fixed vector.  The pairing
$\operatorname{Re}\operatorname{tr}(AB)$ on this real form
complexifies to the nondegenerate matrix trace pairing.  It is
therefore nondegenerate and invariant under conjugation.  Thus
$(\det\operatorname{Ad}(g))^2=1$.  Conjugation multiplies a
left invariant volume density by the absolute determinant of its
derivative at the identity, which is consequently one.  Composing
conjugation with a left translation proves that every right
translation preserves every left Haar normalization.  Therefore
$\pi_w(g)=\pi_{c_0^{-1}w}(gc_0)$ gives identical pushforward
measures.  Applying the block calculation proves the general case.
The coordinate isometry is used only in this computation; the
measure on the original group and the actual orbit are unchanged.
\end{proof}
% END AOS8

\subsection{The source frame and all infinite-place witnesses}

% BEGIN AOS9
\paragraph{AOS9. The source frame has the prescribed determinant and compact stabilizer.}
\label{claim:AOS9}
For every original $D,X,u'$, the tuple (AOS.3) is independent,
has nondegenerate Hermitian Gram matrix, and at every real place
\[
          |\det\varphi_v(T_0)|=T_w(u').
 \tag{AOS.23}
\]
Its orthogonal complement is positive definite of dimension $k$.
The actual infinite pointwise stabilizer $H_\infty$ is a compact
group isomorphic to $\prod_{v\mid\infty}U(k)$.  In particular this
proves the original \texttt{Hinf\allowbreak Compact} at its full scope; neither
measure compatibility nor a test function is required for compactness.

\begin{proof}
The setting's frame identities give, in the order $(e,z,f)$,
\[
       T_0=\begin{pmatrix}
           0&0&I_r\\0&-T_w(u')&0\\I_r&0&0
          \end{pmatrix}.
 \tag{AOS.24}
\]
These identities are transported by the supplied coordinate map and
field identification, rather than by a newly chosen frame.
Reordering as $(e,f,z)$ shows
$\det T_0=(-1)^{r+1}T_w(u')$.  The integer $T_w(u')$ is positive,
so (AOS.23) follows at every real embedding.  Invertibility of the
Gram matrix implies independence: a relation between the columns
would be in the kernel of $T_0$.

Each hyperbolic pair has one positive and one negative direction,
as is seen from $(e_i+f_i)/\sqrt2$ and $(e_i-f_i)/\sqrt2$.
The vector $z$ is negative, of squared length $-T_w(u')$.
Thus the tuple span has signature $(r,r+1)$.  The original setting
has exactly $r+1$ negative directions.  Its rank is
$n=k+m-1=k+2r+1$, because its $m$ is even and at least two,
and $r=m/2-1$.  Sylvester's law of inertia, applied to the orthogonal
direct sum with this nondegenerate tuple span, makes its complement
positive definite and of dimension $k$.

An isometry fixing each tuple vector fixes the tuple span pointwise
and acts arbitrarily unitarily on its orthogonal complement.
Conversely every such complementary unitary extends by the identity
on the span.  These two matrix constructions are inverse continuous
homomorphisms.  At each real place the stabilizer is therefore
isomorphic to $U(k)$; the completed CM coordinates give the finite
product assertion for the actual $H_\infty$.  Compactness follows
from compactness of $U(k)$.  The argument includes $r=0$ (one
negative vector) and the zero-dimensional complementary group.
\end{proof}
% END AOS9

% BEGIN AOS10
\paragraph{AOS10. Products preserve the exact fiber measure at every parameter.}
\label{claim:AOS10}
Under the completed coordinates of AOS1, let $\omega_{v,T_v}$ be
the local fiber measures.  For every $T=(T_v)$,
\[
                         \omega_T=\bigotimes_v\omega_{v,T_v}.
 \tag{AOS.25}
\]
This equality holds for every permitted global self-dual pair and
every permitted global fiber-family witness.  At nondegenerate $T$,
the product of the orbit maps pushes the prescribed product Haar
measure to the product of the local orbit measures.

\begin{proof}
The independent locus over the product algebra is the product of the
local independent loci: multiplying a linear relation by each
coordinate idempotent proves this equivalence in both directions.
The Gram map is componentwise.  On products
of the local submersion charts of GMR3--GMR4, the ambient Jacobian is
the product of their positive Jacobians.  The fiber density at a
fixed parameter is therefore their product, by ordinary Fubini.
These product charts cover the independent locus and their measures
agree on overlaps, so they construct the product fiber family.
The support and disintegration statements follow by Tonelli and a
countable chart cover.  A compactly supported continuous test is
covered by finitely many product charts; a partition of unity and
a common compact bound in each chart give continuity of its fiber
integral by dominated convergence.  This verifies the full original
fiber-family condition, including singular parameters.

AOS1 identifies every permitted global self-dual measure with the
corresponding product.  GMR5 now identifies any permitted global
fiber family with the product construction pointwise, not merely
almost everywhere.  This proves (AOS.25).  For orbit pushforwards,
equality on measurable rectangles is the definition of product
measure and product pushforward; a monotone-class argument gives
equality of measures.  Compact stabilizers make these orbit
pushforwards Radon by AOS7, so the needed product measures are
sigma finite.  No infinite restricted product occurs in this
archimedean assertion.
\end{proof}
% END AOS10

% BEGIN AOS11
\paragraph{AOS11. The full original \texttt{Push\allowbreak Inf}.}
\label{claim:AOS11}
For every $D,X,h_b,F_{dN},F_{dN'},u,u',\psi$ satisfying the original
\texttt{Hyp\allowbreak X}, \texttt{Meas\allowbreak Compat\,$\psi$}, \texttt{Is\allowbreak Std\allowbreak Char}, and
all original \texttt{Gauge\allowbreak Dens} antecedents, for every witness
$h_T$ of Hermitian symmetry of $T_0$, and for every original
self-dual pair and fiber family,
\[
             (\pi_{Y_0})_*\mu_{G,\infty}
                  =\operatorname{ofReal}(K_X(u'))\,\omega_{T_0}.
 \tag{AOS.26}
\]
Here $T_0$ on the right is the literal \texttt{herm\allowbreak Of} determined by
the rational Gram matrix and $h_T$.  Thus (AOS.26) is the original
\texttt{Push\allowbreak Inf} equality, with its complete witness quantifiers.

\begin{proof}
The completed bridge intertwines the actual action with the product
of the real-model actions.  Explicitly,
$\theta_v(g)=P_vg_vP_v^{-1}$ and
$\theta_v(Y)=P_vW_v$, so
$P_v^{-1}\theta_v(gY)=g_vW_v$.  This uses the supplied bridge;
its equality on matrices determines the action without a further
isomorphism choice.  The first measure tie in (AOS.2) identifies
the prescribed $\mu_{G,\infty}$ with the pushforward of precisely
the product of the original real Haar measures.

AOS9 verifies the hypotheses of AOS8 at every real place.  AOS5
identifies the scalar of its original Haar measure with $g_v$.
Multiplying (AOS.22), using AOS10 and (AOS.23), gives the scalar
\[
       \prod_v g_v\operatorname{volU}(k)
                            |\det\varphi_v(T_0)|^{-k}
          =\left(\prod_vg_v\right)
               \operatorname{volU}(k)^dT_w(u')^{-kd}.
 \tag{AOS.27}
\]
It is exactly (AOS.1).  All its factors are strictly positive and
finite.  Transport back through the same coordinate homeomorphism
proves equality of the measures on the original tuple space.
The proof $h_T$ specifies an element of a fixed Hermitian subspace;
its underlying matrix is unchanged by a different symmetry proof.
Uniqueness in AOS1 and AOS10 proves the assertion for all the
self-dual and fiber witnesses quantified in the original statement.

No other part of the setting is altered.  In particular the actual
bridge, $\mu_{G,\infty}$, $\mu_{N,\infty}$ and
$\mu_{H,\infty}$ are the ones supplied before this argument.
The unused level and fundamental-domain conditions are retained
because they are antecedents of the literal original theorem.
\end{proof}
% END AOS11

% BEGIN AOS12
\paragraph{AOS12. A positive finite test for the full original \texttt{Const\allowbreak Test}.}
\label{claim:AOS12}
For every datum satisfying the original \texttt{Push\allowbreak Hyp}, there is a
measurable $\phi:V_\infty\to[0,\infty]$ such that
\[
 \begin{gathered}
             0<\int\phi\,d\omega_{T_0}<\infty,\\
 \int\phi\,d(\pi_{Y_0*}\mu_{G,\infty})
                  =K_X(u')\int\phi\,d\omega_{T_0}.
 \end{gathered}
 \tag{AOS.28}
\]
In fact $\phi$ may be taken continuous, nonnegative and compactly
supported, with support contained in the independent locus.

\begin{proof}
The independent locus is open and contains $Y_0$ by AOS9.  In any
fixed Euclidean coordinates, take $\delta>0$ such that the closed
ball of radius $2\delta$ about $Y_0$ is inside that locus.  Define
\[
             \phi(Y)=\max(1-\|Y-Y_0\|/\delta,0).
 \tag{AOS.29}
\]
Its support is compact and it is bounded by one.  The fiber measure
is Radon on the independent locus, so the integral is finite.
It is positive because $\phi\geq1/2$ on the relative fiber ball of
radius $\delta/2$, whose measure is positive by GMR4, or by the
positive local density in AOS4.  Regard (AOS.29) as an extended
nonnegative-valued function.  It is measurable and still has those
integral bounds.

The existential witnesses in \texttt{Push\allowbreak Hyp} supply the bridge,
sets and first word needed for AOS11; its other clauses are exactly
the remaining antecedents and witness conditions of that theorem.
Apply (AOS.26) to the just constructed test.  This proves (AOS.28)
with the prescribed constant and no cancellation of a zero or
infinite integral.  It is the full original \texttt{Const\allowbreak Test}.
\end{proof}
% END AOS12

\subsection{Existence, the actual compact Haar mass, and all sections}

% BEGIN AOS13
\paragraph{AOS13. The full original \texttt{Ex\allowbreak Inf}, with all witnesses.}
\label{claim:AOS13}
For every supplied setting, bridge, words, sets and character under
the antecedents of \texttt{Ex\allowbreak Inf}, the nested \texttt{All\allowbreak Some} assertion
in that definition holds: tuple self-dual measures exist; for every
such measure Hermitian self-dual measures exist; and for every such
pair a full continuous fiber family exists.

\begin{proof}
Apply GMR2 at the finite product of real places with the source
pairings and character.  AOS1 gives their exact coordinate forms
and hence nondegeneracy, so it also supplies a direct realization
of the two self-dual measures.  GMR4 constructs a fiber family on
the whole Hermitian parameter space for each prescribed pair of
additive Haar measures.  Thus it applies to every self-dual pair,
not only to one chosen in the proof.  Equivalently one may use the
product construction in AOS10.  It is valid on the independent
locus even at singular parameters.  This verifies the existence
and universal clauses in their original order.  The global measure
compatibility antecedent is retained; as in the original source,
it could also supply these existence clauses directly from its
\texttt{Is\allowbreak Om\allowbreak T} field.
\end{proof}
% END AOS13

% BEGIN AOS14
\paragraph{AOS14. The full original \texttt{Htot\allowbreak Inf} for the unchanged stabilizer measure.}
\label{claim:AOS14}
For every original datum under \texttt{Hyp\allowbreak X},
\[
         \mu_{H,\infty}(H_\infty)=\mu_H(u')(G_\infty),
             \qquad 0<\mu_H(u')(G_\infty)<\infty.
 \tag{AOS.30}
\]
The symbols on the right denote the original ambient real-model
measure and its mass on the whole ambient space, exactly as in
\texttt{Htot\allowbreak Inf}; it is supported on its stabilizer.

\begin{proof}
The actual $H_\infty$ is compact by AOS9.  The setting supplies
the Haar-measure property of its unchanged $\mu_{H,\infty}$.
A Haar measure is finite on compact sets and has positive mass
on a nonempty open set, so its total mass is finite and positive.
Now evaluate the second measure tie in (AOS.2) on the whole ambient
group.  Both pushforward maps are measurable, and the preimage of
the whole space is the whole domain.  The resulting equality is
the first equality of (AOS.30).  It transfers both inequalities to
the original $\mu_H(u')$.  This is the full original statement;
the computation does not normalize either stabilizer measure to
probability and does not replace it by the Frobenius comparison
measure used in AOS6--AOS8.
\end{proof}
% END AOS14

% BEGIN AOS15
\paragraph{AOS15. Compact-stabilizer unfolding for every section.}
\label{claim:AOS15}
Let a locally compact second countable group $G$ act on an orbit
$Y=G/H$ with continuous local sections and compact stabilizer $H$.
Let $\mu_G$ be a right invariant left Haar measure, let $\eta$ be
any Haar measure on $H$, and put $\rho=\pi_*\mu_G$.
For every nonnegative measurable $f$ and every section $s$ on the
orbit, even one not assumed measurable,
\[
       \int_G f\,d\mu_G
        =\eta(H)^{-1}
           \int_Y\int_H f(s(y)h)\,d\eta(h)\,d\rho(y).
 \tag{AOS.31}
\]
The same formula holds if $Y$ is regarded as a measurable orbit in
an ambient space and the measures are extended by zero.

\begin{proof}
Compactness and nonzero Haar mass give $0<\eta(H)<\infty$.
Choose a countable cover by continuous section domains and
disjointify it to obtain a Borel section $s_0$.
For this section the inner integral is measurable by Tonelli.
Since $\rho$ is a pushforward, the double integral equals
\[
      \int_G\int_H f(s_0(\pi(g))h)\,d\eta(h)\,d\mu_G(g).
\]
For every $g$, the two representatives $s_0(\pi(g))$ and $g$
differ by right multiplication by an element of $H$.
Left invariance of $\eta$ therefore makes the inner integral
$\int_H f(gh)\,d\eta(h)$.  Interchanging the two nonnegative
integrals and applying right invariance of $\mu_G$ gives
$\eta(H)\int_G f\,d\mu_G$.  Divide by its finite nonzero
factor to obtain (AOS.31), also when the integral of $f$ is infinite.

For any other section $s$, the same left change of variable in $H$
shows pointwise on the orbit that its inner integral equals the
one for $s_0$.  Thus it is the same measurable function there;
no measurable selection property of $s$ is needed.  In an ambient
space the two inner functions agree outside a set of outer measure
zero, so their lower integrals agree as well.  This verifies the
literal all-section form used in the source.
\end{proof}
% END AOS15

% BEGIN AOS16
\paragraph{AOS16. The full original \texttt{Weil\allowbreak Inf}.}
\label{claim:AOS16}
Under every original \texttt{Hyp\allowbreak X}, for all nonnegative measurable
$f$ on the actual $G_\infty$ and all source sections $s$,
\[
 \int_{G_\infty}f(g)\,d\mu_{G,\infty}(g)
   =\mu_{H,\infty}(H_\infty)^{-1}
        \int_{V_\infty}\int_{H_\infty}
                 f(s(y)h)\,d\mu_{H,\infty}(h)
                     \,d(\pi_{Y_0*}\mu_{G,\infty})(y).
 \tag{AOS.32}
\]
This assertion does not require a choice of self-dual measure or
fiber family.

\begin{proof}
The actual $G_\infty$ is a finite product of real unitary groups
through the supplied completed bridge.  It is locally compact,
second countable and unimodular, by the same trace-form argument
used in AOS8.  The stabilizer is the actual $H_\infty$ of AOS9,
with the supplied Haar measure from AOS14.  The orbit map has
continuous local sections: AOS4 gives them at the base tuple,
and translation gives them at every other point.  Its orbit is
the Gram fiber by Witt extension, hence closed because its Gram
matrix is nondegenerate.  The quotient and orbit topologies agree
in these charts.  All hypotheses of AOS15 now hold.  Its formula
is exactly (AOS.32), including the requirement on $s$ only where
an orbit representative exists.  This proves \texttt{Weil\allowbreak Inf}
without changing the original $\mu_{H,\infty}$.
\end{proof}
% END AOS16

% BEGIN AOS17
\paragraph{AOS17. The complete original archimedean integration clause.}
\label{claim:AOS17}
For every datum, character and Hermitian witness under exactly the
\texttt{Hyp\allowbreak X}, \texttt{Meas\allowbreak Compat\,$\psi$}, \texttt{Is\allowbreak Std\allowbreak Char}
and \texttt{Gauge\allowbreak Dens} antecedents of AOS11, put
\[
             a_\infty(X,u')
                 =\frac{K_X(u')}{\mu_H(u')(G_\infty)}.
 \tag{AOS.33}
\]
The denominator here is its finite real mass.  This number is
strictly positive, and the full original
$\texttt{Ainf\allowbreak Clause}(X,\psi,u',h_T,a_\infty)$ holds.  Equivalently,
for every permitted self-dual pair and fiber family, every
nonnegative measurable $f$, and every source section $s$,
\[
       \int f\,d\mu_{G,\infty}
          =a_\infty(X,u')
             \int_{V_\infty}\int_{H_\infty}
                    f(s(y)h)\,d\mu_{H,\infty}(h)\,d\omega_{T_0}(y).
 \tag{AOS.34}
\]

\begin{proof}
AOS13 supplies all existence clauses in the nested \texttt{All\allowbreak Some}
contract, in its original order.  Fix any witnesses quantified
there.  Substitute (AOS.26) into (AOS.32).  The nonnegative
integral against a scalar multiple of a measure extracts the
scalar $K_X(u')$.  By AOS14 its remaining denominator is exactly
the unchanged finite positive mass $\mu_H(u')(G_\infty)$.
This gives (AOS.34) and the number (AOS.33).  Positivity follows
from positivity of the gauge factors, (AOS.18), $T_w(u')>0$ and
(AOS.30).  For extended nonnegative integrals these operations
remain legitimate because every scalar being divided out is finite
and nonzero.  This proves \texttt{Ainf\allowbreak Value} at its full original
scope, including all fiber witnesses and all sections.
\end{proof}
% END AOS17

% BEGIN AOS18
\paragraph{AOS18. The full original \texttt{Aprime\allowbreak Inf}: uniqueness of the prescribed scalar.}
\label{claim:AOS18}
For every original datum and every $h_T$ under the antecedents of
AOS11, if $a'>0$ satisfies the original \texttt{Ainf\allowbreak Clause}, then
\[
       a'=\frac{\operatorname{gaugeU}(X.B.I)
                         \operatorname{volU}(k)^dT_w(u')^{-kd}}
                        {\mu_H(u')(G_\infty)}.
 \tag{AOS.35}
\]
This is the full original \texttt{Aprime\allowbreak Inf}, not a statement that
an unspecified positive scalar exists.

\begin{proof}
Choose a compact neighborhood $C$ of the identity in the actual
$G_\infty$.  It has positive finite Haar mass; it is Borel because
the group is Hausdorff.  This construction needs only the supplied
locally compact group and Haar measure.  In particular it gives
the full original \texttt{Ginf\allowbreak Test} under \texttt{Hyp\allowbreak X} alone,
without any character or gauge hypothesis.  Choose one self-dual
pair and one fiber family
whose existence is guaranteed by AOS13, and one Borel section.
Apply the clause for $a'$ and the clause proved in AOS17 to
$f=1_C$.  Both use \emph{the same} witnesses, section and original
stabilizer measure, and hence the same nonnegative inner integral
$I$.  They give
\[
              0<\mu_{G,\infty}(C)=a'I
                      =a_\infty(X,u')I<\infty.
\]
Since $a'$ is a finite positive real number, these equalities
force $0<I<\infty$.  Cancelling $I$ proves
$a'=a_\infty(X,u')$, which is (AOS.35).  The source
\texttt{Aprime\allowbreak Inf\allowbreak S} and \texttt{Aprime\allowbreak Inf} have this identical
formula and quantifiers, so this proves the literal endpoint.
\end{proof}
% END AOS18

\subsection{Use on the rescaled setting and the remaining consumers}

Apply the chapter to the actual output of GMR19.  Its one standard
character supplies (AOS.5), its preserved measure tie supplies
(AOS.2), and its compatible measure record supplies the original
\texttt{Meas\allowbreak Compat\,$\psi$}.  The original \texttt{Hyp\allowbreak X} and
\texttt{Gauge\allowbreak Dens} antecedents are kept in the statements of the
archimedean consumers; AOS5 supplies the latter for every sign
function.  Every later word is treated using this same
character, bridge and infinity measures.  The auxiliary Frobenius
measures and orthogonal coordinates used to calculate a scalar
have not become fields of a new setting.

The completed endpoint is the exact archimedean integration scalar,
including its all-witness and all-section clauses.  The original
\texttt{Real\allowbreak Bracket} still requires its unipotent and local
disintegration factors, and the full mass comparison still requires
the finite-place and arithmetic quotient-volume producers.  In
particular the original \texttt{Siegel\allowbreak Weil\allowbreak Unitary\allowbreak Convergent}
premise retained by GMR19 is not proved or eliminated from that
upstream original statement by this chapter.  The source handoff
records those next consumers without claiming a whole selected
RH/GRH obligation has been closed.


\section{The actual Wach coefficient embedding}
\label{sec:wach-inclusion-saturation-v054}

\status{This chapter proves all four clauses of the original
Wach-inclusion theorem: injectivity, reflection of units, saturation
for multiplication by the residue characteristic, and the exact
intersection with the integral Witt tensor ring. The supplied
coefficient series ring, period-ring maps and Frobenius are kept
throughout. This supplies one of the three period-ring inputs
retained in WLD26. The actual Wach prism and the invariant
Frobenius-matrix descent remain separate inputs.}

The coefficient embedding is not an arbitrary map between two fields.
Its source is a completed Laurent-series ring with several coefficient
components. Its target is an integral coefficient extension of a Witt
ring. Neither target is assumed to be a domain after tensoring.
The key observation is that a nonzero component idempotent cannot be
divisible by an element of the Jacobson radical. Together with the
first coefficient of a power series not divisible by a uniformizer,
this reflects divisibility from the Witt tensor ring back to the
given series ring.

We first establish the actual ring properties used in this argument.
The proofs use standard elementary foundations of finite extensions
of $\mathbb Q_p$, tensor products, and the coordinate formulas for
$p$-typical Witt vectors. No Wach module, period comparison functor,
or invariant Frobenius-solution descent is an input.

\subsection{The actual coefficient and Witt tensor rings}

% BEGIN WIS1
\paragraph{WIS1: complete integral coefficients and one uniformizer.}
\label{claim:WIS1}
Let $L/\mathbb Q_p$ be a finite extension and let $O$ be the integral
closure of $\mathbb Z_p$ in $L$. Then $O$ is a complete discrete
valuation ring, finite free of positive rank over $\mathbb Z_p$,
and its fraction field is $L$. There are a uniformizer $\varpi$,
an integer $e\geq1$, and $u\in O^\times$ such that
\[
                         p=u\varpi^e.                     \tag{WIS.1}
\]
Both the $p$-adic and $\varpi$-adic topologies on $O$ are complete
and separated.

\begin{proof}
Use the unique extension of the $p$-adic absolute value to $L$.
An element is integral over $\mathbb Z_p$ exactly when its absolute
value is at most one. Indeed, a monic integral equation precludes
absolute value greater than one by the ultrametric inequality.
Conversely, all conjugates have the same absolute value, so the
coefficients of the minimal polynomial have absolute value at most
one and belong to $\mathbb Z_p$. An integral element is a unit of
$O$ exactly when its absolute value is one: its inverse is integral
in that case, and the two integral factors of a product equal to
one must each have absolute value one. Thus the nonunits form
the ideal of elements with absolute value less than one, and $O$
is local.

For completeness we recall why the integral closure is finite.
Choose a $\mathbb Q_p$-basis $b_1,\ldots,b_d$ of $L$ consisting of
integral elements, after scaling an arbitrary basis. The trace
pairing is nondegenerate, since this characteristic-zero extension
is separable. For $x\in O$, each $\operatorname{Tr}_{L/\mathbb Q_p}
(xb_i)$ belongs to $\mathbb Z_p$: the trace is integral and lies
in $\mathbb Q_p$. Hence $O$ lies in the trace-dual lattice of
$\sum_i\mathbb Z_p b_i$. It is a submodule of a finite free
$\mathbb Z_p$-module and is therefore finite. It is torsion-free,
so the elementary structure theorem over the principal ideal
domain $\mathbb Z_p$ makes it finite free; its rank is $d>0$.
Scaling any element of $L$ by a sufficiently large power of $p$
puts it in $O$, proving the fraction-field assertion.

The finite integral closure is a one-dimensional integrally closed
Noetherian domain. It is local and not a field, because $p$ is
nonzero and has absolute value less than one. The discrete
valuation criterion for a one-dimensional normal Noetherian local
domain gives a uniformizer $\varpi$. Factoring the nonzero nonunit
$p$ gives (WIS.1). A finite free $\mathbb Z_p$-module is $p$-adically
complete and separated by taking limits in each basis coordinate.
The identity $pO=\varpi^eO$ makes the two stated filtrations cofinal,
which proves the remaining completeness assertions.
\end{proof}
% END WIS1

% BEGIN WIS2
\paragraph{WIS2: the Witt-coordinate properties used below.}
\label{claim:WIS2}
Let $S$ be a perfect ring of characteristic $p$. The ring $W(S)$
is $p$-adically complete and separated, multiplication by $p$ is
injective, and the constant-coordinate map has kernel $pW(S)$.
Moreover,
\[
       a\in W(S)^\times\quad\Longleftrightarrow\quad
                         a_0\in S^\times.                 \tag{WIS.2}
\]
For an inclusion of perfect characteristic-$p$ rings $S\hookrightarrow F$,
the induced map $W(S)\to W(F)$ is injective and its image is
$p$-saturated.

\begin{proof}
In Witt coordinates the multiplication formula is
\[
          (pa)_0=0,\qquad (pa)_{r+1}=a_r^p.                 \tag{WIS.3}
\]
Frobenius on $S$ is bijective. Thus a vector is divisible by $p$
exactly when its zeroth coordinate vanishes, and division by $p$
is unique. Iteration says that divisibility by $p^N$ is equivalent
to vanishing of the first $N$ coordinates. In particular,
multiplication by $p$ is injective and the $p$-adic filtration is
separated.

A compatible system modulo $p^N$ determines each of its coordinates
and hence one Witt vector. Addition and multiplication on every
finite truncation are given by the Witt polynomials; the vector
so obtained has the prescribed truncations. This proves
completeness. In particular $pW(S)$ lies in the Jacobson radical:
for any $b$, the series $\sum_{r\geq0}(pb)^r$ is an inverse of
$1-pb$.

Units map to units under the zeroth-coordinate homomorphism.
Conversely, if $a_0$ is a unit, multiply $a$ by the inverse of
the Teichmuller vector $[a_0]$. The result is $1+pb$ for some
$b$, and the preceding geometric series makes it a unit.
This proves (WIS.2).

An inclusion of coefficient rings gives an inclusion coordinate
by coordinate. If the image of $a\in W(S)$ is $pw$ in $W(F)$,
then (WIS.3) gives $w_r^p=a_{r+1}$ in $F$ for every $r$.
Let $b_r$ be the unique $p$th root of $a_{r+1}$ in $S$.
Injectivity of Frobenius in $F$ gives $w_r=b_r$, so
$w$ is the image of the Witt vector $(b_r)_{r\geq0}$ in $W(S)$.
This is the asserted saturation, including the actual divided
element rather than only its existence in a quotient.
\end{proof}
% END WIS2

% BEGIN WIS3
\paragraph{WIS3: the cyclotomic element is in the maximal ideal of $A_{\inf}$.}
\label{claim:WIS3}
Use the original tilt ring
\[
 S=\mathcal O_C^\flat
   =\varprojlim_{x\mapsto x^p}(\mathcal O_C/p),\qquad
 A_{\inf}=W(S).
\]
Then $S$ and $A_{\inf}$ are local. For the given compatible
cyclotomic system $\epsilon$ and its $p^s$th root $\epsilon_s$,
the element
\[
                         \mu_s=[\epsilon_s]-1             \tag{WIS.4}
\]
belongs to the maximal ideal of $A_{\inf}$ for every $s\geq0$.

\begin{proof}
The ring $\mathcal O_C$ is a valuation ring and $p$ is a nonunit,
so $\mathcal O_C/p$ is a nonzero local ring. For a compatible
sequence $x=(x_r)$ in its perfection, one has
$x_r^{p^r}=x_0$. If $x_0$ is a unit, every $x_r$ is a unit,
and the inverses form a compatible sequence. Thus $x$ is a unit
exactly when $x_0$ is a unit. Its nonunits form the inverse
image of the maximal ideal under the zeroth projection, proving
locality of $S$. WIS2 proves locality of $W(S)$ in the same way.

The original cyclotomic data give
$\operatorname{pr}_0(\epsilon)=\zeta_0\bmod p=1$.
Hence $\epsilon-1$ is not a unit.
In characteristic $p$,
\[
                     (\epsilon_s-1)^{p^s}=\epsilon-1.
\]
A power of a unit would be a unit, so $\epsilon_s-1$ is not
a unit. It is the zeroth Witt coordinate of $\mu_s$.
By (WIS.2), $\mu_s$ is a nonunit of the local ring $A_{\inf}$,
as required. This also covers $s=0$; no nonzero-coordinate
assumption is inserted.
\end{proof}
% END WIS3

% BEGIN WIS4
\paragraph{WIS4: finite coefficient coordinates commute with the actual maps.}
\label{claim:WIS4}
Fix a $\mathbb Z_p$-basis $b_1,\ldots,b_d$ of $O$. For every
$\mathbb Z_p$-algebra $D$, each element of $O\otimes_{\mathbb Z_p}D$
has a unique expression
\[
                          \sum_{i=1}^d b_i\otimes d_i.      \tag{WIS.5}
\]
Coefficient maps in these coordinates commute with maps $D\to D'$
and with multiplication by $1\otimes z$. In particular,
\[
 x\in(1\otimes z)^N(O\otimes D)
       \quad\Longleftrightarrow\quad d_i\in z^ND
                      \text{ for every }i.               \tag{WIS.6}
\]

\begin{proof}
Tensor the basis isomorphism $O\simeq\mathbb Z_p^d$ with $D$.
Its inverse and coordinate maps send $o\otimes a$ to the
coordinates of $o$ multiplied by $a$. These formulas prove
the two compatibility assertions first on pure tensors and
then by additivity. If $x=(1\otimes z)^Ny$, its coordinates
are $z^N$ times those of $y$. Conversely, choose a quotient
for each coordinate and use their finite sum (WIS.5).
This proves (WIS.6) without assuming that $D$ is flat.
\end{proof}
% END WIS4

% BEGIN WIS5
\paragraph{WIS5: the actual integral inclusion is injective and saturated.}
\label{claim:WIS5}
Let $F=\operatorname{Frac}(S)$ be the original tilt field and set
\[
 B=O\otimes_{\mathbb Z_p}A_{\inf},\qquad
 W=O\otimes_{\mathbb Z_p}W(F),\qquad c:B\longrightarrow W.
                                                               \tag{WIS.7}
\]
Both rings are $p$-adically complete and separated, multiplication
by $p$ is injective on them, $W\ne0$, and $c$ is injective.
Its image is $p$-saturated:
\[
                  c(b)=pw\ \Longrightarrow\
                        w=c(b')\text{ for some }b'\in B.   \tag{WIS.8}
\]

\begin{proof}
The supplied tilt construction makes $S$ a perfect integral
domain. Its fraction field is perfect as well, by taking
$p$th roots in the numerator and denominator. Apply WIS2
to $S$ and $F$, and then WIS4 to the fixed integral basis
of $O$. Completeness, separatedness, absence of $p$-torsion,
and injectivity of $c$ follow in each coordinate.
The underlying $W(F)$-module of $W$ is $W(F)^d$ with
$d>0$, so it is nonzero. For (WIS.8), each coordinate
of $c(b)=pw$ lies in the $p$-saturated image of
$A_{\inf}\to W(F)$ by WIS2. The finitely many divided
coordinates assemble to $b'$ using (WIS.5).
\end{proof}
% END WIS5

% BEGIN WIS6
\paragraph{WIS6: a finite matrix proves the required Jacobson-radical property.}
\label{claim:WIS6}
If $D$ is a commutative ring, $z\in\operatorname{Jac}(D)$,
and $C$ is a finite free commutative $D$-algebra, then
the image of $z$ lies in $\operatorname{Jac}(C)$. In particular,
for the rings (WIS.7),
\[
             \pi_B=1\otimes\mu_s\in\operatorname{Jac}(B),
             \qquad p\in\operatorname{Jac}(W).             \tag{WIS.9}
\]

\begin{proof}
For $b\in C$, multiplication by $1-zb$ on a fixed $D$-basis
has matrix $1-zM_b$. Its determinant is $1+zq$ for some
$q\in D$, and is a unit because $z$ belongs to the Jacobson
radical. The adjugate matrix gives an inverse linear map.
Its value at $1$ is an element $v\in C$ satisfying
$(1-zb)v=1$. Thus $1-zb$ is a unit for every $b$, which
is exactly the Jacobson-radical criterion.

The $A_{\inf}$-algebra $B$ is finite free by WIS4, and
$\mu_s\in\operatorname{Jac}(A_{\inf})$ by WIS3, proving
the first assertion in (WIS.9). The second follows directly
from $p$-adic completeness of $W$: the geometric series
in $pw$ inverts $1-pw$ for every $w\in W$.

Only (WIS.9) is needed in the embedding proof. This argument
does not assert that $B$ is a domain or local, and does not
assert the stronger $\mu_s$-adic completeness theorem.
\end{proof}
% END WIS6

\subsection{Power-series components and the abstract embedding problem}

% BEGIN WIS7
\paragraph{WIS7: the first coefficient not divisible by a uniformizer.}
\label{claim:WIS7}
Let $O$ be a local ring in which every nonunit is divisible by
$\varpi$. For each $a\in O[[T]]$, either $a=\varpi b$,
or there are $b,u\in O[[T]]$ and $k\geq0$ with
\[
                    u\in O[[T]]^\times,\qquad
                    a=\varpi b+T^ku.                      \tag{WIS.10}
\]
For $A=\prod_{i\in J}O[[T]]$, put
$t=(\varpi)_i$, $\pi=(T)_i$, and let $e_i$ be the component
idempotents. For every $a\in A$ and $i\in J$, either
$e_i a=tb$, or
\[
                     e_i a=tb+\pi^kue_i,\qquad u\in A^\times.
                                                               \tag{WIS.11}
\]

\begin{proof}
If every coefficient is divisible by $\varpi$, divide
the coefficients individually. Otherwise let $k$ be the
smallest index of a coefficient not divisible by $\varpi$.
Divide the coefficients below $k$ by $\varpi$ to form a
polynomial $b$, and take
$u=\sum_{r\geq0}a_{k+r}T^r$. Its constant coefficient
is a unit. The usual coefficient recursion constructs an
inverse power series for $u$, proving (WIS.10).

Apply this construction to $a_i$. Place $b$ in its $i$th
component and zero elsewhere. For the second alternative,
take the same unit series $u$ in every component; its product
with $e_i$ has only the needed component. This proves
(WIS.11), with a unit of the entire product ring.
\end{proof}
% END WIS7

% BEGIN WIS8
\paragraph{WIS8: coefficientwise completeness of the integral series ring.}
\label{claim:WIS8}
For the complete ring $O$ of WIS1, the ring
$A=\prod_{i\in J}O[[T]]$ is $p$-adically complete and separated.
In particular, if $a_k\in A$ and
\[
           a_k-a_{k'}\in p^kA\qquad(k\leq k'),
                                                               \tag{WIS.12}
\]
there exists $a\in A$ such that $a_k-a\in p^kA$ for every $k$.

\begin{proof}
For every component $i$ and degree $r$, take the $p$-adic
limit in $O$ of the coefficient of $T^r$ in $a_k(i)$.
These limits define a power series in each component.
Reducing modulo $p^k$ shows that the difference from $a_k$
has every coefficient divisible by $p^k$. Choosing those
quotients coefficientwise proves divisibility in $A$ itself.
If an element lies in every $p^kA$, each coefficient lies
in $\bigcap_kp^kO=0$, giving separatedness. The argument
even works for an arbitrary index set $J$; later its
finiteness is used to sum the component idempotents.
\end{proof}
% END WIS8

% BEGIN WIS9
\paragraph{WIS9: Frobenius keeps every component visible.}
\label{claim:WIS9}
Let $A=\prod_{i\in\mathbb Z/n\mathbb Z}O[[T]]$ with $n>0$,
let $j:A\to E$ and $f:E\to W$ be unital ring maps with
$W\ne0$, and suppose operators $\phi_A,\phi_E,\phi_W$ satisfy
\[
 \phi_A(e_i)=e_{i+1},\qquad
 \phi_Ej=j\phi_A,\qquad f\phi_E=\phi_Wf,\qquad \phi_W(0)=0.
                                                               \tag{WIS.13}
\]
Then $f(j(e_i))\ne0$ for every $i$.

\begin{proof}
If one image were zero, iterating (WIS.13) would make
the images of $e_i,e_{i+1},e_{i+2},\ldots$ all zero.
These run through all components. But the idempotents sum
to $1$, whose image is $1\ne0$ in $W$.
For the original series Frobenius,
$\phi_A(a)_i=\phi_0(a_{i-1})$, the first identity in
(WIS.13) follows because $\phi_0$ preserves zero and one.
No injectivity or surjectivity of a Frobenius operator
is needed for this argument.
\end{proof}
% END WIS9

We now formulate the algebraic problem in a way that keeps
the four original maps distinct. Let
\[
 A\xrightarrow{j}E\xrightarrow{f}W,\qquad
 A\xrightarrow{g}B\xrightarrow{c}W,\qquad fj=cg              \tag{WIS.14}
\]
be a commutative square of unital maps. In the next eight claims
we use precisely the following data.

\begin{enumerate}
\item In $A$ there are $t,\pi,u$ with $u$ a unit and
      $p=ut^e$ for an integer $e\geq1$.
\item A finite family of idempotents $e_i$ sums to one,
      every $f(j(e_i))$ is nonzero, and the dichotomy
      (WIS.11) holds for every $a$ and $i$.
\item $A$ has the completeness property (WIS.12).
      The ring $E$ is $p$-adically complete and separated,
      and $j(\pi)$ is a unit.
\item For every $x\in E$ there are $a\in A$, $m\geq0$ such that
      $xj(\pi)^m-j(a)\in pE$. Also, for every $a\in A$ and $k$,
      $j(a)\in p^kE$ implies $a\in p^kA$.
\item In $W$, multiplication by $p$ is injective and
      $p\in\operatorname{Jac}(W)$. The image of $c$ is
      $p$-saturated, in the sense of (WIS.8), and
      $g(\pi)\in\operatorname{Jac}(B)$.
\end{enumerate}

Write $t_E=j(t)$, $\pi_E=j(\pi)$, and
$t_W=f(t_E)$, $\pi_W=f(\pi_E)$; a subscript on another
element of $A$ denotes its image by the same maps.
At this stage neither $f$ nor $g$ is assumed injective.
For a subring image, notation such as $c(B)$ means its
set-theoretic image.

% BEGIN WIS10
\paragraph{WIS10: the uniformizer is radical and regular where needed.}
\label{claim:WIS10}
Under (WIS.14) and the listed hypotheses,
\[
 t_E\in\operatorname{Jac}(E),\qquad
 t_W\in\operatorname{Jac}(W),\qquad
 t_Ww=0\Longrightarrow w=0.                                \tag{WIS.15}
\]
No nonzero idempotent $f(j(e_i))$ is divisible by $t_W$.
Every $x\in E$ has a representation
\[
                     x\pi_E^m=j(a)+t_Ez                   \tag{WIS.16}
\]
for some $a\in A$, $z\in E$, $m\geq0$.

\begin{proof}
Completeness puts $p$ in $\operatorname{Jac}(E)$ by the
geometric series argument. Since $p=u_Et_E^e$ and $u_E$
is a unit, $t_E^e$ belongs to every maximal ideal.
Each maximal ideal is prime, hence contains $t_E$.
The same proof in $W$ uses the supplied radical property
of $p$. If $t_Ww=0$, multiplying by $u_Wt_W^{e-1}$
gives $pw=0$, so $w=0$.

An idempotent $h$ in the Jacobson radical is zero:
$1-h$ is a unit and $h(1-h)=0$. Thus if $t_W$ divided
$f(j(e_i))$, that nonzero idempotent would be zero.
Finally, the supplied density condition writes
$x\pi_E^m-j(a)=pr=u_Et_E^er$.
Take $z=u_Et_E^{e-1}r$ to obtain (WIS.16).
\end{proof}
% END WIS10

% BEGIN WIS11
\paragraph{WIS11: reflection of one uniformizer factor.}
\label{claim:WIS11}
For every $x\in E$,
\[
                         t_W\mid f(x)
                    \quad\Longrightarrow\quad t_E\mid x. \tag{WIS.17}
\]

\begin{proof}
Choose (WIS.16). Its image in $W$ shows that $t_W$
divides $f(j(a))$. Fix a component $i$. If the second
alternative of (WIS.11) held, mapping it to $W$ would give
\[
       \pi_W^kv_W f(j(e_i))
          =f(j(e_i))f(j(a))-t_Wf(j(b)).
\]
The right side is divisible by $t_W$, while
$\pi_W^kv_W$ is a unit. This would make the nonzero
idempotent $f(j(e_i))$ divisible by $t_W$, contrary
to WIS10. Thus $e_i a$ is divisible by $t$ for every $i$.
Summing over the finite family gives $a=tb$. Equation
(WIS.16) becomes $x\pi_E^m=t_E(j(b)+z)$.
Cancel the unit $\pi_E^m$ to obtain (WIS.17).
\end{proof}
% END WIS11

% BEGIN WIS12
\paragraph{WIS12: reflection of every uniformizer power.}
\label{claim:WIS12}
For every $k\geq0$ and $x\in E$,
\[
                      t_W^k\mid f(x)
                 \quad\Longrightarrow\quad t_E^k\mid x.   \tag{WIS.18}
\]

\begin{proof}
The case $k=0$ is immediate. If
$f(x)=t_W^{k+1}y$, WIS11 gives $x=t_Ex'$.
Therefore
$t_W(f(x')-t_W^ky)=0$. By the regularity in WIS10,
$f(x')=t_W^ky$. Apply the induction hypothesis to $x'$
and multiply its resulting divisibility by $t_E$.
The cancellation takes place in $W$ using its proved
regularity; it does not presume that $f$ is injective.
\end{proof}
% END WIS12

% BEGIN WIS13
\paragraph{WIS13: injectivity and all $p$-power saturations.}
\label{claim:WIS13}
The map $f$ is injective. For all $r\geq0$ and $x\in E$,
\[
 f(x)\in p^rW\quad\Longrightarrow\quad x\in p^rE.
                                                               \tag{WIS.19}
\]
If $f(x)=p^rw$, there is $x'\in E$ with
$x=p^rx'$ and $f(x')=w$.

\begin{proof}
If $f(x)=0$, WIS12 makes $x$ divisible by every
$t_E^{er}$. Since $p^r=u_E^rt_E^{er}$ with $u_E$
a unit, $x$ lies in every $p^rE$. Separatedness gives
$x=0$, proving injectivity.

For (WIS.19), the equality $p^r=u_W^rt_W^{er}$
shows that $t_W^{er}$ divides $f(x)$. Apply WIS12
and the same equality in $E$. Choose $x'$ with
$x=p^rx'$. The relation $p^r(f(x')-w)=0$ and
the absence of $p$-torsion in $W$ give $f(x')=w$.
This includes $r=0$.
\end{proof}
% END WIS13

% BEGIN WIS14
\paragraph{WIS14: reflection of units.}
\label{claim:WIS14}
If $f(x)$ is a unit of $W$, then $x$ is a unit of $E$.

\begin{proof}
Suppose $x$ is not a unit, and choose a maximal ideal
$M$ of $E$ containing $x$. By WIS10, $t_E\in M$.
Choose (WIS.16); it gives $j(a)\in M$.

For any $i$, the first alternative $e_i a=tb$ in
(WIS.11) is impossible. Indeed, multiplying (WIS.16)
by $j(e_i)$ and mapping to $W$ would give
\[
 f(j(e_i)) f(x)\pi_W^m
       =t_W\bigl(f(j(b))+f(j(e_i))f(z)\bigr).
 \tag{WIS.20}
\]
The factor $f(x)\pi_W^m$ is a unit, contradicting WIS10.
Hence the second alternative holds for every $i$.
Its image in $E$ is
\[
         \pi_E^kv_Ej(e_i)=j(e_i)j(a)-t_Ej(b)\in M.
\]
Canceling the unit on the left gives $j(e_i)\in M$
for every $i$. Their sum is one, contradicting that
$M$ is proper. Therefore $x$ is a unit.
\end{proof}
% END WIS14

\subsection{The integral intersection}

% BEGIN WIS15
\paragraph{WIS15: removing one integral uniformizer digit.}
\label{claim:WIS15}
For $x\in E$ with $f(x)\in c(B)$, there are $a\in A$
and $x'\in E$ such that
\[
                   x=j(a)+t_Ex',\qquad f(x')\in c(B).      \tag{WIS.21}
\]

\begin{proof}
Write $f(x)=c(y)$ and choose (WIS.16).
We first prove, for every $i$, an expression
\[
                         e_i a=tb_i+\pi^md_i.             \tag{WIS.22}
\]
The first alternative in (WIS.11) gives this with $d_i=0$.
In the second alternative write
$e_i a=tb_i+\pi^{k_i}v_i e_i$ with $v_i$ a unit.
If $k_i\geq m$, take $d_i=\pi^{k_i-m}v_i e_i$.

Suppose instead that $k_i<m$, and put $h_i=f(j(e_i))$.
Multiplying (WIS.16) by the component idempotent and
using $f(x)=c(y)$ gives
\[
 \pi_W^{k_i}(v_i)_W h_i
    \left(1-c(y)\pi_W^{m-k_i}(v_i)_W^{-1}\right)
        =-t_W\bigl(f(j(b_i))+h_i f(z)\bigr).                \tag{WIS.23}
\]
The factor in parentheses is
\[
          c\left(1-y\,g(\pi)^{m-k_i}g(v_i)^{-1}\right).
\]
Because $m-k_i\geq1$ and $g(\pi)\in\operatorname{Jac}(B)$,
this factor is a unit. The two other factors preceding
$h_i$ in (WIS.23) are units as well. Consequently $t_W$
would divide $h_i$, contradicting WIS10. This proves
(WIS.22) in every case.

Summing (WIS.22) gives $a=tb+\pi^md$.
Substitution into (WIS.16), followed by cancellation
of $\pi_E^m$, gives
\[
         x=j(d)+t_Ex',\qquad
         x'=(j(b)+z)\pi_E^{-m}.                            \tag{WIS.24}
\]
Its image says $c(y-g(d))=t_W f(x')$. Multiply by
$u_Wt_W^{e-1}$ and use $p=u_Wt_W^e$:
\[
       c\bigl(g(u)g(t)^{e-1}(y-g(d))\bigr)=p f(x').
\]
The $p$-saturation of $c(B)$ now proves $f(x')\in c(B)$.
This is the second assertion in (WIS.21), which is
essential for iterating the construction.
\end{proof}
% END WIS15

% BEGIN WIS16
\paragraph{WIS16: the exact integral intersection.}
\label{claim:WIS16}
Under the abstract hypotheses,
\[
             f(x)\in c(B)\quad\Longrightarrow\quad
                          x=j(a)\text{ for some }a\in A. \tag{WIS.25}
\]
Consequently
\[
                       f(E)\cap c(B)=f(j(A))=c(g(A)).     \tag{WIS.26}
\]

\begin{proof}
Iterate WIS15 finitely many times. At every integer $N$
it gives an $a^{(N)}\in A$ such that
$x-j(a^{(N)})\in t_E^NE$. Take $N=ek$ and use
$p^k=u_E^kt_E^{ek}$ to choose $a_k$ with
\[
                         x-j(a_k)\in p^kE.               \tag{WIS.27}
\]
For $k\leq k'$, subtraction shows
$j(a_k-a_{k'})\in p^kE$. The stated saturation of
$j:A\to E$ implies $a_k-a_{k'}\in p^kA$.
Completeness (WIS.12) supplies $a\in A$ with
$a_k-a\in p^kA$ for all $k$.

Maps of rings preserve divisibility by the integer $p^k$,
so $j(a_k)-j(a)\in p^kE$. Adding this to (WIS.27)
puts $x-j(a)$ in every $p^kE$. Separatedness of $E$
gives $x=j(a)$.
No continuity of $j$ for an additional topology is
required: the relevant congruences are literal ring
divisibilities. The reverse containment in (WIS.26)
is the commutative-square identity $fj=cg$.
\end{proof}
% END WIS16

% BEGIN WIS17
\paragraph{WIS17: the full abstract four-clause embedding theorem.}
\label{claim:WIS17}
For the square (WIS.14) with the hypotheses printed after it,
$f$ is injective, reflects units, satisfies
\[
 f(x)=py\quad\Longrightarrow\quad x=px'
                                  \text{ for some }x'\in E,
\]
and satisfies the intersection assertion (WIS.25).
These are exactly the four conclusions of the original
abstract theorem in the Wach-inclusion proof.

\begin{proof}
Injectivity and the displayed saturation are WIS13,
unit reflection is WIS14, and the last assertion is WIS16.
Their proofs use the listed finite family of idempotents,
its dichotomy and nonvanishing, not any additional
orthogonality, basis, or domain hypothesis on $E,B,W$.
The completeness assumption on $A$ is precisely the
existence of limits for (WIS.12); the proof did not
assume its separatedness.
\end{proof}
% END WIS17

% BEGIN WIS18
\paragraph{WIS18: the full original Wach-inclusion theorem.}
\label{claim:WIS18}
Fix every original allowed tuple
$(p,P,AI,K,s,L)$, its condition $\mathrm{Good}(K,s,L)$,
and the supplied records $\Gamma,R,I$ with $\mathrm{IoW}(I)$.
Set $O=\mathcal O_L$, $n=[K_0:\mathbb Q_p]$,
$A=\prod_{\mathbb Z/n\mathbb Z}O[[T]]$, $E=R.EE$,
and let $B,W,c$ be the actual rings and map (WIS.7).
With the original maps
\[
 j:A\longrightarrow E,\qquad g=I.io:A\longrightarrow B,
                   \qquad f=I.ioE:E\longrightarrow W,
\]
all four original conclusions hold:
\begin{enumerate}
\item $f$ is injective.
\item If $f(x)$ is a unit, then $x$ is a unit.
\item If $f(x)=py$ in $W$, then $x=px'$ for some $x'\in E$.
\item If $f(x)=c(y)$ for some $y\in B$, then $x=j(a)$
      for some $a\in A$.
\end{enumerate}

\begin{proof}
We check every input of WIS17 on the actual data.
By WIS1 choose $\varpi,u,e$ with $p=u\varpi^e$.
In $A$ take the constant tuples $t=(\varpi)_i$ and
$u_A=(u)_i$, and $\pi=(T)_i$.
The original period datum identifies its residue
characteristic with $p$; thus the completeness,
density, saturation and invertibility clauses of
$R$ are precisely
\[
 \begin{gathered}
 E\text{ is }p\text{-adically complete and separated},
 \qquad j(\pi)\in E^\times,\\
 \forall x\in E,\ \exists a,m:\ xj(\pi)^m-j(a)\in pE,\\
 j(a)\in p^kE\Longrightarrow a\in p^kA .
 \end{gathered}                                            \tag{WIS.28}
\]
The density clause here is its original clause at
precision one; no stronger description of $E$ is used.
The original $I.ioE_{\mathrm{ext}}$ is $fj=cg$, and
$I.io_\pi$ identifies $g(\pi)$ with $1\otimes\mu_s$.

The component idempotents sum to one and WIS7 gives
their dichotomy. WIS8 gives completeness of $A$.
The integer $n$ is positive because $K_0$ is a finite
field extension of $\mathbb Q_p$. Its component
Frobenius shifts the indices, and the two original
Frobenius-compatibility clauses give (WIS.13).
WIS5 makes $W$ nonzero, so WIS9 gives nonvanishing
of every component image.

Finally WIS5 supplies absence of $p$-torsion and
$p$-saturation of $c(B)$, and WIS6 supplies both
required Jacobson-radical assertions. WIS17 therefore
applies to the exact original maps, proving all four
clauses. The hypotheses $\mathrm{Good}$ and
$\mathrm{IoW}$ are retained at the original endpoint;
this particular proof does not need their additional
arithmetic consequences once the displayed records
have been supplied.
\end{proof}
% END WIS18

\subsection{Closed images and fixed elements}

% BEGIN WIS19
\paragraph{WIS19: the original $p$-adic closed-image theorem.}
\label{claim:WIS19}
Let $j_0:E_0\to R_0$ be a unital map of commutative rings.
Assume $R_0$ is $p$-torsion-free and $p$-adically separated,
$E_0$ is $p$-adically complete and separated, and
\[
 j_0(e)=pr\Longrightarrow e=pe'
                              \text{ for some }e'\in E_0. \tag{WIS.29}
\]
If $x\in R_0$ admits approximations
$x-j_0(e_k)\in p^kR_0$ for every $k$, then
$x=j_0(e)$ for some $e\in E_0$.
In particular the actual image $I.ioE(E)$ is $p$-adically
closed in $W$, and the induced map
$E/p^kE\to W/p^kW$ is injective for every $k$.

\begin{proof}
First iterate (WIS.29). If $j_0(e)=p^{k+1}r$,
write $e=pe'$ using (WIS.29). Cancellation of $p$
in $R_0$ gives $j_0(e')=p^kr$, and induction gives
$e\in p^{k+1}E_0$. Thus $j_0$ reflects every
$p$-power divisibility, without an assumption that
$j_0$ is injective.

For $k\leq k'$ the approximations imply
$j_0(e_k-e_{k'})\in p^kR_0$. The reflected
divisibility makes $(e_k)$ a $p$-adic Cauchy system.
Choose its limit $e$. Both
$x-j_0(e_k)$ and $j_0(e_k-e)$ lie in $p^kR_0$,
so $x-j_0(e)$ lies in every $p^kR_0$ and is zero.
This proves the full original generic
\texttt{padic\_\allowbreak closed} implication.

Apply it to $I.ioE$ using (WIS.28), WIS5 and WIS18.
The kernel of the map modulo $p^k$ is zero by
(WIS.19), proving the last assertion as well.
\end{proof}
% END WIS19

% BEGIN WIS20
\paragraph{WIS20: the full fixed-subring $p$-adic supplier.}
\label{claim:WIS20}
For any set $H$ of elements of the original Galois group, let
\[
             W^H=\{w\in W:\sigma(w)=w\text{ for all }\sigma\in H\}.
\]
No subgroup or closedness hypothesis on $H$ is required.
The following original properties hold:
\begin{enumerate}
\item $pw\in W^H$ implies $w\in W^H$.
\item $w\in W^H$ implies $\phi(w)\in W^H$.
\item If $w_i\in W^H$ for every $i$, there exists $w\in W^H$
      such that for every $k$,
      \[
                  w-\sum_{i<k}p^iw_i\in p^kW.             \tag{WIS.30}
      \]
\end{enumerate}

\begin{proof}
The Galois action is the identity on $O$ and the
Witt-vector lift of its action on $F$.
It fixes $p$. Hence $pw$ fixed gives
$p(\sigma(w)-w)=0$, and WIS5 permits cancellation.

On $F$ every ring automorphism commutes with $x\mapsto x^p$.
Applying this equality in Witt coordinates and then
tensoring with the identity on $O$ shows
$\sigma\phi=\phi\sigma$. This proves the second clause.

The partial sums in (WIS.30) form a $p$-adic Cauchy
system in the complete ring $W$, so they have a
limit $w$ satisfying its stated congruences.
For every $\sigma\in H$, applying $\sigma$ to each
congruence is legitimate algebraically: a ring map
preserves $p^kW$. Since the partial sum is fixed,
$\sigma(w)-w\in p^kW$ for every $k$.
Separatedness of $W$ gives $\sigma(w)=w$.
This proves the full original
\texttt{fix\allowbreak Sub\_\allowbreak ladic} statement, including its
summation conclusion for arbitrary $H$.
\end{proof}
% END WIS20

% BEGIN WIS21
\paragraph{WIS21: an actual fixed lift of an element fixed modulo $p$.}
\label{claim:WIS21}
For any set $H$ as above, suppose $x\in W$ satisfies
$\sigma(x)-x\in pW$ for all $\sigma\in H$.
Then there is $y\in W^H$ with $y-x\in pW$.
Using the same integral basis as WIS4, it can be chosen as
\[
            y=\sum_i b_i\otimes[(x_i)_0],
            \qquad x=\sum_i b_i\otimes x_i.                \tag{WIS.31}
\]

\begin{proof}
By WIS4 and WIS2, the congruence for $\sigma(x)-x$
is equivalent to
$\sigma((x_i)_0)=(x_i)_0$ in $F$ for each $i$.
The Teichmuller lift commutes with every coefficient
automorphism: $\sigma([a])=[\sigma(a)]$.
Thus each term in (WIS.31) is fixed and $y\in W^H$.
Its coordinates have the same zeroth Witt coordinates
as those of $x$, so $y-x\in pW$ by WIS2 and WIS4.
The finite coefficient basis is fixed by the original
action; averaging over $H$ is neither needed nor used.
This proves the full original fixed-modulo-prime
lifting statement for every permitted set $H$.
\end{proof}
% END WIS21

\subsection{The exact cyclotomic approximation filtration}

For the rest of this subsection keep the original $P,AI,L,s$,
but no series record $R$ or embedding $I$ is required.
Put
\[
 \tau=\epsilon_s-1\in S,\qquad
 \pi=c(1\otimes\mu_s)\in W,\qquad
 \mathcal N_m=\pi^mc(B)+pW\quad(m\geq0).                    \tag{WIS.32}
\]
Each $\mathcal N_m$ is an additive subgroup. It is not being
treated as the ideal $\pi^mW+pW$; in the embedding situation
$\pi$ is a unit of $W$, so that ideal would lose the
filtration. The restriction of the coefficient to $c(B)$
in (WIS.32) is essential.

% BEGIN WIS22
\paragraph{WIS22: separation and complete summation in the actual tilt.}
\label{claim:WIS22}
For every $s\geq0$, the original tilt ring $S$ satisfies
\[
 \bigcap_{m\geq0}\tau^mS=0,\qquad
 \forall(a_i)_{i\geq0}\subset S,\quad
 \exists a\in S:\quad
 a-\sum_{i<m}\tau^ia_i\in\tau^mS\quad\text{for every }m.
                                                               \tag{WIS.33}
\]
These are the two clauses of the original
tilt separation-and-completeness theorem.

\begin{proof}
In the inverse-limit description of $S$, the $r$th
coordinate of $\epsilon_s$ is
$\zeta_{s+r}\bmod p$: inverse Frobenius shifts the
compatible coordinates. The cyclotomic identities give
$\zeta_{s+r}^{p^{s+r}}=1$. Thus, in $\mathcal O_C/p$,
\[
                   \tau_r^{M(r)}=0,\qquad M(r)=p^{s+r}.
                                                               \tag{WIS.34}
\]
If $a$ is divisible by every power of $\tau$, use
$\tau^{M(r)}\mid a$ to conclude that its $r$th
coordinate is zero. This for all $r$ proves separation.

For a sequence $(a_i)$, the $r$th coordinate of
$\sum_{i<N}\tau^ia_i$ is constant once $N\geq M(r)$.
Define $a_r$ to be that stable value. These coordinates
are compatible under Frobenius: compare adjacent
coordinates at any common cutoff at least
$\max(M(r),M(r+1))$, and use compatibility of the
finite sum. Hence they define $a\in S$.

For a fixed $m$, repeat this construction for the tail
sequence $a_m,a_{m+1},\ldots$, obtaining $b_m\in S$.
At each coordinate, evaluate the two sums at a common
large finite cutoff. Cancellation of the first $m$
terms then gives
\[
                     a-\sum_{i<m}\tau^ia_i=\tau^m b_m.
 \tag{WIS.35}
\]
This proves the complete summation assertion, including
literal divisibility of every tail. No convergence of
a chosen sequence in the characteristic-zero period
ring is assumed here; all limits were constructed
by stabilized finite coordinates in the tilt.
\end{proof}
% END WIS22

% BEGIN WIS23
\paragraph{WIS23: the original approximation condition in coordinates.}
\label{claim:WIS23}
For $x=\sum_i b_i\otimes x_i\in W$, write
$q_i(x)=(x_i)_0\in F$. Let $\iota:S\hookrightarrow F$
be the original inclusion. Then
\[
 \begin{split}
 x\in pW&\quad\Longleftrightarrow\quad q_i(x)=0
                                      \text{ for all }i,\\
 x\in\mathcal N_m&\quad\Longleftrightarrow\quad
 q_i(x)\in\iota(\tau^mS)\text{ for all }i.
 \end{split}                                               \tag{WIS.36}
\]
These are the actual \texttt{dvd\_\allowbreak iff\_\allowbreak crd} and
\texttt{near\allowbreak Pi\_\allowbreak iff\_\allowbreak crd} criteria, with their prescribed
maps and parameter.

\begin{proof}
The first assertion follows from WIS2 and WIS4.
The zeroth Witt coordinate of $\mu_s$ is $\tau$.
Hence the coordinates of $\pi^mc(b)$ at degree zero
are $\iota(\tau^m d_i)$, where $d_i$ are the zeroth
coordinates of $b$. The term $pw$ contributes zero,
proving the forward implication of the second assertion.

Conversely, write $q_i(x)=\iota(\tau^m d_i)$.
Choose $b=\sum_i b_i\otimes[d_i]\in B$.
The degree-zero coordinates of $x-\pi^mc(b)$
are all zero, so this difference equals $pw$
for some $w\in W$. This is precisely the original
definition of membership in $\mathcal N_m$.
All coordinate maps and inclusions used here are
the actual tensor maps, not independently selected
identifications of quotient rings.
\end{proof}
% END WIS23

% BEGIN WIS24
\paragraph{WIS24: the full original approximation separation and summation.}
\label{claim:WIS24}
For the subgroups (WIS.32),
\[
       \bigcap_{m\geq0}\mathcal N_m=pW.                    \tag{WIS.37}
\]
If $k\geq0$ and $f_i\in\mathcal N_{i+k}$ for every $i$,
there is $x\in W$ such that
\[
                      x-\sum_{i<m}f_i\in\mathcal N_{m+k}
                          \quad\text{for all }m\geq0.     \tag{WIS.38}
\]
This proves both clauses of the original
\texttt{near\allowbreak Pi\_\allowbreak sep\_\allowbreak complete} theorem for every $s$.

\begin{proof}
Certainly $pW$ lies in every $\mathcal N_m$.
If $x$ lies in their intersection, (WIS.36) at $m=0$
writes $q_i(x)=\iota(d_i)$ with $d_i\in S$.
At any other $m$, the injectivity of $\iota$ says
$d_i\in\tau^mS$. WIS22 makes each $d_i$ zero,
and (WIS.36) gives $x\in pW$.

For (WIS.38), choose $d_{i,j}\in S$ with
\[
                         q_j(f_i)=\iota(\tau^{i+k}d_{i,j}).
\]
For each of the finitely many coefficient indices $j$,
WIS22 constructs $a_j\in S$ such that
$a_j-\sum_{i<m}\tau^id_{i,j}\in\tau^mS$.
Take the actual element
\[
                      x=\sum_j b_j\otimes[\iota(\tau^k a_j)]
                               \in W.
 \tag{WIS.39}
\]
Its difference from the partial sum has every
degree-zero coordinate in $\iota(\tau^{m+k}S)$.
WIS23 gives (WIS.38). This proves the exact
approximation statement; it does not replace it
with an assertion about a different topology on $W$.
\end{proof}
% END WIS24

% BEGIN WIS25
\paragraph{WIS25: the full original fixed-subring approximation supplier.}
\label{claim:WIS25}
Let $H$ be any set of elements of the original Galois group
and suppose $\pi\in W^H$. If $f_i\in W^H\cap
\mathcal N_{i+k}$ for every $i$, there exists $y\in W^H$
such that
\[
                    y-\sum_{i<m}f_i\in\mathcal N_{m+k}
                               \quad\text{for every }m.  \tag{WIS.40}
\]
Thus the complete original
\texttt{fix\allowbreak Sub\_\allowbreak near\allowbreak Pi\_\allowbreak complete} conclusion holds.

\begin{proof}
Apply WIS24 first without a fixedness requirement
and obtain $x$ satisfying (WIS.38). Each original
Galois action preserves $c(B)$, because it acts
on the tilt and its Witt vectors and fixes $O$.
For $\sigma\in H$, it also fixes $\pi$. Thus it
preserves each $\mathcal N_m$.

Let $S_m=\sum_{i<m}f_i$. This partial sum is fixed, so
\[
              \sigma(x)-x
                =\sigma(x-S_m)-(x-S_m)\in\mathcal N_{m+k}.
 \tag{WIS.41}
\]
Since $\pi\in c(B)$, the subgroups $\mathcal N_m$
decrease with $m$. It follows that the left side
of (WIS.41) lies in every $\mathcal N_m$.
By WIS24 it lies in $pW$. WIS21 therefore supplies
an actual fixed $y$ with $y-x\in pW$.
As $pW\subseteq\mathcal N_{m+k}$ for every $m$,
adding this correction to (WIS.38) proves (WIS.40).
The same basis and the same parameter $\pi$ have
been used throughout. In particular no additional
restriction on $H$ or on the allowed index $s$
has been imposed.
\end{proof}
% END WIS25

% BEGIN WIS26
\paragraph{WIS26: the precise return to the Wach-lattice construction.}
\label{claim:WIS26}
On every original permitted layer with its supplied
$\Gamma,R,I,E_{\mathrm{et}}$, retain the original
$\operatorname{LocTied}(p,P,AI)$ hypothesis (hence the
$\operatorname{IsStd}$ instance used for realization), together
with the original Bhatt--Scholze crystalline-realization input. Suppose
the following two original suppliers are available:
\begin{enumerate}
\item The actual Wach prism with its maps, coefficient
      coordinates and scalar-map identities.
\item Descent of the original cyclotomic-kernel-fixed
      homogeneous Frobenius matrix solutions, at every
      required rank on that layer.
\end{enumerate}
Then the full original Wach-lattice existence conclusion
of WLD26 holds. A separate assumption of injectivity or
unit reflection for $I.ioE$ is no longer needed.

\begin{proof}
Apply WIS18 to the supplied $I$ to obtain the two
clauses used by WLD26, namely injectivity and unit
reflection. The two retained suppliers are exactly
the other two period-ring inputs in that theorem.
All its original coefficient data and realizing
crystals remain unchanged. Thus WLD26 applies,
producing the same inverse-image lattice
$N_T=j_T^{-1}(L_T)$ inside the supplied
\'etale module, with all the fields already
proved in WLD18--WLD25.

This implication is uniform over the original
allowed layer indices and over the supplied
\'etale records. No transition maps between
different outer indices are added. The independent
faithful-invariant, \'etale-basis, Herr-class-map
and quantitative comparison inputs needed after
Wach existence are not consequences of this
two-supplier return.
\end{proof}
% END WIS26

\subsection{The remaining period-ring work}

The embedding theorem WIS18 has its full four-clause original
scope. Its first two clauses feed the lattice construction;
its saturation clause also feeds the original invariant
Frobenius-matrix proof. WIS19 supplies that proof's
$p$-adic closed-image input, WIS20 its fixed-subring
$p$-adic input, and WIS23--WIS25 its original approximation
separation and completeness inputs.

The matrix-solution theorem itself is still to be reconstructed.
Its remaining argument involves the Frobenius behavior of the
approximation filtration, approximation of actual cyclotomic
invariants by the specified coefficients, the closedness
modulo $p$ in that filtration, and the matrix lifting step.
The exact source endpoint is
\texttt{Fixed\allowbreak By\allowbreak Cyclotomic\allowbreak Kernel\allowbreak Etale\allowbreak Phi\allowbreak Solution\allowbreak Mem\allowbreak Wach}.
These requirements are not supplied merely by the existence
of the limits in WIS24 or WIS25.

The actual prism in
\texttt{Exists\allowbreak Wach\allowbreak Prism\allowbreak Over\allowbreak Integers} is the other remaining
period-ring supplier. Both have concrete maps and identities
that must be verified before WIS26 becomes an unconditional
Wach-existence result. The later full local comparison package
and the common family of quantitative local bounds are
separate subsequent tasks. No whole selected global obligation
is declared closed in this chapter.



\endgroup
\end{document}
